% Lesson 1 — Listening: Texts about undergoing a job interview — Anglais Tle A — Cours signé OnBuch+
% Programme officiel MINESEC. Compiler avec : tectonic EN01-listening-texts-about-undergoing-a-job-interview.tex
\newcommand{\DOCMATIERE}{Anglais}
\newcommand{\DOCNIVEAU}{Tle A}
\newcommand{\DOCMODULE}{Chapter 1 : Family and Social Life}
\newcommand{\DOCLECON}{EN01}
\newcommand{\DOCTITRE}{Lesson 1 — Listening: Texts about undergoing a job interview}
% OnBuch+ — préambule « cours signés » ENRICHI (compilé avec tectonic / XeLaTeX)
% Système visuel « Pop » d'OnBuch+ : fond crème (jamais blanc), bordures encre
% épaisses, ombres « dures » décalées, titres Archivo Black en capitales.
% Version MATHÉMATIQUES : graphes (pgfplots/tikz), boîtes pédagogiques
% (définition, propriété, méthode, attention, le savais-tu, exercice résolu,
% exercice, TP, bilan), page de garde et signature OnBuch+ ancrée en bas.
%
% Macros attendues dans main.tex AVANT \input{preamble} :
%   \DOCMATIERE  \DOCNIVEAU  \DOCMODULE  \DOCLECON (numéro)  \DOCTITRE
\documentclass[11pt]{article}

\usepackage{fontspec}
\usepackage[english,french]{babel} % français = langue principale ; l'anglais via \eng{..} / textebox
\usepackage[a4paper,margin=2cm,top=3.4cm,bottom=2.8cm,headheight=1.4cm,footskip=1.3cm]{geometry}
\usepackage{amsmath,amssymb,amsfonts}
\usepackage{mathtools} % \xmapsto, \coloneqq... souvent utilisés par les modèles
\usepackage{cancel} % simplifications barrées (bilans de forces)
% \abs{z} (module, valeur absolue) et \norm{u} (norme), utilisés par les modèles
\providecommand{\abs}{}\let\abs\relax\DeclarePairedDelimiter{\abs}{\lvert}{\rvert}
\providecommand{\norm}{}\let\norm\relax\DeclarePairedDelimiter{\norm}{\lVert}{\rVert}
\usepackage[table]{xcolor} % [table] : fournit aussi \rowcolors (alternance de lignes)
\usepackage{pagecolor}
\usepackage{tikz}
\usetikzlibrary{arrows.meta,positioning,calc,shapes.geometric,shapes.misc,shapes.arrows,decorations.pathmorphing,decorations.pathreplacing,decorations.markings,patterns,shadows,mindmap,trees,fit,backgrounds,matrix,intersections,angles,quotes}
\usepackage{pgfplots}
\usepackage[version=4]{mhchem} % équations nucléaires \ce{^{235}_{92}U}
\usepackage[european,siunitx]{circuitikz} % schémas électriques
\pgfplotsset{compat=1.18}
\usepgfplotslibrary{fillbetween,groupplots} % aires entre courbes (calcul intégral)
\usepackage{enumitem}
\usepackage{fancyhdr}
% [most] charge toutes les bibliothèques tcolorbox (listings, minted, raster,
% documentation...) et gonfle inutilement la mémoire TeX sur un document long
% (~300 boîtes) : on ne charge que ce qui est réellement utilisé.
\usepackage[skins,breakable]{tcolorbox}
\usepackage{graphicx}
\usepackage{colortbl}
\usepackage{booktabs}
\usepackage{tabularx}
\usepackage{diagbox} % case d'en-tête diagonale des tableaux à double entrée
% Colonnes extensibles centrée / gauche / droite (C, L, R), souvent utilisées par les modèles
\newcolumntype{C}{>{\centering\arraybackslash}X}
\newcolumntype{L}{>{\raggedright\arraybackslash}X}
\newcolumntype{R}{>{\raggedleft\arraybackslash}X}
\usepackage{array}
\usepackage{multicol}
\usepackage{multirow}
\usepackage{siunitx}
% parse-numbers=false : accepte aussi \SI{2,5 \times 10^{-3}}{...}
\sisetup{locale=FR,output-decimal-marker={,},group-minimum-digits=4,parse-numbers=false}
\DeclareSIUnit\ohm{\ensuremath{\symup{\Omega}}} % U+2126 absent de la police
\DeclareSIUnit\division{div} % réglages d'oscilloscope (ms/div, V/div)
\DeclareSIUnit\tr{tr} % tours (vitesse de rotation, ex. \SI{1500}{\tr\per\minute})
\DeclareSIUnit\molar{M} % concentration molaire (ex. \SI{3}{\molar}, \SI{1}{\micro\molar})
\DeclareSIUnit\an{an} % année (le modèle confond parfois avec \year, un compteur TeX) — ex. \SI{5}{\kilo\meter\per\an}
\DeclareSIUnit\FCFA{FCFA} % franc CFA (ex. \SI{500}{\FCFA\per\kilo\gram})
\DeclareSIUnit\degreeBrix{\textdegree Brix} % teneur en sucre d'un sirop/jus (réfractométrie)
\DeclareSIUnit\month{mois} % le modèle utilise parfois \month, un compteur TeX, comme unité de durée
\usepackage{qrcode}
% \degree : commande fréquemment utilisée par les modèles pour un angle ou une
% température en dehors de \SI{}{} (non fournie par les paquets chargés).
\providecommand{\degree}{^{\circ}}
% \qed (fin de démonstration), utilisé par les modèles sans amsthm
\providecommand{\qed}{\hfill$\square$}
% \barre : double barre des valeurs interdites dans un tableau de variations (tkz-tab)
\providecommand{\barre}{\ensuremath{\|}}
% environnement proof (amsthm non chargé) utilisé par les modèles
\ifdefined\proof\else\newenvironment{proof}[1][Démonstration]{\par\noindent\textit{#1.}\ }{\qed\par}\fi
\usepackage{lastpage}
\usepackage{hyperref}
\hypersetup{hidelinks}

% ============================================================
% Palette « Pop » OnBuch+ — mêmes hex que la classe P (plus_ui.dart)
% ============================================================
\definecolor{poporange}{HTML}{F59321}
\definecolor{popdark}{HTML}{E07A0C}
\definecolor{popcream}{HTML}{FFF6E8}
\definecolor{popink}{HTML}{1C1714}
\definecolor{poppurple}{HTML}{7A5AE0}
\definecolor{popgreen}{HTML}{2E9E6A}
\definecolor{poppink}{HTML}{E0457B}
\definecolor{popblue}{HTML}{2F7DE1}
\definecolor{popgold}{HTML}{C98A12}
\definecolor{popmuted}{HTML}{8A7F76}
\definecolor{popline}{HTML}{EADFD2}
\definecolor{popgray}{HTML}{8A7F76}
\colorlet{popgrey}{popgray}
% Alias tolérants (le modèle nomme parfois les couleurs différemment)
\colorlet{popwhite}{white}
\colorlet{popblack}{popink}
\colorlet{popred}{poppink}
\colorlet{popyellow}{popgold}
% Teintes claires (fonds de tableaux/encadrés) — le modèle nomme parfois
% les couleurs avec un suffixe L (ex. popblueL) sans les avoir définies.
\colorlet{poporangeL}{poporange!20}
\colorlet{popdarkL}{popdark!20}
\colorlet{poppurpleL}{poppurple!20}
\colorlet{popgreenL}{popgreen!20}
\colorlet{poppinkL}{poppink!20}
\colorlet{popblueL}{popblue!20}
\colorlet{popgoldL}{popgold!20}
\colorlet{popredL}{popred!20}
\colorlet{popgrayL}{popgray!20}
% Teintes claires pour fonds de tableaux / zones
\colorlet{popblueL}{popblue!10!white}
\colorlet{poporangeL}{poporange!14!white}
\colorlet{popgreenL}{popgreen!12!white}
\colorlet{poppurpleL}{poppurple!12!white}
\colorlet{poppinkL}{poppink!12!white}
\colorlet{popgoldL}{popgold!14!white}

\pagecolor{popcream}

% ============================================================
% Typographie
% ============================================================
\defaultfontfeatures{Ligatures=TeX}
\setmainfont{PlusJakartaSans-Medium.ttf}[Path=../../../fonts/,BoldFont=PlusJakartaSans-Bold.ttf]
% Maths en sans-serif (Fira Math) avec chiffres et lettres droites de Plus
% Jakarta, pour que les formules se fondent dans le texte.
\usepackage[math-style=ISO,bold-style=ISO]{unicode-math}
\setmathfont{FiraMath-Regular.otf}[Path=../../../fonts/]
\setmathfont{PlusJakartaSans-Medium.ttf}[Path=../../../fonts/,range={up/num,up/{Latin,latin},\mathup}]
\usepackage{icomma} % virgule décimale sans espace en mode math
% Caractères mathématiques Unicode absents des polices de texte (Plus Jakarta,
% Archivo Black) : rendus par leur équivalent mathématique, en texte comme en
% formule — sinon ils s'affichent en carré vide (ex. « Arithmétique dans ℤ »).
\usepackage{newunicodechar}
\newunicodechar{ℝ}{\ensuremath{\mathbb{R}}}
\newunicodechar{ℤ}{\ensuremath{\mathbb{Z}}}
\newunicodechar{ℂ}{\ensuremath{\mathbb{C}}}
\newunicodechar{ℕ}{\ensuremath{\mathbb{N}}}
\newunicodechar{ℚ}{\ensuremath{\mathbb{Q}}}
\newunicodechar{𝔻}{\ensuremath{\mathbb{D}}}
\newunicodechar{α}{\ensuremath{\alpha}}
\newunicodechar{β}{\ensuremath{\beta}}
\newunicodechar{γ}{\ensuremath{\gamma}}
\newunicodechar{δ}{\ensuremath{\delta}}
\newunicodechar{ε}{\ensuremath{\varepsilon}}
\newunicodechar{θ}{\ensuremath{\theta}}
\newunicodechar{λ}{\ensuremath{\lambda}}
\newunicodechar{μ}{\ensuremath{\mu}}
\newunicodechar{σ}{\ensuremath{\sigma}}
\newunicodechar{τ}{\ensuremath{\tau}}
\newunicodechar{φ}{\ensuremath{\varphi}}
\newunicodechar{ψ}{\ensuremath{\psi}}
\newunicodechar{ω}{\ensuremath{\omega}}
\newunicodechar{Σ}{\ensuremath{\Sigma}}
% majuscules produites par \MakeUppercase dans les titres (α-aminés -> Α)
\newunicodechar{Α}{\ensuremath{\alpha}}
\newunicodechar{Β}{\ensuremath{\beta}}
\newunicodechar{Γ}{\ensuremath{\Gamma}}
\newunicodechar{Δ}{\ensuremath{\Delta}}
\newunicodechar{Ε}{\ensuremath{\varepsilon}}
\newunicodechar{Θ}{\ensuremath{\Theta}}
\newunicodechar{Λ}{\ensuremath{\Lambda}}
\newunicodechar{Μ}{\ensuremath{\mu}}
\newunicodechar{Τ}{\ensuremath{\tau}}
\newunicodechar{Φ}{\ensuremath{\Phi}}
\newunicodechar{Ψ}{\ensuremath{\Psi}}
\newunicodechar{Ω}{\ensuremath{\Omega}}
\newunicodechar{↦}{\ensuremath{\mapsto}}
\newunicodechar{∈}{\ensuremath{\in}}
\newunicodechar{∉}{\ensuremath{\notin}}
\newunicodechar{∧}{\ensuremath{\wedge}}
\newunicodechar{⇔}{\ensuremath{\Leftrightarrow}}
\newunicodechar{⟺}{\ensuremath{\Longleftrightarrow}}
\newunicodechar{⇒}{\ensuremath{\Rightarrow}}
\newunicodechar{⟹}{\ensuremath{\Longrightarrow}}
\newunicodechar{∘}{\ensuremath{\circ}}
\newunicodechar{∪}{\ensuremath{\cup}}
\newunicodechar{∩}{\ensuremath{\cap}}
\newunicodechar{≡}{\ensuremath{\equiv}}
\newunicodechar{⊂}{\ensuremath{\subset}}
\newunicodechar{⊥}{\ensuremath{\perp}}
\newunicodechar{∃}{\ensuremath{\exists}}
\newunicodechar{∀}{\ensuremath{\forall}}
\newunicodechar{∛}{\ensuremath{\sqrt[3]{\,}}}
\newunicodechar{∜}{\ensuremath{\sqrt[4]{\,}}}
\newunicodechar{⃗}{\ensuremath{{}^{\to}}}
\newunicodechar{ⁿ}{\ifmmode^{n}\else\textsuperscript{\ensuremath{n}}\fi}
\newunicodechar{ˣ}{\ifmmode^{x}\else\textsuperscript{\ensuremath{x}}\fi}
\newunicodechar{ᵇ}{\ifmmode^{b}\else\textsuperscript{\ensuremath{b}}\fi}
\newunicodechar{ʳ}{\ifmmode^{r}\else\textsuperscript{\ensuremath{r}}\fi}
\newunicodechar{ᵘ}{\ifmmode^{u}\else\textsuperscript{\ensuremath{u}}\fi}
\newunicodechar{ʸ}{\ifmmode^{y}\else\textsuperscript{\ensuremath{y}}\fi}
\newunicodechar{ᵃ}{\ifmmode^{a}\else\textsuperscript{\ensuremath{a}}\fi}
\newunicodechar{ᵉ}{\ifmmode^{e}\else\textsuperscript{\ensuremath{e}}\fi}
\newunicodechar{ᵏ}{\ifmmode^{k}\else\textsuperscript{\ensuremath{k}}\fi}
\newunicodechar{ᵖ}{\ifmmode^{p}\else\textsuperscript{\ensuremath{p}}\fi}
\newunicodechar{ᴺ}{\ifmmode^{N}\else\textsuperscript{\ensuremath{N}}\fi}
\newunicodechar{ᶜ}{\ifmmode^{c}\else\textsuperscript{\ensuremath{c}}\fi}
\newunicodechar{ʰ}{\ifmmode^{h}\else\textsuperscript{\ensuremath{h}}\fi}
\newunicodechar{ᵠ}{\ifmmode^{q}\else\textsuperscript{\ensuremath{q}}\fi}
\newunicodechar{ₙ}{\ifmmode_{n}\else\textsubscript{\ensuremath{n}}\fi}
\newunicodechar{ₐ}{\ifmmode_{a}\else\textsubscript{\ensuremath{a}}\fi}
\newunicodechar{ᵢ}{\ifmmode_{i}\else\textsubscript{\ensuremath{i}}\fi}
\newunicodechar{ᵣ}{\ifmmode_{r}\else\textsubscript{\ensuremath{r}}\fi}
\newunicodechar{ₖ}{\ifmmode_{k}\else\textsubscript{\ensuremath{k}}\fi}
\newunicodechar{ᵦ}{\ifmmode_{\beta}\else\textsubscript{\ensuremath{\beta}}\fi}
\newunicodechar{ₓ}{\ifmmode_{x}\else\textsubscript{\ensuremath{x}}\fi}
\newfontfamily\popdisplay{ArchivoBlack-Regular.ttf}[Path=../../../fonts/]
\newfontfamily\poplogo{SpaceGrotesk-Bold.ttf}[Path=../../../fonts/]
\newfontfamily\popsemi{PlusJakartaSans-SemiBold.ttf}[Path=../../../fonts/]
\newfontfamily\popxbold{PlusJakartaSans-ExtraBold.ttf}[Path=../../../fonts/]
\newfontfamily\popxboldls{PlusJakartaSans-ExtraBold.ttf}[Path=../../../fonts/,LetterSpace=3.0]

\setlist[enumerate]{leftmargin=1.7em,itemsep=5pt,topsep=4pt}
\setlist[itemize]{leftmargin=1.7em,itemsep=4pt,topsep=4pt,label={\color{poporange}\textbullet}}
\setlength{\parindent}{0pt}
\setlength{\parskip}{5pt}
\renewcommand{\arraystretch}{1.25}

% pgfplots : style Pop par défaut
\pgfplotsset{
  every axis/.append style={axis line style={popink,line width=1pt},tick style={popink},
    grid style={popline,line width=0.5pt},label style={font=\small},tick label style={font=\footnotesize}},
  popcurve/.style={poporange,line width=1.8pt,no markers,smooth},
}
% Même style disponible dans un \draw TikZ simple (hors pgfplots)
\tikzset{popcurve/.style={poporange,line width=1.8pt}}
\tikzset{popgrid/.style={popline,very thin}}
\colorlet{popcurve}{poporange} % le nom du style est parfois utilisé comme couleur

% ============================================================
% Pill / squiggle
% ============================================================
\newcommand{\poppill}[3]{%
  \tikz[baseline=-0.55ex]\node[draw=popink,line width=1.2pt,rounded corners=2.6pt,%
    fill=#2,inner xsep=6.5pt,inner ysep=3pt]{%
    \popxboldls\fontsize{7}{7}\selectfont\color{#3}\MakeUppercase{#1}};%
}
\newcommand{\popsquiggle}[1]{%
  \tikz[baseline=-0.4ex]{
    \draw[#1,line width=2.6pt,line cap=round]
      plot [smooth] coordinates {(0,0.09) (0.34,0.01) (0.68,0.21) (1.02,0.01) (1.36,0.10)};
  }%
}
% Mot-clé surligné (marqueur orange sous le texte)
\newcommand{\cle}[1]{{\popxbold\color{popdark}#1}}

% ============================================================
% En-tête / pied
% ============================================================
\pagestyle{fancy}
\fancyhf{}
\renewcommand{\headrulewidth}{0pt}
\renewcommand{\footrulewidth}{0pt}
\lhead{\raisebox{-0.15ex}{\poplogo\fontsize{13}{13}\selectfont\color{popink}OnBuch\color{poporange}+}}
\rhead{\poppill{\DOCMATIERE\ \textbullet\ \DOCNIVEAU\ \textbullet\ Leçon \DOCLECON}{white}{popink}}
\lfoot{\popxbold\fontsize{7.6}{7.6}\selectfont\color{popmuted}\MakeUppercase{Cours signé OnBuch+}\ \textbullet\ {\popsemi\DOCTITRE}}
\rfoot{\tikz[baseline=-0.6ex]\node[rounded corners=8.5pt,fill=popink,minimum height=17pt,minimum width=17pt,inner xsep=5pt,inner sep=0]{\popxbold\fontsize{8}{8}\selectfont\color{popcream}\thepage\,/\,\pageref*{LastPage}};}

% ============================================================
% Titres
% ============================================================
\newcommand{\chaptitre}[2]{%
  \begin{tcolorbox}[enhanced,colback=poporange,colframe=popink,boxrule=2.6pt,arc=11pt,
      drop shadow=popink,left=15pt,right=15pt,top=12pt,bottom=11pt,before skip=3pt,after skip=18pt]
    \poppill{#2}{popink}{popcream}\par\vspace{9pt}
    {\raggedright\hyphenpenalty=10000\exhyphenpenalty=10000\popdisplay\fontsize{22}{24}\selectfont\color{popcream}\MakeUppercase{#1}\par}
  \end{tcolorbox}
}
% Section numérotée : pastille orange avec le numéro + titre Archivo
\newcounter{popsec}
\newcommand{\coursec}[1]{%
  \stepcounter{popsec}\setcounter{popsub}{0}%
  \par\vspace{16pt}\noindent
  \begin{minipage}{\linewidth}
  \tikz[baseline=-0.7ex]\node[draw=popink,line width=1.6pt,fill=poporange,rounded corners=4pt,minimum size=22pt,inner sep=0]{\popdisplay\fontsize{12}{12}\selectfont\color{popcream}\thepopsec};%
  \hspace{8pt}{\popdisplay\fontsize{15}{17}\selectfont\color{popink}\MakeUppercase{#1}}\par
  \vspace{4pt}\noindent{\color{poporange}\rule{36pt}{3.6pt}}
  \end{minipage}\par\nopagebreak\vspace{8pt}
}
\newcounter{popsub}
\newcommand{\courssub}[1]{%
  \stepcounter{popsub}%
  \par\vspace{9pt}\noindent{\popxbold\fontsize{11.5}{13}\selectfont\color{popink}{\color{poporange}\thepopsec.\thepopsub}\hspace{6pt}#1}\par\nopagebreak\vspace{4pt}
}

% ============================================================
% Boîtes pédagogiques — même recette : fond blanc, bordure épaisse colorée,
% ombre dure encre, badge plein. Titre optionnel [#1].
% ============================================================
\tcbset{popbase/.style={breakable,enhanced,colback=white,boxrule=2.3pt,arc=9pt,
  shadow={3pt}{-3pt}{0pt}{popink},left=14pt,right=14pt,top=11pt,bottom=11pt,before skip=11pt,after skip=15pt,
  fontupper=\popsemi\fontsize{10.6}{14.5}\selectfont\color{popink},
  fontlower=\popsemi\fontsize{10.6}{14.5}\selectfont\color{popink}}}
% Coche dessinée (le glyphe ✓ n'existe pas dans les polices du document)
\newcommand{\popcheck}[1]{\tikz[baseline=-0.5ex]\draw[#1,line width=1.6pt,line cap=round,line join=round](-0.13,0.0)--(-0.04,-0.09)--(0.13,0.1);}
\providecommand{\coche}{\popcheck{popgreen}}
\providecommand{\resultat}[1]{\par\smallskip\noindent\fcolorbox{popgreen}{popgreenL}{\parbox{\dimexpr\linewidth-2\fboxsep-2\fboxrule\relax}{\textbf{Résultat.} #1}}\par\smallskip}
\newcommand{\popboxhead}[3]{\poppill{#1}{#2}{white}\ifx\relax#3\relax\else\hspace{6pt}{\popxbold\fontsize{10.6}{12}\selectfont\color{popink}#3}\fi\par\vspace{7pt}}

\newtcolorbox{aretenir}[1][]{popbase,colframe=popgreen,before upper={\popboxhead{À retenir}{popgreen}{#1}}}
% Texte en anglais (document de lecture, dialogue, extrait) : cadre
% d'étude. Un titre optionnel (source, auteur) s'affiche dans l'en-tête.
\newtcolorbox{textebox}[1][]{popbase,colframe=popblue,before upper={\popboxhead{Text}{popblue}{#1}\begin{otherlanguage}{english}},after upper={\end{otherlanguage}}}
% Marques ✓ / ✗ (forme correcte / fautive) souvent utilisées par les modèles
\providecommand{\cross}{\textcolor{poppink}{\ensuremath{\times}}}
\providecommand{\xmark}{\cross}\providecommand{\crossmark}{\cross}\providecommand{\cmark}{\ensuremath{\checkmark}}
% Ligne de réponse / justification à compléter (inventée par les modèles)
\providecommand{\justification}[1]{\par\noindent\textit{Justification~:} \dotfill\par}
% \textipa (package tipa non chargé) : la police ne couvre pas l'API -> simple passage
\providecommand{\textipa}[1]{#1}
% Soulignement (ulem non chargé) : \uline, \ul
\providecommand{\uline}[1]{\underline{#1}}\providecommand{\ul}[1]{\underline{#1}}
% \glossaire{\item ...} inventée par les modèles : liste de glossaire
\providecommand{\glossaire}[1]{\begin{itemize}#1\end{itemize}}
% \alert (beamer) inventée par les modèles : texte mis en évidence
\providecommand{\alert}[1]{\textbf{\textcolor{poppink}{#1}}}
% variantes ASCII d'ordinaux écrites en macro
\providecommand{\iere}{\textsuperscript{re}}\providecommand{\ieme}{\textsuperscript{e}}\providecommand{\ere}{\textsuperscript{re}}\providecommand{\eme}{\textsuperscript{e}}\providecommand{\er}{\textsuperscript{er}}\providecommand{\ie}{\textsuperscript{e}}
% Ordinaux français écrits en macro par les modèles : 1\ière, 3\ième
\newcommand{\ière}{\textsuperscript{re}}\newcommand{\ième}{\textsuperscript{e}}
% \exemple (inventée par les modèles) : étiquette « Exemple : »
\providecommand{\exemple}{\textbf{Exemple~:}\ }
% \square utilisable aussi hors mode math (cases à cocher)
\AtBeginDocument{\let\squareMath\square\renewcommand{\square}{\ensuremath{\squareMath}}}
% Mot ou phrase en anglais à l'intérieur d'un paragraphe français
\newcommand{\eng}[1]{\ifmmode\text{\textit{#1}}\else\begingroup\selectlanguage{english}\itshape #1\endgroup\fi}
\let\esp\eng % alias toléré
\newtcolorbox{exemplebox}[1][]{popbase,colframe=poporange,before upper={\popboxhead{Exemple}{poporange}{#1}}}
\newtcolorbox{definition}[1][]{popbase,colframe=popblue,before upper={\popboxhead{Définition}{popblue}{#1}}}
\newtcolorbox{propriete}[1][]{popbase,colframe=poppurple,before upper={\popboxhead{Propriété}{poppurple}{#1}}}
\newtcolorbox{methode}[1][]{popbase,colframe=popgold,before upper={\popboxhead{Méthode}{popgold}{#1}}}
\newtcolorbox{attention}[1][]{popbase,colframe=poppink,before upper={\popboxhead{Attention}{poppink}{#1}}}
\newtcolorbox{experience}[1][]{popbase,colframe=popblue,colback=popblueL,before upper={\popboxhead{Expérience}{popblue}{#1}}}
% « Le savais-tu ? » : carte violette pleine (contexte, culture, Cameroun)
\newtcolorbox{savaistu}[1][]{popbase,colback=poppurple,colframe=popink,
  fontupper=\popsemi\fontsize{10.4}{14.2}\selectfont\color{white},
  before upper={\poppill{Le savais-tu\,?}{white}{poppurple}\ifx\relax#1\relax\else\hspace{6pt}{\popxbold\color{white}#1}\fi\par\vspace{7pt}}}
% Exercice résolu : énoncé en haut, \tcblower puis la solution sur fond crème
\newcounter{popexo}\newcounter{popexr}
\newtcolorbox{exoresolu}[1][]{popbase,colframe=poporange,colbacklower=popcream,
  segmentation style={popink,line width=1.2pt,dashed},
  before upper={\stepcounter{popexr}\popboxhead{Exercice résolu \thepopexr}{poporange}{#1}},
  before lower={\poppill{Solution}{popink}{popcream}\par\vspace{6pt}}}
% Exercice (non résolu en place) — la correction est dans la partie Corrigés
\newtcolorbox{exercice}[1][]{popbase,colframe=popink,
  before upper={\stepcounter{popexo}\popboxhead{Exercice \thepopexo}{popink}{#1}}}
% Corrigé
\newtcolorbox{corrige}[1][]{popbase,colframe=popgreen,colback=popgreenL,
  before upper={\popboxhead{Corrigé}{popgreen}{#1}}}
% Objectifs / prérequis / bilan
\newtcolorbox{objectifs}{popbase,colframe=popink,colback=poporangeL,
  before upper={\popboxhead{Objectifs de la leçon}{popink}{}}}
\newtcolorbox{prerequis}{popbase,colframe=popink,colback=popblueL,
  before upper={\popboxhead{Prérequis}{popblue}{}}}
\newtcolorbox{bilan}{popbase,colframe=popink,colback=white,boxrule=2.8pt,
  before upper={\popboxhead{Fiche bilan}{poporange}{L'essentiel à connaître pour l'examen}}}
% Figure encadrée avec légende
\newtcolorbox{popfigure}[1][]{enhanced,breakable=false,colback=white,colframe=popline,boxrule=1.4pt,arc=8pt,
  left=10pt,right=10pt,top=10pt,bottom=8pt,before skip=10pt,after skip=12pt,halign=center,
  lower separated=false,halign lower=center,
  fontlower=\popsemi\fontsize{9}{11.5}\selectfont\color{popmuted},
  #1}
\newcommand{\legende}[1]{\par\vspace{6pt}{\popsemi\fontsize{9}{11.5}\selectfont\color{popmuted}#1}}

% ============================================================
% Signature OnBuch+
% ============================================================
\newcommand{\poplogoinline}[1]{{\poplogo\fontsize{#1}{#1}\selectfont\color{popink}OnBuch\color{poporange}+}}
\newcommand{\signatureonbuch}{%
  \begin{tcolorbox}[enhanced,colback=popink,colframe=popink,boxrule=2.2pt,arc=10pt,
      drop shadow=poporange,left=14pt,right=14pt,top=10pt,bottom=10pt,before skip=0pt,after skip=0pt]
    \tikz[baseline=-0.7ex]\node[circle,fill=popgreen,draw=popcream,line width=1.3pt,minimum size=22pt,inner sep=0]{\popcheck{white}};%
    \hspace{9pt}\begin{minipage}[c]{0.62\linewidth}
      {\popxbold\fontsize{12}{13}\selectfont\color{popcream}Signé\ \textbullet\ {\poplogo OnBuch\color{poporange}+}}\par\vspace{2pt}
      {\raggedright\popsemi\fontsize{8}{10.5}\selectfont\color{popline}\MakeUppercase{Équipe pédagogique OnBuch+}\\\MakeUppercase{Conforme au programme officiel MINESEC}\par}
    \end{minipage}\hfill
    {\poplogo\fontsize{22}{22}\selectfont\color{popcream}OnBuch\color{poporange}+}
  \end{tcolorbox}%
}

% ============================================================
% Page de garde
% #1 titre · #2 matière · #3 niveau · #4 module · #5 n° leçon · #6 durée ·
% #7 description / objectifs (texte ou itemize)
% ============================================================
\newcommand{\pagegarde}[7]{%
  \thispagestyle{empty}
  \newgeometry{a4paper,margin=1.7cm,top=1.6cm,bottom=1.4cm}%
  \begin{tikzpicture}[remember picture,overlay]
    \fill[poporange!18] ([xshift=-3.2cm,yshift=-3.6cm]current page.north east) circle (4.6cm);
    \fill[poppurple!14] ([xshift=2.4cm,yshift=7.6cm]current page.south west) circle (3.8cm);
  \end{tikzpicture}%
  {\poplogo\fontsize{30}{30}\selectfont\color{popink}OnBuch\color{poporange}+}\hfill
  \raisebox{0.6ex}{\poppill{Cours signé \textbullet\ Premium}{popink}{popcream}}\par
  \vspace{1pt}\popsquiggle{poppurple}\par
  \vspace{8pt}
  \begin{tcolorbox}[enhanced,colback=poporange,colframe=popink,boxrule=2.8pt,arc=13pt,
      drop shadow=popink,left=18pt,right=18pt,top=12pt,bottom=13pt,after skip=10pt]
    \poppill{#2 \textbullet\ #3}{popink}{popcream}\hspace{5pt}\poppill{#4}{white}{popink}\par\vspace{8pt}
    {\popdisplay\fontsize{14}{14}\selectfont\color{popink}LEÇON #5}\par\vspace{4pt}
    {\raggedright\hyphenpenalty=10000\exhyphenpenalty=10000\popdisplay\fontsize{25}{27}\selectfont\color{popcream}\MakeUppercase{#1}\par}
    \vspace{8pt}
    {\popxbold\fontsize{9.5}{9.5}\selectfont\color{popink}\MakeUppercase{Durée conseillée : #6}}
  \end{tcolorbox}
  \begin{tcolorbox}[enhanced,colback=white,colframe=popink,boxrule=2.2pt,arc=10pt,
      drop shadow=popink,left=16pt,right=16pt,top=10pt,bottom=10pt,after skip=8pt]
    \poppill{Dans cette leçon}{poppurple}{white}\par\vspace{5pt}
    \popsemi\fontsize{9.3}{12.6}\selectfont\color{popink}#7
  \end{tcolorbox}
  \noindent\begin{minipage}[t]{0.487\linewidth}
    \begin{tcolorbox}[enhanced,colback=popgreen,colframe=popink,boxrule=2.2pt,arc=10pt,
        drop shadow=popink,left=11pt,right=11pt,top=10pt,bottom=10pt,height=3.1cm,valign=center]
      \tcbox[colback=white,colframe=popink,boxrule=1.3pt,arc=5pt,boxsep=0pt,nobeforeafter,
        left=3pt,right=3pt,top=3pt,bottom=3pt,valign=center]{\qrcode[height=1.9cm]{https://play.google.com/store/apps/details?id=com.luvvix.onbuch&hl=fr}}%
      \hspace{8pt}\begin{minipage}[c]{0.5\linewidth}
        {\popdisplay\fontsize{11}{12.5}\selectfont\color{white}\MakeUppercase{Télécharge OnBuch}}\par\vspace{4pt}
        {\raggedright\popsemi\fontsize{8.4}{11}\selectfont\color{white}Gratuit \textbullet\ Google Play\\Tuteur IA, annales, concours, résultats\par}
      \end{minipage}
    \end{tcolorbox}
  \end{minipage}\hfill
  \begin{minipage}[t]{0.487\linewidth}
    \begin{tcolorbox}[enhanced,colback=poppurple,colframe=popink,boxrule=2.2pt,arc=10pt,
        drop shadow=popink,left=11pt,right=11pt,top=10pt,bottom=10pt,height=3.1cm,valign=center]
      \tikz[baseline=-0.7ex]\node[circle,fill=white,draw=popink,line width=1.3pt,minimum size=1.6cm,inner sep=0]{\popdisplay\fontsize{24}{24}\selectfont\color{poppurple}+};%
      \hspace{8pt}\begin{minipage}[c]{0.6\linewidth}
        {\popdisplay\fontsize{11}{12.5}\selectfont\color{white}\MakeUppercase{Passe à OnBuch+}}\par\vspace{4pt}
        {\raggedright\popsemi\fontsize{8.4}{11}\selectfont\color{white}Tous les cours signés, Tuteur IA illimité, fiches PDF — onglet OnBuch+ de l'app\par}
      \end{minipage}
    \end{tcolorbox}
  \end{minipage}
  \vspace{8pt}
  \popcontenu
  \vfill
  \signatureonbuch
  \restoregeometry
  \newpage
}

% Rangée « Ce cours contient » (page de garde)
\newcommand{\poptile}[3]{% #1 couleur #2 gros texte #3 légende
  \begin{tcolorbox}[enhanced,colback=#1,colframe=popink,boxrule=2pt,arc=9pt,shadow={2.5pt}{-2.5pt}{0pt}{popink},
      left=6pt,right=6pt,top=9pt,bottom=9pt,halign=center,valign=center,height=1.75cm,width=\linewidth,nobeforeafter]
    {\popdisplay\fontsize{12.5}{13}\selectfont\color{white}\MakeUppercase{#2}}\par\vspace{3pt}
    {\centering\popsemi\fontsize{7.8}{9.5}\selectfont\color{white}#3\par}
  \end{tcolorbox}}
\newcommand{\popcontenu}{%
  \par\noindent{\popxboldls\fontsize{7.5}{7.5}\selectfont\color{popmuted}CE COURS SIGNÉ CONTIENT}\par\vspace{7pt}
  \noindent\begin{minipage}[t]{0.235\linewidth}\poptile{popblue}{Cours}{complet, illustré, pas à pas}\end{minipage}\hfill
  \begin{minipage}[t]{0.235\linewidth}\poptile{poporange}{Méthodes}{et exercices résolus}\end{minipage}\hfill
  \begin{minipage}[t]{0.235\linewidth}\poptile{poppink}{Exercices}{type examen + corrigés}\end{minipage}\hfill
  \begin{minipage}[t]{0.235\linewidth}\poptile{popgreen}{Fiche bilan}{l'essentiel à retenir}\end{minipage}\par}

% Sommaire
\newtcolorbox{sommaire}{enhanced,colback=white,colframe=popink,boxrule=2.2pt,arc=10pt,
  drop shadow=popink,left=18pt,right=18pt,top=13pt,bottom=13pt,before skip=6pt,after skip=20pt,
  before upper={\poppill{Sommaire}{poppurple}{white}\par\vspace{9pt}\popsemi\fontsize{10.6}{14.5}\selectfont\color{popink}}}

% Signature de fin de cours — poussée tout en bas de la dernière page
\newcommand{\findecours}{%
  \par\vfill\nopagebreak
  \begin{center}\popsquiggle{poporange}\end{center}\vspace{4pt}
  \signatureonbuch
}
\AtBeginDocument{\def\llbracket{\mathopen{[\mkern-3.5mu[}}\def\rrbracket{\mathclose{]\mkern-3.5mu]}}} % intervalles d'entiers (glyphe ⟦ absent de la police)
% Glyphes absents de Fira Math : redéfinis à partir de glyphes disponibles
\AtBeginDocument{%
  \def\setminus{\mathbin{\backslash}}\let\smallsetminus\setminus
  \def\vdots{\mathord{\vbox{\baselineskip3.2pt\lineskiplimit0pt\kern4pt\hbox{.}\hbox{.}\hbox{.}}}}%
  \def\ddots{\mathinner{\mkern1mu\raise6.5pt\hbox{.}\mkern2mu\raise3.6pt\hbox{.}\mkern2mu\raise0.7pt\hbox{.}\mkern1mu}}%
  \def\triangle{\mathord{\scalebox{0.9}{$\symup{\Delta}$}}}%
  \def\nabla{\mathord{\raisebox{\depth}{\scalebox{1}[-1]{$\symup{\Delta}$}}}}%
  \def\checkmark{\popcheck{popdark}}%
  \def\star{\mathbin{\ast}}\def\bigstar{\mathord{\ast}}%
  \def\diamond{\mathbin{\scalebox{0.6}{$\Diamond$}}}%
  \def\bigcup{\mathop{\mathchoice{\raisebox{-0.3em}{\scalebox{1.7}{$\cup$}}}{\scalebox{1.25}{$\cup$}}{\cup}{\cup}}\displaylimits}%
  \def\bigcap{\mathop{\mathchoice{\raisebox{-0.3em}{\scalebox{1.7}{$\cap$}}}{\scalebox{1.25}{$\cap$}}{\cap}{\cap}}\displaylimits}%
  \def\lesseqqgtr{\mathrel{\lessgtr}}\def\bigoplus{\mathop{\oplus}\displaylimits}%
}
\providecommand{\ssi}{si et seulement si} % abréviation fréquente
\AtBeginDocument{\providecommand{\wideparen}{\overparen}} % arc de cercle

\begin{document}

\pagegarde{Lesson 1 — Listening: Texts about undergoing a job interview}{Anglais}{Tle A}{Chapter 1 : Family and Social Life}{EN01}{2 h}{Cette leçon de compréhension orale entraîne les élèves à écouter et comprendre des dialogues et monologues sur le thème de l'entretien d'embauche. Elle développe les stratégies d'écoute (repérage du thème, des mots-clés, des chiffres et des noms propres), enrichit le lexique du monde du travail et prépare à la réutilisation de ce lexique dans de courtes productions.
\vspace{4pt}
\begin{multicols}{2}\small
\begin{itemize}
\item À la fin de cette leçon, je sais identifier le thème général d'un enregistrement sur un entretien d'embauche dès la première écoute.
\item À la fin de cette leçon, je sais repérer les mots-clés, les chiffres (âge, salaire, horaires, dates) et les noms propres dans un dialogue d'embauche.
\item À la fin de cette leçon, je sais relever les informations essentielles (poste, qualités, expérience, disponibilité) et les détails secondaires.
\item À la fin de cette leçon, je sais répondre à des questions de compréhension de types variés : vrai/faux, QCM et réponses courtes.
\item À la fin de cette leçon, je connais et j'emploie le lexique de l'entretien d'embauche : applicant, interviewer, CV, qualifications, salary, to hire...
\item À la fin de cette leçon, je sais distinguer les questions typiques du recruteur et les réponses attendues du candidat.
\end{itemize}\end{multicols}}

\begin{sommaire}
\begin{multicols}{2}
\begin{enumerate}
\item Mise en route et anticipation
\item Le texte support : script d'écoute
\item Compréhension : questions variées
\item Lexique thématique : le monde du recrutement
\item Méthode d'écoute et exemple traité
\item Production guidée : réemploi du lexique
\item Activité d'intégration
\item Méthodes et exercices résolus
\item Exercices
\item Corrigés des exercices
\item Fiche bilan
\end{enumerate}
\end{multicols}
\end{sommaire}

\begin{objectifs}
\begin{itemize}
\item À la fin de cette leçon, je sais identifier le thème général d'un enregistrement sur un entretien d'embauche dès la première écoute.
\item À la fin de cette leçon, je sais repérer les mots-clés, les chiffres (âge, salaire, horaires, dates) et les noms propres dans un dialogue d'embauche.
\item À la fin de cette leçon, je sais relever les informations essentielles (poste, qualités, expérience, disponibilité) et les détails secondaires.
\item À la fin de cette leçon, je sais répondre à des questions de compréhension de types variés : vrai/faux, QCM et réponses courtes.
\item À la fin de cette leçon, je connais et j'emploie le lexique de l'entretien d'embauche : applicant, interviewer, CV, qualifications, salary, to hire...
\item À la fin de cette leçon, je sais distinguer les questions typiques du recruteur et les réponses attendues du candidat.
\item À la fin de cette leçon, je sais réemployer le lexique appris dans de courtes phrases personnelles sur mon projet professionnel.
\item À la fin de cette leçon, je sais anticiper le contenu d'un enregistrement à partir du titre et des questions avant l'écoute.
\end{itemize}
\end{objectifs}

\begin{prerequis}
\begin{itemize}
\item Le lexique de base des métiers et professions vu en classe de Première (job, worker, teacher, nurse...)
\item Les questions en Wh- (who, what, where, when, why, how) et l'inversion du sujet
\item Le present simple et le present perfect pour parler d'expérience (I have worked...)
\item Les nombres, les dates et les heures en anglais
\item La prise de notes rapide pendant l'écoute (abréviations, mots-clés)
\end{itemize}
\end{prerequis}

\newpage

\coursec{Mise en route et anticipation}

\courssub{Situation d'accroche et problématique}

Imaginez : après le Bac, vous postulez pour un petit job d'été dans une ONG à Yaoundé ou à Buea. Le recruteur vous accueille : \eng{Good morning! Please, tell me about yourself.} Votre cœur bat fort, les mots se bousculent... Au Royaume-Uni, aux États-Unis, au Nigeria ou au Ghana, l'entretien d'embauche (\eng{job interview}) est une épreuve codifiée que tout jeune diplômé doit affronter. Savoir écouter et comprendre un entretien, c'est la première étape pour un jour le réussir soi-même. Cette leçon vous entraîne à capter l'essentiel d'un dialogue d'embauche.

\courssub{Brainstorming : qu'est-ce qu'un entretien d'embauche ?}

Avant d'écouter, activons nos connaissances. Un \cle{job interview} est un échange structuré entre un employeur (ou son représentant) et un candidat. On y évalue les compétences, la motivation et l'adéquation avec le poste.

\begin{definition}[Job interview]
A \textbf{job interview} is a formal meeting in which an employer asks questions to a person who wants a job, in order to decide if that person is suitable for the position. \\
Un \emph{entretien d'embauche} est une rencontre formelle au cours de laquelle un employeur pose des questions à une personne qui souhaite un emploi, afin de décider si elle convient au poste.
\end{definition}

\begin{center}
\begin{tabularx}{\linewidth}{|X|X|X|}
\hline
\rowcolor{popblueL} \textbf{Français} & \textbf{English} & \textbf{Remarque} \\ \hline
Entretien d'embauche & Job interview / Interview & \eng{Interview} seul suffit souvent. \\ \hline
Candidat / Postulant & Applicant / Candidate / Interviewee & \eng{Interviewee} = celui qui subit l'entretien (suffixe \cle{-ee}). \\ \hline
Recruteur / Employeur & Interviewer / Employer / Recruiter & \eng{Interviewer} = celui qui mène l'entretien (suffixe \cle{-er}). \\ \hline
Jury de recrutement & Interview panel / Selection panel & Fréquent dans la fonction publique et les ONG. \\ \hline
CV & CV / Resume (US) & \eng{Resume} se prononce « ré-zu-mé ». \\ \hline
Lettre de motivation & Cover letter & Indispensable en pays anglophones. \\ \hline
\end{tabularx}
\end{center}

\courssub{Observation d'une situation : un jeune Camerounais devant un jury}

Regardons la scène ci-dessous. Elle se passe dans une salle de réunion à Douala. Un jeune homme, \eng{Thomas}, vêtu d'une chemise bien repassée, fait face à trois personnes assises derrière une table : un \eng{HR manager}, un \eng{technical supervisor} et un \eng{team leader}. Sur la table : son \eng{CV}, une \eng{cover letter} et un verre d'eau.

\begin{popfigure}
\begin{tikzpicture}[
  mindmap,
  concept color=popblue,
  level 1/.style={level distance=3.5cm, sibling angle=90},
  level 2/.style={level distance=2.8cm, sibling angle=45},
  every node/.style={font=\small},
  text=white
]
\node[concept] {JOB INTERVIEW}
  child[concept color=poporange] { node[concept] {PEOPLE}
    child { node[concept, scale=0.6] {Interviewer} }
    child { node[concept, scale=0.6] {Applicant / Interviewee} }
    child { node[concept, scale=0.6] {Panel members} }
  }
  child[concept color=popgreen] { node[concept] {DOCUMENTS}
    child { node[concept, scale=0.6] {CV / Resume} }
    child { node[concept, scale=0.6] {Cover letter} }
    child { node[concept, scale=0.6] {Application form} }
    child { node[concept, scale=0.6] {Certificates} }
  }
  child[concept color=poppurple] { node[concept] {QUESTIONS}
    child { node[concept, scale=0.6] {Experience} }
    child { node[concept, scale=0.6] {Skills / Qualifications} }
    child { node[concept, scale=0.6] {Strengths / Weaknesses} }
    child { node[concept, scale=0.6] {Motivation} }
  }
  child[concept color=poppink] { node[concept] {PRACTICAL DETAILS}
    child { node[concept, scale=0.6] {Salary / Wages} }
    child { node[concept, scale=0.6] {Working hours} }
    child { node[concept, scale=0.6] {Start date} }
    child { node[concept, scale=0.6] {Contract type} }
  };
\end{tikzpicture}
\legende{Carte mentale lexicale : les éléments clés d'un entretien d'embauche}
\end{popfigure}

\courssub{Prédiction à partir du titre \eng{Undergoing a job interview}}

Le titre de la leçon est \eng{Undergoing a job interview}. Le verbe \eng{to undergo} (irrégulier : undergo, underwent, undergone) signifie « subir », « passer par », « traverser » une épreuve. Cela suggère que l'entretien est un moment de tension, une étape à franchir.

\begin{methode}[Comment anticiper le contenu d'un enregistrement]
\begin{enumerate}
\item \textbf{Lisez le titre} et soulignez les mots-clés (\eng{job interview}, \eng{undergoing}).
\item \textbf{Lisez les questions} de compréhension si elles sont données (mots-clés : \eng{name}, \eng{qualifications}, \eng{salary}...).
\item \textbf{Prédisez 3 à 5 mots ou expressions} que vous allez probablement entendre (ex. \eng{experience}, \eng{skills}, \eng{teamwork}, \eng{availability}).
\item \textbf{Visualisez la scène} : qui parle ? où ? dans quel but ?
\end{enumerate}
\end{methode}

Appliquons la méthode. En groupe, élaborons une liste de questions que le recruteur posera presque certainement :

\begin{center}
\begin{tabularx}{\linewidth}{|X|X|}
\hline
\rowcolor{popblueL} \textbf{Thème probable} & \textbf{Exemples de questions attendues} \\ \hline
Identité & \eng{What is your full name? / How old are you?} \\ \hline
Formation & \eng{What qualifications do you have? / Tell me about your education.} \\ \hline
Expérience & \eng{Have you ever worked before? / Describe your previous experience.} \\ \hline
Compétences & \eng{What are your main skills? / What languages do you speak?} \\ \hline
Motivation & \eng{Why do you want this job? / Why should we hire you?} \\ \hline
Disponibilité & \eng{When can you start? / Are you available on weekends?} \\ \hline
Salaire & \eng{What are your salary expectations?} \\ \hline
Centres d'intérêt & \eng{What are your hobbies? / What do you do in your free time?} \\ \hline
\end{tabularx}
\end{center}

\courssub{Objectifs d'écoute : sens global puis détails}

Dans cette leçon, nous allons travailler deux niveaux de compréhension :

\begin{aretenir}[Objectifs d'écoute]
\begin{itemize}
\item \textbf{Écoute globale (gist listening)} : comprendre le sujet principal, le ton, le résultat (le candidat est-il accepté ?). On n'écoute pas chaque mot.
\item \textbf{Écoute détaillée (listening for detail)} : relever des informations précises : un chiffre (salaire, âge), un nom propre (nom de l'entreprise, nom du candidat), une date, un mot de vocabulaire ciblé.
\end{itemize}
\end{aretenir}

\begin{attention}[Piège fréquent : vouloir tout comprendre]
Ne cherchez \textbf{pas} à comprendre chaque mot. En examen comme dans la vie réelle, \textbf{le message passe par les mots-clés} (noms, verbes, chiffres). Les mots-outils (\eng{the}, \eng{a}, \eng{is}, \eng{well}, \eng{you know}) peuvent être ignorés. Concentrez-vous sur les \cle{mots de contenu} (content words) qui portent le sens.
\end{attention}

\courssub{Le savais-tu ? : politesse et codes culturels}

\begin{savaistu}[Salutations et posture en entretien]
Au Royaume-Uni et aux États-Unis, on serre la main fermement (\eng{to shake hands firmly}) en regardant son interlocuteur dans les yeux (\eng{eye contact}) et on attend d'être invité à s'asseoir (\eng{Please, have a seat / Take a seat}). Au Cameroun anglophone (Buea, Bamenda, Limbe), la politesse envers le \eng{panel} est tout aussi codifiée : on salue par ordre de préséance, on utilise \eng{Sir} ou \eng{Madam}, et on évite de croiser les bras (signe de fermeture). Le \eng{small talk} (bavardage d'entrée : \eng{How was your journey?}, \eng{Nice weather today}) est fréquent pour détendre l'atmosphère avant les questions techniques.
\end{savaistu}

\courssub{Attention : \eng{Interviewer} vs \eng{Interviewee}}

C'est l'erreur n°1 des francophones. En anglais, le suffixe \cle{-er} (ou \cle{-or}) désigne l'agent (celui qui fait l'action), tandis que le suffixe \cle{-ee} désigne le patient (celui qui subit l'action).

\begin{attention}[Interviewer $\neq$ Interviewee]
\begin{itemize}
\item \eng{The \textbf{interviewer} asks the questions.} (Le recruteur pose les questions.) $\rightarrow$ Agent.
\item \eng{The \textbf{interviewee} answers the questions.} (Le candidat répond aux questions.) $\rightarrow$ Patient.
\end{itemize}
Même logique : \eng{employer} (employeur) / \eng{employee} (employé) ; \eng{payer} (payeur) / \eng{payee} (bénéficiaire) ; \eng{trainer} (formateur) / \eng{trainee} (stagiaire).
\end{attention}

\courssub{Exercice résolu : identifier les rôles}

\begin{exoresolu}[Qui fait quoi ?]
Dans les phrases ci-dessous, soulignez le sujet et dites s'il s'agit de l'\eng{interviewer} ou de l'\eng{interviewee}.
\begin{enumerate}
\item \eng{The \textbf{interviewer} welcomed the candidate with a smile.}
\item \eng{The nervous \textbf{interviewee} forgot to mention his computer skills.}
\item \eng{``Tell me about yourself,'' said the \textbf{interviewer}.}
\item \eng{The \textbf{interviewee} asked about the working hours.}
\end{enumerate}
\tcblower
\textbf{Solution détaillée :}
\begin{enumerate}
\item \eng{The interviewer} (sujet) $\rightarrow$ c'est celui qui accueille : \textbf{interviewer} (recruteur).
\item \eng{The nervous interviewee} (sujet) $\rightarrow$ c'est celui qui est nerveux et oublie : \textbf{interviewee} (candidat).
\item \eng{the interviewer} (sujet de \eng{said}) $\rightarrow$ c'est celui qui parle/ordonne : \textbf{interviewer}.
\item \eng{The interviewee} (sujet) $\rightarrow$ c'est celui qui pose une question à son tour : \textbf{interviewee}.
\end{enumerate}
\end{exoresolu}

\courssub{Texte support d'anticipation : « First Impressions »}

Avant d'aborder l'écoute principale, lisons ce court texte original qui met en scène Thomas (notre jeune Camerounais) juste avant d'entrer dans la salle. Il illustre le vocabulaire clé et la tension du moment.

\begin{textebox}[First Impressions — Original text]
Thomas adjusted his tie for the third time while waiting in the corridor of the NGO headquarters in Douala. His CV and cover letter were neatly arranged in a black folder. He took a deep breath, reminding himself of the advice his elder brother gave him: ``Be confident, make eye contact, and listen carefully before you answer.'' The heavy wooden door opened. A tall woman with a warm smile stepped out. ``Good morning, Thomas. I am Mrs. Ayuk, the HR Manager. Please, come in. The panel is ready for you.'' Thomas stood up, offered a firm handshake, and walked into the room with his head high. Three faces turned towards him. The interview had begun.
\end{textebox}

\begin{center}
\begin{tabularx}{\linewidth}{|l|X|}
\hline
\rowcolor{popblueL} \textbf{Mot / Expression} & \textbf{Traduction / Glossaire} \\ \hline
to adjust & ajuster, arranger \\ \hline
tie & cravate \\ \hline
corridor & couloir \\ \hline
headquarters & siège (social) \\ \hline
neatly & soigneusement, proprement \\ \hline
folder & chemise, dossier \\ \hline
to take a deep breath & prendre une grande respiration \\ \hline
elder brother & grand frère \\ \hline
confident & confiant, sûr de soi \\ \hline
eye contact & contact visuel \\ \hline
heavy wooden door & lourde porte en bois \\ \hline
stepped out & est sortie (verbe \eng{step out}) \\ \hline
firm handshake & poignée de main ferme \\ \hline
with his head high & la tête haute \\ \hline
panel & jury, commission \\ \hline
\end{tabularx}
\end{center}

\begin{exercice}[Compréhension du texte « First Impressions »]
Répondez en anglais par des phrases courtes.
\begin{enumerate}
\item Where does the scene take place?
\item What documents does Thomas have in his folder?
\item What advice did his brother give him?
\item Who is Mrs. Ayuk?
\item How does Thomas greet Mrs. Ayuk?
\end{enumerate}
\end{exercice}

\begin{corrige}[Exercice n°1]
\begin{enumerate}
\item The scene takes place in the corridor of the NGO headquarters in Douala.
\item He has his CV and cover letter.
\item His brother advised him to be confident, make eye contact and listen carefully before answering.
\item She is the HR Manager.
\item He offers a firm handshake.
\end{enumerate}
\end{corrige}

\courssub{Synthèse de la mise en route}

Vous avez maintenant :
\begin{itemize}
\item identifié les acteurs (\eng{interviewer} / \eng{interviewee}) et les documents clés (\eng{CV}, \eng{cover letter}) ;
\item prédit les questions types (formation, expérience, compétences, salaire) ;
\item compris la stratégie : \textbf{écouter d'abord le sens global, puis traquer les détails} ;
\item noté les pièges linguistiques (\eng{-er} vs \eng{-ee}) et culturels (poignée de main, \eng{eye contact}).
\end{itemize}

Nous sommes prêts à passer à l'écoute du script principal dans la section suivante.

\coursec{Le texte support : script d'écoute}

Après la phase de mise en route où vous avez activé vos connaissances sur le monde professionnel et formulé des hypothèses sur le déroulement d'un entretien, nous allons maintenant étudier le document sonore authentique. Ce script reproduit fidèlement l'enregistrement que vous allez écouter. Il met en scène un échange réaliste entre une responsable des ressources humaines d'une ONG locale et un jeune diplômé camerounais. Observez attentivement la structure de la conversation, le registre de langue (formel mais accessible) et les informations factuelles (chiffres, noms, dates) qui sont les cibles principales de l'exercice de compréhension.

\courssub{Présentation du contexte et des personnages}

\begin{definition}[Contexte de l'entretien]
L'entretien se déroule au siège de la \textbf{Hope Foundation}, une organisation non gouvernementale (ONG) basée à \textbf{Buea}, chef-lieu de la région du Sud-Ouest du Cameroun. \textbf{Mrs Bih} est la \emph{Human Resources Manager} (Directrice des Ressources Humaines). Elle reçoit \textbf{Tabi Ngong}, âgé de 22 ans, diplômé en Sociologie de l'\textbf{University of Buea}, candidat au poste de \emph{Project Assistant} (Assistant de projet). L'échange suit le protocole standard d'un entretien de recrutement formel en contexte anglophone.
\end{definition}

\begin{center}
\begin{tabularx}{\linewidth}{|X|X|X|}
\hline
\rowcolor{popblueL}
\textbf{Élément} & \textbf{Détail} & \textbf{Rôle dans l'écoute} \\
\hline
Lieu & Hope Foundation, Buea & Repère spatial (nom propre à noter) \\
Recruteur & Mrs Bih (HR Manager) & Locutrice principale, pose les questions \\
Candidat & Tabi Ngong & Locuteur secondaire, apporte les infos factuelles \\
Poste visé & Project Assistant & Thème central du dialogue \\
\hline
\end{tabularx}
\end{center}

\courssub{Script intégral de l'enregistrement}

Écoutez d'abord l'enregistrement sans prendre de notes pour vous imprégner de l'intonation et du rythme. Ensuite, lisez le script ci-dessous pour confirmer votre compréhension globale.

\begin{textebox}[A JOB INTERVIEW AT HOPE FOUNDATION]
Mrs Bih: Good morning, Mr Ngong. Please sit down.
Tabi: Good morning, Madam. Thank you.
Mrs Bih: So, you are applying for the post of project assistant. Can you tell me about yourself?
Tabi: Well, my name is Tabi Ngong. I am twenty-two years old and I come from Bamenda. I graduated from the University of Buea last year with a degree in Sociology.
Mrs Bih: Interesting. Do you have any work experience?
Tabi: Yes, I have worked as a volunteer with the Red Cross for two years. I helped to organise health campaigns in rural areas.
Mrs Bih: Very good. What are your strengths?
Tabi: I am hard-working, punctual and I speak English, French and Pidgin fluently.
Mrs Bih: The job pays 150,000 CFA francs a month, from 8 a.m. to 4 p.m., Monday to Friday. When could you start?
Tabi: I could start on the first of September.
Mrs Bih: Fine. We will call you next week. Thank you for coming.
Tabi: Thank you very much, Madam. Have a nice day.
\end{textebox}

\courssub{Analyse structurelle du dialogue}

Le dialogue suit une progression logique classique (\emph{canonical structure}) que tout candidat doit maîtriser pour anticiper les questions. Voici l'organigramme du déroulement :

\begin{popfigure}
\begin{tikzpicture}[
    node distance=1.2cm and 1.5cm,
    box/.style={rectangle, rounded corners, draw=popblue, thick, fill=popblueL, text width=2.8cm, align=center, minimum height=1.2cm, font=\small},
    arrow/.style={-Stealth, thick, popdark}
]
\node[box, align=center] (greet){Greeting \&\\ Introduction};
\node[box, below=of greet, align=center] (intro){Présentation\\ du candidat\\ (Background)};
\node[box, below=of intro, align=center] (exp){Expérience \&\\ Compétences\\ (Skills)};
\node[box, below=of exp, align=center] (qual){Qualités personnelles\\ (Strengths)};
\node[box, below=of qual, align=center] (prac){Détails pratiques\\ Salaire, Horaires,\\ Date de début};
\node[box, below=of prac, align=center] (close){Conclusion \&\\ Prochaines étapes};

\draw[arrow] (greet) -- (intro);
\draw[arrow] (intro) -- (exp);
\draw[arrow] (exp) -- (qual);
\draw[arrow] (qual) -- (prac);
\draw[arrow] (prac) -- (close);
\end{tikzpicture}
\legende{Organigramme du déroulement standard d'un entretien d'embauche}
\end{popfigure}

\courssub{Glossaire thématique : le vocabulaire du recrutement}

Les mots ci-dessous sont essentiels pour comprendre le script et pour répondre aux questions de compréhension. Apprenez-les avec leur traduction et leur prononciation approximative (entre guillemets français). Les mots-clés du thème sont mis en évidence.

\begin{definition}[Glossaire thématique — Vocabulaire clé de l'entretien]
\begin{tabular}{@{}ll@{}}
\cle{to apply for} [tu é-plaï for] & : postuler à \\
\cle{post} [poust] & : poste, emploi \\
\cle{to graduate} [tu græ-djou-eït] & : obtenir son diplôme (université) \\
\cle{degree} [di-gri] & : diplôme universitaire (licence, master) \\
\cle{volunteer} [vo-lon-ti-eu] & : bénévole (nom) ; faire du bénévolat (verbe) \\
\cle{strengths} [stren(g)ths] & : points forts, atouts (antonyme : \eng{weaknesses} = faiblesses) \\
\cle{hard-working} [hard wé-rou-king] & : travailleur, acharné au travail \\
\cle{punctual} [ponk-tchou-el] & : ponctuel (qui arrive à l'heure) \\
\cle{fluently} [flou-ent-li] & : couramment, avec aisance \\
\cle{to pay} [tu peï] & : payer (un salaire) ; \eng{the job pays...} = le poste rapporte... \\
\cle{per month} [peu manth] & : par mois \\
\cle{from ... to ...} [from ... tou ...] & : de ... à ... (horaires) \\
\cle{to start} [tu start] & : commencer, débuter (le travail)
\end{tabular}
\end{definition}

\begin{propriete}[Repérage des informations chiffrées et noms propres — Cibles prioritaires]
Dans un exercice de \emph{listening} de type « job interview », les informations factuelles sont les plus faciles à noter et rapportent le plus de points. Entraînez-vous à les isoler immédiatement :
\begin{itemize}
    \item \textbf{Âge :} \eng{twenty-two years old} $\rightarrow$ 22 ans.
    \item \textbf{Origine :} \eng{Bamenda} (ville), \eng{University of Buea} (institution).
    \item \textbf{Durée expérience :} \eng{two years} (bénévolat Red Cross).
    \item \textbf{Salaire :} \eng{150,000 CFA francs a month} (attention : \eng{150,000} s'entend « one hundred and fifty thousand »).
    \item \textbf{Horaires :} \eng{8 a.m. to 4 p.m.}, \eng{Monday to Friday}.
    \item \textbf{Date de début :} \eng{first of September} (ou \eng{September 1st}).
    \item \textbf{Délai de réponse :} \eng{next week}.
\end{itemize}
\end{propriete}

\courssub{Méthode : Les trois écoutes stratégiques}

La méthode officielle pour réussir l'épreuve de compréhension de l'oral repose sur trois passages obligatoires. Ne cherchez pas à tout écrire à la première écoute.

\begin{methode}[Stratégie des trois écoutes]
\textbf{1ʳᵉ écoute : Saisir le thème global et identifier les personnages.}
\begin{enumerate}
    \item Ne prenez \textbf{aucune note} détaillée. Notez seulement : Qui parle ? Où ? De quoi s'agit-il globalement ? (Ex. : \emph{Job interview, NGO, Buea, young graduate}).
    \item Repérez le ton : formel, poli, professionnel.
\end{enumerate}

\textbf{2ᵉ écoute : Relever les détails factuels (chiffres, noms, qualités).}
\begin{enumerate}
    \item Utilisez un système d'abréviations personnelles : \emph{150k / mo}, \emph{8h-16h}, \emph{Sept 1st}, \emph{Univ Buea}, \emph{Red Cross 2y}.
    \item Concentrez-vous sur les \emph{Wh-questions} posées par le recruteur (\eng{What...? When...? How much...?}) : les réponses du candidat contiennent l'information.
    \item Cochez les éléments du glossaire que vous entendez.
\end{enumerate}

\textbf{3ᵉ écoute (si temps) : Vérification et complétion.}
\begin{enumerate}
    \item Relisez vos notes, vérifiez la cohérence (le salaire correspond-il au poste ?).
    \item Complétez les trous éventuels sur les qualités (\eng{hard-working, punctual, fluent}) ou la conclusion (\eng{call next week}).
    \item Préparez les réponses aux questions de l'examen (Vrai/Faux, QCM, Réponses courtes).
\end{enumerate}
\end{methode}

\courssub{Focus phonologique : Le piège des nombres en \eng{-teen} vs \eng{-ty}}

C'est l'erreur n°1 des candidats francophones à l'écoute. La différence ne se joue pas sur la voyelle mais sur \textbf{l'accent tonique} (la syllabe forte).

\begin{attention}[Piège auditif : Thirteen (13) vs Thirty (30) — Fifteen (15) vs Fifty (50)]
\begin{itemize}
    \item \textbf{Nombres en \eng{-teen} (13 à 19) :} L'accent tonique tombe sur la \textbf{dernière syllabe} (\textbf{-TEEN}).
    \eng{thir-TEEN} (prononcé « zeur-\textbf{TINE} »), \eng{fif-TEEN} (prononcé « fif-\textbf{TINE} »).
    \item \textbf{Nombres en \eng{-ty} (20, 30, 40...) :} L'accent tonique tombe sur la \textbf{première syllabe}.
    \eng{THIR-ty} (prononcé « \textbf{ZEUR}-ti »), \eng{FIF-ty} (prononcé « \textbf{FIF}-ti »).
\end{itemize}
\textbf{Astuce de survie :} Si vous entendez un son long et aigu à la fin (\eng{...TINE}), c'est un nombre entre 13 et 19. Si le mot est « court » et tombe tout de suite, c'est une dizaine (30, 40, 50). Dans notre script, \eng{150,000} contient \eng{fifty} (\eng{one hundred and \textbf{FIF}-ty thousand}) $\rightarrow$ accent sur \eng{FIF}.
\end{attention}

\courssub{Tableau récapitulatif : Informations clés du script (Pour prise de notes)}

Ce tableau modélise la fiche de notes idéale à construire lors de la 2ᵉ écoute.

\begin{center}
\begin{tabularx}{\linewidth}{|l|X|X|}
\hline
\rowcolor{popblueL}
\textbf{Information demandée} & \textbf{Extrait audio (anglais)} & \textbf{Note prise (abréviation)} \\
\hline
Nom du candidat & Tabi Ngong & Tabi Ngong \\
Âge & twenty-two years old & 22 ans \\
Ville d'origine & Bamenda & Bamenda \\
Diplôme / Université & degree in Sociology, University of Buea & Socio / Univ Buea \\
Expérience & volunteer, Red Cross, two years & Volontaire Croix-Rouge 2 ans \\
Qualités (Strengths) & hard-working, punctual, speak 3 lang. fluently & Travailleur, Ponctuel, 3 langues \\
Salaire mensuel & 150,000 CFA francs a month & 150k FCFA/mois \\
Horaires & 8 a.m. to 4 p.m., Monday to Friday & 8h-16h, Lun-Ven \\
Date de début & first of September & 1er Sept \\
Suite & call you next week & Appel sem. proch. \\
\hline
\end{tabularx}
\end{center}

\courssub{Exercice résolu : Questions de compréhension type}

Voici un exemple de questions exactement comme vous en trouverez à l'examen. Appliquez la méthode : lisez les questions \textbf{avant} la 2ᵉ écoute.

\begin{exoresolu}[Compréhension détaillée du script]
\textbf{Énoncé :} Écoutez le dialogue (ou relisez le script) et répondez aux questions suivantes en anglais ou en français selon la consigne.

\begin{enumerate}
    \item \textbf{True or False?} Justify by quoting the text.
    \begin{itemize}
        \item[a)] Tabi is 25 years old.
        \item[b)] Tabi studied at the University of Yaoundé.
        \item[c)] Tabi has worked for the Red Cross.
        \item[d)] The job is 7 days a week.
    \end{itemize}
    \item \textbf{Multiple Choice.} Choose the correct answer.
    \begin{itemize}
        \item[a)] What is the monthly salary?
            A) 15,000 FCFA \quad B) 150,000 FCFA \quad C) 1,500,000 FCFA
        \item[b)] Which language is \textbf{NOT} mentioned by Tabi?
            A) English \quad B) French \quad C) German \quad D) Pidgin
    \end{itemize}
    \item \textbf{Short Answers.} Answer in full sentences.
    \begin{itemize}
        \item[a)] When can Tabi start working?
        \item[b)] What will Mrs Bih do next week?
    \end{itemize}
\end{enumerate}

\tcblower
\textbf{Solution détaillée :}

\begin{enumerate}
    \item \textbf{True / False + Justification}
    \begin{itemize}
        \item[a)] \textbf{False.} \eng{« I am twenty-two years old. »} (L. 4)
        \item[b)] \textbf{False.} \eng{« I graduated from the University of Buea. »} (L. 5)
        \item[c)] \textbf{True.} \eng{« I have worked as a volunteer with the Red Cross for two years. »} (L. 7-8)
        \item[d)] \textbf{False.} \eng{« Monday to Friday. »} (L. 11) $\rightarrow$ 5 jours, pas 7.
    \end{itemize}
    \item \textbf{Multiple Choice}
    \begin{itemize}
        \item[a)] \textbf{B) 150,000 FCFA.} \eng{« The job pays 150,000 CFA francs a month. »} (L. 10-11). \emph{Piège : ne pas confondre avec 15,000 (fifteen thousand).}
        \item[b)] \textbf{C) German.} Tabi cite : \eng{English, French and Pidgin}. L'allemand n'est pas mentionné.
    \end{itemize}
    \item \textbf{Short Answers (Phrases complètes attendues)}
    \begin{itemize}
        \item[a)] \eng{Tabi can start on the first of September.} (Ou : \eng{He could start on September 1st.})
        \item[b)] \eng{Mrs Bih will call him next week.} (Ou : \eng{They will call Tabi next week.})
    \end{itemize}
\end{enumerate}

\textbf{Conseil de rédaction :} Pour les \emph{Short Answers}, reprenez toujours le sujet de la question et le verbe au bon temps. Évitez les réponses télégraphiques (\emph{« On September 1st »}) sauf si la consigne dit « \emph{Answer briefly} ».
\end{exoresolu}

\courssub{Production guidée : Réemploi immédiat du lexique}

Avant de passer à la section « Compréhension : questions variées », entraînez-vous à réutiliser le vocabulaire actif. Transformez les notes du tableau récapitulatif en phrases complètes à l'oral ou à l'écrit.

\begin{exercice}[Mise en bouche lexicale — À faire sur cahier]
Complétez les phrases suivantes avec un mot du \textbf{Glossaire} (forme correcte requise).
\begin{enumerate}
    \item Tabi \rule{2.5cm}{0.4pt} from the University of Buea last year.
    \item He worked as a \rule{2.5cm}{0.4pt} for the Red Cross.
    \item His main \rule{2.5cm}{0.4pt} are punctuality and fluency in three languages.
    \item The \rule{2.5cm}{0.4pt} pays 150,000 FCFA \rule{2.5cm}{0.4pt} month.
    \item Working hours are \rule{2.5cm}{0.4pt} 8 a.m. \rule{2.5cm}{0.4pt} 4 p.m.
    \item Mrs Bih will \rule{2.5cm}{0.4pt} him next week to give the result.
\end{enumerate}
\end{exercice}

\begin{corrige}[Exercice : Mise en bouche lexicale]
\begin{enumerate}
    \item \textbf{graduated} (passé simple / \emph{past simple} : action terminée dans le passé daté \eng{last year}).
    \item \textbf{volunteer} (nom après \eng{as a}).
    \item \textbf{strengths} (pluriel car \eng{are} + plusieurs qualités listées).
    \item \textbf{job / post} — \textbf{per} (ou \textbf{a}) : \eng{per month} / \eng{a month}.
    \item \textbf{from} — \textbf{to} (structure fixe \eng{from ... to ...}).
    \item \textbf{call} (futur simple \eng{will call} après \eng{will}).
\end{enumerate}
\end{corrige}

\begin{savaistu}[Le contexte camerounais : Le bilinguisme, atout majeur]
Dans le script, Tabi met en avant : \eng{« I speak English, French and Pidgin fluently. »} Au Cameroun, le bilinguisme officiel (anglais/français) est une compétence professionnelle \textbf{exigeante} pour les ONG, l'administration et les entreprises internationales. Le \emph{Pidgin English} (ou \emph{Kamtok}), langue véhiculaire très parlée dans les régions Nord-Ouest et Sud-Ouest, est un atout terrain pour le travail communautaire (\eng{rural areas}, \eng{health campaigns}). Lors d'un entretien, citez toujours vos langues avec le niveau (ex. \eng{fluent}, \eng{working knowledge}, \eng{basic}).
\end{savaistu}

\courssub{Transition vers la section suivante}

Vous maîtrisez maintenant le script, son vocabulaire, sa structure et la méthode pour en extraire l'information. La section suivante (« \emph{Compréhension : questions variées} ») vous proposera une batterie complète d'exercices types examen (Vrai/Faux justifié, QCM pièges, Questions ouvertes, Complétion de tableau) pour vous entraîner en conditions réelles. Prenez le temps de refaire l'exercice résolu ci-dessus sans regarder la correction avant d'y aller.

\coursec{Compréhension : questions variées}

Après avoir découvert le script d'écoute dans la section précédente, nous allons maintenant exploiter ce texte à travers une batterie de questions progressives : de la compréhension globale aux détails précis, en passant par le réemploi lexical. Chaque type de question appelle une stratégie spécifique que vous devez maîtriser pour l'examen.

\courssub{Questions de compréhension globale (Global understanding)}

Ces questions portent sur le \cle{schéma communicationnel} : qui parle ? où ? de quoi ? Elles se répondent dès la première écoute.

\begin{exemplebox}[Questions globales types]
\begin{itemize}
    \item \eng{Who are the speakers?} « Qui sont les interlocuteurs ? »
    \item \eng{Where does the interview take place?} « Où se déroule l'entretien ? »
    \item \eng{What job is discussed?} « Quel poste est concerné ? »
    \item \eng{What is the outcome of the interview?} « Quel est l'issue de l'entretien ? »
\end{itemize}
\end{exemplebox}

\begin{propriete}[Stratégie : la première écoute]
\emph{Ne prenez pas de notes détaillées.} Concentrez-vous sur :
\begin{enumerate}
    \item Les \cle{mots-clés} du thème (interview, candidate, manager, position, salary, experience).
    \item Les \cle{noms propres} (Tabi, Mr. Ndi, Yaoundé, Red Cross).
    \item Les \cle{chiffres} (âge, salaire, années d'expérience, date de début).
    \item Le \cle{ton} général (formel, stressé, confiant, concluant).
\end{enumerate}
Une deuxième écoute servira aux détails.
\end{propriete}

\courssub{Vrai / Faux avec justification (True / False + Justification)}

C'est l'exercice roi du Bac. La réponse « True » ou « False » ne vaut \textbf{aucun point} sans la justification textuelle exacte.

\begin{methode}[Réussir le Vrai/Faux en 4 étapes]
\begin{enumerate}
    \item \textbf{Lisez l'énoncé} avant la deuxième écoute : repérez les mots à surveiller (négation, comparatif, chiffre).
    \item \textbf{Écoutez activement} : quand vous entendez l'information, \textbf{soulignez la phrase exacte} dans votre script (ou notez le time-code si pas de script).
    \item \textbf{Décidez} : True (l'énoncé résume fidèlement le texte) ou False (contraire, absence d'info, déformation).
    \item \textbf{Justifiez par citation} : recopiez le segment du texte qui prouve votre choix. Traduisez-le si la consigne l'exige.
\end{enumerate}
\end{methode}

\begin{exemplebox}[Modèles de Vrai/Faux corrigés]
\textbf{Énoncé 1 :} \eng{Tabi has never worked before.} \\
\textbf{Réponse :} \textbf{FALSE}. \\
\textbf{Justification :} \eng{“I worked as a volunteer with the Red Cross for two years.”} \\
« Il a travaillé comme bénévole à la Croix-Rouge pendant deux ans. »

\textbf{Énoncé 2 :} \eng{The salary offered is 150,000 FCFA per month.} \\
\textbf{Réponse :} \textbf{TRUE}. \\
\textbf{Justification :} \eng{“The starting salary is 150,000 FCFA monthly.”} \\
« Le salaire de départ est de 150 000 FCFA par mois. »

\textbf{Énoncé 3 :} \eng{Mr. Ndi asks Tabi about his family.} \\
\textbf{Réponse :} \textbf{FALSE}. \\
\textbf{Justification :} \eng{“Mr. Ndi: ‘Tell me about your professional experience.’”} (Aucune question sur la famille n'est posée.)
\end{exemplebox}

\begin{attention}[Piège : l'absence d'information ≠ Faux]
Au Bac, si le texte \textbf{ne dit rien} sur un point, l'énoncé est \textbf{FALSE} (ou « Not mentioned » selon la grille). Ne déduisez jamais.
\begin{itemize}
    \item Texte : \eng{“I speak English and French.”}
    \item Énoncé : \eng{“Tabi does not speak German.”} $\rightarrow$ \textbf{FALSE} (information absente, pas contredite).
\end{itemize}
De même, \textbf{attention aux négations} : \eng{“I don't have a driving licence.”} $\neq$ \eng{“I have a driving licence.”}
\end{attention}

\courssub{QCM sur détails chiffrés (Multiple Choice : numbers, dates, amounts)}

Les distracteurs (fausses options) jouent sur la \cle{proximité phonétique} ou l'\cle{ordre de grandeur}.

\begin{exemplebox}[QCM type — Salaire]
\textbf{Question :} \eng{How much does the job pay per month?} \\
\begin{itemize}
    \item[a)] 15,000 FCFA
    \item[b)] \textbf{150,000 FCFA}
    \item[c)] 500,000 FCFA
    \item[d)] 1,500,000 FCFA
\end{itemize}
\textbf{Bonne réponse : b)} \eng{150,000 FCFA}.
\textbf{Justification audio :} \eng{“The starting salary is one hundred fifty thousand francs CFA.”}
\end{exemplebox}

\begin{exemplebox}[QCM type — Âge et expérience]
\textbf{Question :} \eng{How old is Tabi?} \\
a
\end{exemplebox}


\coursec{Lesson 1 — Listening: Texts about undergoing a job interview}

\courssub{Mise en route et anticipation}

Avant d'aborder l'écoute, mettez-vous en situation. Observez l'image mentale suivante : un jeune diplômé, vêtu d'un costume soigné, attend dans le couloir d'une entreprise à Douala (ex. : SONARA, CDC, ou une banque à Bonanjo). Il tient un dossier sous le bras.

\begin{experience}[Brainstorming guidé]
En classe, répondez oralement à ces questions (en anglais si possible) :
\begin{enumerate}
    \item How does the candidate feel? (Stressed, confident, nervous, hopeful...)
    \item What documents are in his folder? (CV, cover letter, diplomas, certificates, ID card...)
    \item What questions will the interviewer likely ask? (Tell me about yourself. Why do you want this job? What are your strengths/weaknesses? What are your salary expectations?)
    \item What qualities does the employer look for? (Punctual, hard-working, reliable, skilled, fluent in English/French...)
\end{enumerate}
\par\medskip
\textbf{Objectif :} Activer votre schème de connaissances sur le recrutement et anticiper le vocabulaire de l'enregistrement.
\end{experience}

\courssub{Script d'écoute : Entretien d'embauche à Douala}

Voici le script de l'enregistrement que vous allez entendre (ou lire). C'est un dialogue authentique entre un recruteur (\eng{Interviewer}) et une candidate (\eng{Interviewee}) pour un poste d'\eng{Assistant Marketing} dans une entreprise agro-industrielle à Douala.

\begin{textebox}[Audio Script: Job Interview at Cameroon Cocoa Corp.]
\textbf{Interviewer (Mr. Kouam):} Good morning, Miss Mendo. Please, have a seat. Thank you for coming on time.
\textbf{Interviewee (Miss Mendo):} Good morning, Sir. Thank you for the opportunity.
\textbf{Interviewer:} First, let me introduce the panel. I am the HR Manager. Beside me is Ms. Ngono, the Sales Director. We have your CV and cover letter here. You applied for the Marketing Assistant position advertised last month.
\textbf{Interviewee:} Yes, Sir. I hold a Bachelor's degree in Marketing from the University of Douala. I also completed a six-month internship at a cocoa exporting company in Kribi, where I gained experience in market research and client relations.
\textbf{Interviewer:} Good. The job requires someone who is computer-literate and fluent in both official languages. Are you comfortable with digital marketing tools?
\textbf{Interviewee:} Absolutely. I master the Microsoft Office Suite, and I have a certification in Google Analytics. I am fluent in English and French, and I speak basic Duala.
\textbf{Interviewer:} Excellent. This is a full-time position with a three-month probation period. The salary is 250,000 FCFA monthly, plus transport allowance and health insurance. There are real promotion prospects for dynamic employees. Do you have any questions?
\textbf{Interviewee:} Yes, Sir. What are the main challenges the marketing team is currently facing?
\textbf{Interviewer:} We are launching a new product line next quarter. We need someone reliable and proactive to handle the launch campaign. We will shortlist candidates by Friday. If you are selected, we will contact you for a second interview.
\textbf{Interviewee:} I understand. I am very interested in this vacancy and I am ready to attend any further interviews. Thank you for your time.
\textbf{Interviewer:} Thank you, Miss Mendo. We will be in touch.
\end{textebox}

\noindent\textbf{Compréhension de l'écoute (Questions variées) :}

\begin{exercice}[Vrai ou Faux ? Justifiez par une phrase du texte.]
\begin{enumerate}
    \item Miss Mendo arrived late for the interview.
    \item The interview panel consists of two people.
    \item Miss Mendo has a Master's degree in Marketing.
    \item She did an internship in Kribi.
    \item The contract offered is a permanent contract (CDI) immediately.
    \item The monthly salary is 300,000 FCFA.
    \item Miss Mendo speaks three languages.
    \item The company is launching a new product next quarter.
\end{enumerate}
\end{exercice}

\begin{exercice}[QCM : Choisissez la bonne réponse.]
\begin{enumerate}
    \item What position did Miss Mendo apply for?
    \begin{itemize}
        \item[a] HR Manager
        \item[b] Sales Director
        \item[c] Marketing Assistant
        \item[d] Accountant
    \end{itemize}
    \item What specific digital skill does she mention?
    \begin{itemize}
        \item[a] Coding in Python
        \item[b] Google Analytics certification
        \item[c] Graphic Design
        \item[d] Video Editing
    \end{itemize}
    \item What is the duration of the probation period?
    \begin{itemize}
        \item[a] One month
        \item[b] Six months
        \item[c] Three months
        \item[d] One year
    \end{itemize}
\end{enumerate}
\end{exercice}

\begin{exercice}[Réponses courtes : Répondez en phrases complètes.]
\begin{enumerate}
    \item Who are the two interviewers on the panel?
    \item What degree does Miss Mendo hold?
    \item Name two benefits mentioned in the contract.
    \item Why does the company need a proactive person right now?
    \item When will the shortlisting be done?
\end{enumerate}
\end{exercice}


\begin{corrige}[Exercices de compréhension - Script d'écoute]
\textbf{Vrai/Faux :}
\begin{enumerate}
    \item \textbf{False.} \eng{Thank you for coming on time.}
    \item \textbf{True.} \eng{I am the HR Manager. Beside me is Ms. Ngono, the Sales Director.}
    \item \textbf{False.} \eng{I hold a Bachelor's degree in Marketing...}
    \item \textbf{True.} \eng{I also completed a six-month internship at a cocoa exporting company in Kribi...}
    \item \textbf{False.} \eng{This is a full-time position with a three-month probation period.}
    \item \textbf{False.} \eng{The salary is 250,000 FCFA monthly...}
    \item \textbf{True.} \eng{I am fluent in English and French, and I speak basic Duala.}
    \item \textbf{True.} \eng{We are launching a new product line next quarter.}
\end{enumerate}
\textbf{QCM :} 1-c, 2-b, 3-c.
\textbf{Réponses courtes :}
\begin{enumerate}
    \item The interviewers are Mr. Kouam (HR Manager) and Ms. Ngono (Sales Director).
    \item She holds a Bachelor's degree in Marketing.
    \item Transport allowance and health insurance. (Also: promotion prospects).
    \item Because they are launching a new product line next quarter and need someone to handle the launch campaign.
    \item The shortlisting will be done by Friday.
\end{enumerate}
\end{corrige}

\courssub{Méthode d'écoute : Repérer l'essentiel}

\begin{methode}[Stratégies pour réussir l'épreuve de Listening (Bac)]
Suivez ces 4 étapes à chaque écoute :
\begin{enumerate}
    \item \textbf{Exploiter le temps de lecture (1 à 2 min) :} Lisez le titre, les consignes et les questions. Soulignez les \cle{mots-clés} (noms propres, chiffres, verbes d'action, mots interrogatifs \eng{Who, What, When, Where, Why, How}). Prédisez le thème et le vocabulaire attendu.
    \item \textbf{Première écoute (Globalité) :} Ne cherchez pas à tout comprendre. Identifiez le \cle{thème} général, le \cle{ton} (formel/informel), le \cle{nombre d'interlocuteurs} et la \cle{structure} (dialogue, monologue, annonce). Cochez les réponses évidentes (Vrai/Faux, QCM).
    \item \textbf{Deuxième écoute (Détails) :} Concentrez-vous sur les informations manquantes : \cle{chiffres} (dates, salaires, âges, pourcentages), \cle{noms propres} (lieux, entreprises, personnes), \cle{mots de liaison} (\eng{but, however, because, so}) qui signalent des arguments. Complétez les réponses courtes.
    \item \textbf{Vérification (30 sec) :} Relisez vos réponses. Vérifiez la cohérence grammaticale (temps, nombre) et l'orthographe des noms propres entendus. Ne laissez aucune case vide.
\end{enumerate}
\par\medskip
\textbf{Astuce Bac :} Pour les chiffres, écoutez la prononciation : \eng{thirteen (13) / thirty (30)}, \eng{fifteen (15) / fifty (50)}. Notez-les en chiffres arabes dès la 1ère écoute.
\end{methode}

\begin{exemplebox}[Application de la méthode sur le script ci-dessus]
\textbf{Étape 1 (Lecture questions) :} Mots-clés repérés : \eng{panel, degree, internship, salary, probation, launch, shortlist, Friday}.
\textbf{Étape 2 (1ère écoute) :} Thème = Entretien d'embauche. 2 interlocuteurs. Ton formel. Réponses V/F 1, 2, 4, 8 cochées.
\textbf{Étape 3 (2ème écoute) :} Chiffres : \eng{six-month}, \eng{250,000 FCFA}, \eng{three-month}, \eng{next quarter}, \eng{Friday}. Noms propres : \eng{Mr. Kouam, Ms. Ngono, Kribi, Duala, Cameroon Cocoa Corp}. Détails : \eng{Bachelor's}, \eng{Google Analytics}, \eng{transport allowance, health insurance}.
\textbf{Étape 4 :} Vérification orthographe \eng{Kribi, Duala, Kouam, Ngono}.
\end{exemplebox}

\coursec{Lexique thématique : le monde du recrutement}

Après avoir travaillé la compréhension détaillée des textes d'écoute sur les entretiens d'embauche, il est indispensable d'ancrer le vocabulaire spécialisé qui structure ce domaine. Cette section vous permet de passer de la compréhension passive à l'appropriation active du lexique : vous pourrez ainsi non seulement comprendre un recruteur, mais aussi vous présenter avec précision lors d'un futur entretien, que ce soit à Douala, Yaoundé, Buea ou dans un contexte international.

\courssub{Les personnes impliquées (People)}

Observons d'abord comment les acteurs sont désignés dans les extraits écoutés : \eng{The interviewer asked difficult questions.} / \eng{The applicant arrived early.} / \eng{The Human Resources manager welcomed us.}
En anglais, les suffixes \cle{-er} / \cle{-or} (celui qui fait l'action) et \cle{-ee} (celui qui subit l'action) sont productifs pour former les noms d'agent.

\begin{center}
\begin{tabularx}{\linewidth}{|X|X|X|}
\hline
\rowcolor{popblueL} \textbf{Anglais} & \textbf{Français} & \textbf{Remarque / Contexte camerounais} \\
\hline
\cle{interviewer} & l'intervieweur, le recruteur & Celui qui mène l'entretien. \\
\cle{interviewee} & l'interviewé, le candidat & Celui qui passe l'entretien (suffixe \textit{-ee}). \\
\cle{applicant} / \cle{candidate} & le candidat, le postulant & \textit{Applicant} insiste sur l'acte de postuler ; \textit{candidate} sur la sélection. \\
\cle{employer} & l'employeur & Personne physique ou morale qui embauche. \\
\cle{employee} & l'employé, le salarié & Attention à l'orthographe : double \textit{e} final. \\
\cle{Human Resources manager} / \cle{HR manager} & le responsable des ressources humaines & Souvent abrégé \textit{HR} à l'oral. \\
\cle{recruiter} & le recruteur (chasseur de têtes) & Spécialiste du recrutement (cabinet ou interne). \\
\cle{panel} & le jury, la commission & \eng{A panel of three interviewers.} (Un jury de trois recruteurs.) \\
\hline
\end{tabularx}
\end{center}

\begin{exemplebox}[Mise en situation : Présentation des acteurs]
\textbf{Contexte :} Au début d'un entretien à la SONARA (Limbe).
\eng{Good morning. I am Mr. Epie, the \textbf{interviewer}. Beside me sits Ms. Ngoh, our \textbf{Human Resources manager}. We will be assessing you, the \textbf{interviewee}, for the position of maintenance technician.}
\par\medskip
\textit{Traduction :} Bonjour. Je suis M. Epie, l'\textbf{intervieweur}. À mes côtés siège Mme Ngoh, notre \textbf{responsable des ressources humaines}. Nous allons évaluer vous, l'\textbf{interviewé}, pour le poste de technicien de maintenance.
\end{exemplebox}

\courssub{Les documents du recrutement (Documents)}

Les textes d'écoute font souvent référence aux pièces du dossier. Maîtriser ces termes évite les malentendus sur ce qu'il faut apporter le jour J.

\begin{definition}[Vacancy]
\cle{Vacancy} (nom) : un poste vacant, une offre d'emploi disponible.
\par\medskip
\eng{The company advertised two \textbf{vacancies} in the national newspaper.}
\par « L'entreprise a annoncé deux \textbf{postes vacants} dans le journal national. »
\end{definition}

\begin{center}
\begin{tabularx}{\linewidth}{|X|X|X|}
\hline
\rowcolor{poporangeL} \textbf{Anglais} & \textbf{Français} & \textbf{Remarque / Collocation} \\
\hline
\cle{CV} / \cle{curriculum vitae} & curriculum vitæ & Usage \textbf{britannique} (standard au Cameroun). \\
\cle{résumé} & CV (résumé) & Usage \textbf{américain}. Les deux sont acceptés au Bac. \\
\cle{cover letter} & lettre de motivation & \eng{To write a cover letter.} (Rédiger une lettre de motivation.) \\
\cle{application form} & formulaire de candidature & \eng{To fill in / fill out an application form.} \\
\cle{diploma} / \cle{degree} & diplôme & \textit{Degree} = diplôme universitaire (Bachelor's, Master's). \textit{Diploma} = souvent technique/professionnel. \\
\cle{certificate} & certificat, attestation & Ex. : \eng{GCE A Level certificate}, \eng{Baccalaureate certificate}. \\
\cle{reference letter} / \cle{recommendation letter} & lettre de recommandation & Écrite par un ancien employeur ou un professeur. \\
\cle{testimonial} & attestation de service / témoignage & Plus formel, souvent pour la fin de contrat. \\
\hline
\end{tabularx}
\end{center}

\begin{savaistu}[CV ou Résumé ?]
Au Cameroun, comme dans tout l'espace Commonwealth, on utilise \textbf{CV} (Curriculum Vitae). Aux États-Unis, on dit \textbf{résumé} (prononcé « ré-zou-mé »). La différence principale : le CV britannique est souvent plus long (2 pages), détaillé, incluant loisirs et références ; le résumé américain est concis (1 page), centré sur les \textbf{achievements} (réalisations). Pour le Bac, les deux termes sont corrects.
\end{savaistu}

\courssub{Les actions clés du processus (Actions)}

Le recrutement suit une chronologie précise. Les verbes ci-dessous sont les « moteurs » du récit d'un entretien. Observez leurs régimes (prépositions) et leur voix (souvent passive pour le candidat).

\begin{center}
\begin{tabularx}{\linewidth}{|X|X|X|}
\hline
\rowcolor{popgreenL} \textbf{Verbe (Action)} & \textbf{Sens / Régime} & \textbf{Exemple traduit} \\
\hline
\cle{to advertise a vacancy} & publier une offre & \eng{The bank \textbf{advertised a vacancy} for a cashier.} (La banque a publié une offre pour caissier.) \\
\cle{to apply for a job / a post} & postuler à un emploi (\textbf{for}) & \eng{She \textbf{applied for the post} of secretary.} (Elle a postulé au poste de secrétaire.) \\
\cle{to fill in / fill out a form} & remplir un formulaire & \eng{Please \textbf{fill in} the application form in capital letters.} (Remplissez en majuscules.) \\
\cle{to shortlist} & présélectionner (faire une liste restreinte) & \eng{They \textbf{shortlisted} five candidates.} (Ils ont présélectionné cinq candidats.) \\
\cle{to interview} & faire passer un entretien & \eng{He was \textbf{interviewed} by a panel of three.} (Il a passé un entretien devant un jury de trois.) \\
\cle{to hire} / \cle{to recruit} & embaucher, recruter & \eng{He was \textbf{hired} last month.} (Il a été embauché le mois dernier.) \\
\cle{to appoint} & nommer (poste officiel/responsabilité) & \eng{She was \textbf{appointed} manager.} (Elle a été nommée manager.) \\
\cle{to reject} / \cle{to turn down} & refuser une candidature / refuser une offre & \eng{My application was \textbf{rejected}.} / \eng{He \textbf{turned down} the offer.} (Il a décliné l'offre.) \\
\cle{to resign} & démissionner & \eng{He \textbf{resigned} from his post.} (Il a démissionné de son poste.) \\
\cle{to retire} & prendre sa retraite & \eng{She will \textbf{retire} next year.} (Elle partira à la retraite l'an prochain.) \\
\hline
\end{tabularx}
\end{center}

\begin{propriete}[Collocations indispensables (Associations de mots figées)]
Ces blocs lexicaux s'apprennent par cœur pour parler naturellement.
\begin{itemize}
    \item \eng{to \textbf{gain experience}} (acquérir de l'expérience) — \textbf{pas} \eng{*win experience} ni \eng{*earn experience}.
    \item \eng{to \textbf{meet the requirements}} (répondre aux exigences / satisfaire aux critères).
    \item \eng{to \textbf{hold a degree / a diploma}} (être titulaire d'un diplôme) — plus formel que \eng{have}.
    \item \eng{to \textbf{submit an application}} (déposer une candidature).
    \item \eng{to \textbf{attend an interview}} (se présenter à un entretien / assister à).
    \item \eng{to \textbf{sign a contract}} (signer un contrat).
\end{itemize}
\end{propriete}

\begin{exemplebox}[Verbes au passé composé / passif : récits d'expérience]
\begin{itemize}
    \item \eng{I \textbf{applied for} the job online last week.} (J'ai postulé en ligne la semaine dernière.)
    \item \eng{My application \textbf{was shortlisted}.} (Ma candidature a été présélectionnée.)
    \item \eng{I \textbf{attended} the interview on Monday.} (Je me suis présenté à l'entretien lundi.)
    \item \eng{I \textbf{was offered} the position.} (On m'a proposé le poste / Le poste m'a été offert.)
    \item \eng{I \textbf{signed} a two-year contract.} (J'ai signé un contrat de deux ans.)
\end{itemize}
\end{exemplebox}

\courssub{Les qualités professionnelles recherchées (Qualities)}

Les recruteurs utilisent des adjectifs précis pour décrire le profil idéal. Notez la formation des contraires par préfixes (\cle{un-}, \cle{in-}, \cle{ir-}, \cle{dis-}).

\begin{center}
\begin{tabularx}{\linewidth}{|X|X|X|}
\hline
\rowcolor{poppurpleL} \textbf{Qualité (Adjectif)} & \textbf{Contraire} & \textbf{Exemple en contexte} \\
\hline
\cle{hard-working} / \cle{industrious} & lazy / idle & \eng{We need a \textbf{hard-working} team member.} \\
\cle{punctual} & late / unpunctual & \eng{Be \textbf{punctual} for the interview.} \\
\cle{reliable} / \cle{dependable} & unreliable & \eng{He is a \textbf{reliable} employee; he never misses a deadline.} \\
\cle{honest} / \cle{trustworthy} & dishonest & \eng{Honesty is a core value in our company.} \\
\cle{dynamic} / \cle{proactive} & passive / slow & \eng{We are looking for a \textbf{dynamic} sales representative.} \\
\cle{skilled} / \cle{qualified} & unskilled / unqualified & \eng{A \textbf{skilled} technician is required.} \\
\cle{experienced} & inexperienced / green & \eng{An \textbf{experienced} driver with a clean licence.} \\
\cle{computer-literate} & computer-illiterate & \eng{Applicants must be \textbf{computer-literate}.} \\
\cle{fluent in English/French} & \textit{(pas d'adjectif unique)} & \eng{Fluent in both official languages.} \\
\hline
\end{tabularx}
\end{center}

\begin{propriete}[Familles de mots : Employ \& Apply]
Connaître les dérivés permet de varier les phrases (nom, verbe, adjectif) et d'éviter les répétitions.
\begin{center}
\begin{tabular}{|c|c|c|c|}
\hline
\rowcolor{popblueL} \textbf{Verbe} & \textbf{Nom (Agent)} & \textbf{Nom (Action/État)} & \textbf{Adjectif / Participe} \\
\hline
to \textbf{employ} & \textbf{employer} (employeur) & \textbf{employment} (emploi) / \textbf{unemployment} (chômage) & \textbf{employed} / \textbf{unemployed} \\
& \textbf{employee} (employé) & & \textbf{employable} (employable) \\
\hline
to \textbf{apply} & \textbf{applicant} (candidat) & \textbf{application} (candidature) & \textbf{applicable} (applicable) \\
& & & \textbf{applied} (appliqué - ex. \eng{applied linguistics}) \\
\hline
\end{tabular}
\end{center}
\eng{Unemployment is high among youth.} (Le chômage est élevé chez les jeunes.) \\
\eng{She sent her \textbf{application} yesterday.} (Elle a envoyé sa candidature hier.)
\end{propriete}

\courssub{Les conditions de travail (Conditions)}

Ces noms sont essentiels pour négocier ou comprendre une offre de contrat.

\begin{center}
\begin{tabularx}{\linewidth}{|X|X|X|}
\hline
\rowcolor{poppinkL} \textbf{Terme} & \textbf{Sens} & \textbf{Précision / Exemple} \\
\hline
\cle{salary} & salaire (mensuel/annuel, cadre) & \eng{A monthly \textbf{salary} of 300,000 FCFA.} \\
\cle{wages} & salaire (horaire/hebdo, ouvrier) & \eng{Hourly \textbf{wages}.} (Salaire horaire.) \\
\cle{full-time} & temps plein & \eng{A \textbf{full-time} position.} (Poste à temps plein.) \\
\cle{part-time} & temps partiel & \eng{She works \textbf{part-time}.} (Elle travaille à mi-temps.) \\
\cle{contract} & contrat & \eng{A \textbf{permanent contract} (CDI) / A \textbf{fixed-term contract} (CDD).} \\
\cle{working hours} & heures de travail & \eng{Standard \textbf{working hours}: 8am–4pm.} \\
\cle{shift work} & travail par roulement / équipes & \eng{Night \textbf{shift}.} (Équipe de nuit.) \\
\cle{promotion} & promotion, avancement & \eng{Good \textbf{promotion} prospects.} (Perspectives d'avancement.) \\
\cle{benefits} / \cle{perks} & avantages (sociaux, primes) & \eng{Fringe \textbf{benefits}: housing, transport, health insurance.} \\
\cle{probation period} & période d'essai & \eng{A three-month \textbf{probation period}.} \\
\hline
\end{tabularx}
\end{center}

\courssub{Synthèse visuelle : Carte mentale du recrutement}

\begin{popfigure}
\begin{tikzpicture}[
    mindmap,
    grow cyclic,
    every node/.style=concept,
    concept color=popblue,
    level 1/.append style={level distance=4.5cm, sibling angle=72},
    level 2/.append style={level distance=3cm, sibling angle=45},
    text=white,
    font=\small\bfseries
]
\node {RECRUITMENT}
    child[concept color=poporange] { node {PEOPLE}
        child { node {Interviewer} }
        child { node {Interviewee} }
        child { node {Applicant} }
        child { node {Employer} }
        child { node {HR Manager} }
    }
    child[concept color=popgreen] { node {DOCUMENTS}
        child { node {CV / Résumé} }
        child { node {Cover Letter} }
        child { node {Application Form} }
        child { node {Degree / Diploma} }
        child { node {Reference Letter} }
    }
    child[concept color=poppurple] { node {ACTIONS}
        child { node {Apply for} }
        child { node {Advertise} }
        child { node {Shortlist} }
        child { node {Interview} }
        child { node {Hire / Reject} }
    }
    child[concept color=poppink] { node {QUALITIES}
        child { node {Hard-working} }
        child { node {Reliable} }
        child { node {Skilled} }
        child { node {Experienced} }
        child { node {Computer-literate} }
    }
    child[concept color=popgold] { node {CONDITIONS}
        child { node {Salary / Wages} }
        child { node {Full-time / Part-time} }
        child { node {Contract (CDI/CDD)} }
        child { node {Working Hours} }
        child { node {Promotion} }
    };
\end{tikzpicture}
\legende{Carte mentale lexicale : les 5 piliers du vocabulaire du recrutement.}
\end{popfigure}

\courssub{Texte d'entraînement : Témoignage d'un jeune diplômé}

Le texte suivant réemploie l'essentiel du vocabulaire vu ci-dessus. Lisez-le activement : soulignez les mots des 5 champs lexicaux.

\begin{textebox}[Témoignage : Mon premier entretien à Douala]
My name is Arnaud and I recently graduated with a \textbf{Bachelor's degree} in Management from the University of Douala. Like many young graduates, I faced the challenge of \textbf{unemployment} for six months. I spent hours scanning newspapers and websites like \textit{Emploi.cm} for \textbf{vacancies}. 

Last March, a major logistics company \textbf{advertised a vacancy} for a \textbf{Junior Project Officer}. I decided to \textbf{apply for the post}. I carefully \textbf{filled in} the online \textbf{application form}, attached my \textbf{CV} and a tailored \textbf{cover letter} highlighting my internship experience. I made sure to \textbf{meet the requirements} listed in the job description: being \textbf{computer-literate}, \textbf{fluent in English and French}, and holding a \textbf{degree} in a relevant field.

Two weeks later, I received an email informing me that my \textbf{application} had been \textbf{shortlisted}. I was invited to \textbf{attend an interview} at their headquarters in Bonanjo. The \textbf{panel} consisted of the \textbf{HR manager} and two technical \textbf{interviewers}. They asked behavioural questions to test if I was \textbf{reliable}, \textbf{dynamic} and able to work under pressure. 

Fortunately, I \textbf{was hired} on the spot! I signed a one-year \textbf{fixed-term contract} with a competitive \textbf{salary} and good \textbf{benefits} (transport allowance, health insurance). There is a three-month \textbf{probation period}, but the \textbf{promotion} prospects are real. My advice to other \textbf{applicants}: prepare your documents well, be \textbf{punctual}, and stay \textbf{honest} during the interview.
\end{textebox}

\noindent\textbf{Glossaire des mots difficiles (selon l'ordre d'apparition) :}
\begin{itemize}
    \item \cle{graduated} (verbe) : être diplômé, obtenir son diplôme.
    \item \cle{scanning} (verbe -ing) : parcourir rapidement (les annonces).
    \item \cle{logistics} (nom) : logistique (secteur d'activité).
    \item \cle{tailored} (adj.) : adapté, sur mesure (pour le poste).
    \item \cle{highlighting} (verbe -ing) : mettre en valeur.
    \item \cle{internship} (nom) : stage professionnel.
    \item \cle{headquarters} (nom pl.) : siège social.
    \item \cle{behavioural questions} (nom) : questions comportementales (sur le vécu).
    \item \cle{under pressure} (loc.) : sous pression.
    \item \cle{on the spot} (loc.) : sur-le-champ, immédiatement.
    \item \cle{allowance} (nom) : indemnité, allocation.
    \item \cle{prospects} (nom pl.) : perspectives (d'avenir).
\end{itemize}

\noindent\textbf{Questions de compréhension (Réponses courtes) :}
\begin{enumerate}
    \item What degree does Arnaud hold?
    \item Where did he look for vacancies?
    \item What documents did he submit for his application?
    \item Who composed the interview panel?
    \item What type of contract did he sign?
    \item What qualities were the interviewers testing?
\end{enumerate}

\begin{corrige}[Questions texte Arnaud]
\begin{enumerate}
    \item He holds a Bachelor's degree in Management.
    \item He looked in newspapers and on websites like Emploi.cm.
    \item He submitted his CV and a tailored cover letter.
    \item The panel consisted of the HR manager and two technical interviewers.
    \item He signed a one-year fixed-term contract.
    \item They were testing if he was reliable, dynamic and able to work under pressure.
\end{enumerate}
\end{corrige}

\courssub{Pièges à éviter : Les faux-amis et erreurs fréquentes}

\begin{attention}[Faux-amis critiques pour le Bac]
\begin{itemize}
    \item \textbf{Formation $\neq$ Formation (professionnelle).}
    \begin{itemize}
        \item \eng{Formation} = \textit{formation géologique, formation militaire, constitution d'un groupe}. 
        \item Pour « formation professionnelle / études » : \eng{training}, \eng{education}, \eng{educational background}, \eng{qualifications}.
        \item \textbf{Incorrect :} \eng{*I have a good formation.} $\rightarrow$ \textbf{Correct :} \eng{I have a good \textbf{training} / \textbf{educational background}.}
    \end{itemize}
    \item \textbf{Diplôme = Degree / Certificate / Diploma (pas \eng{diploma} seul pour tout).}
    \begin{itemize}
        \item \eng{Degree} : Diplôme universitaire (Licence, Master, Doctorat).
        \item \cle{Diploma} : Souvent technique (BTS, DUT, HND) ou scolaire (High School Diploma).
        \item \cle{Certificate} : Attestation, certificat (ex. GCE Certificate).
    \end{itemize}
    \item \textbf{Expérience $\neq$ Experiment.}
    \begin{itemize}
        \item \eng{Experience} (non comptable) = expérience professionnelle, vécu. \eng{I have 5 years' \textbf{experience}.}
        \item \eng{An experiment} (comptable) = expérience scientifique en laboratoire. \eng{A physics \textbf{experiment}.}
        \item \textbf{Incorrect :} \eng{*I made an experience in marketing.} $\rightarrow$ \textbf{Correct :} \eng{I \textbf{gained experience} in marketing.}
    \end{itemize}
    \item \textbf{Salaire : Salary vs Wages.}
    \begin{itemize}
        \item \eng{Salary} : fixe, mensuel/annuel (cadres, employés).
        \item \eng{Wages} : horaire, hebdomadaire (ouvriers, payés à la tâche). Souvent au pluriel.
    \end{itemize}
    \item \textbf{Postuler : Apply FOR.}
    \begin{itemize}
        \item \textbf{Incorrect :} \eng{*I applied the job.} $\rightarrow$ \textbf{Correct :} \eng{I \textbf{applied for} the job.}
    \end{itemize}
\end{itemize}
\end{attention}

\courssub{Exercice résolu : Emploi du lexique en contexte}

\begin{exoresolu}[Compléter le récit de recrutement]
\textbf{Consigne :} Complétez le paragraphe ci-dessous avec les mots de la liste. Attention à la forme (temps verbal, nombre, voix passive).
\par\medskip
\eng{advertise \quad apply \quad shortlist \quad interview \quad hire \quad CV \quad cover letter \quad requirements \quad salary \quad contract}
\par\medskip
\textbf{Texte :}
Last year, a local NGO (1) \_\_\_\_\_\_\_\_ a vacancy for a Field Coordinator. I (2) \_\_\_\_\_\_\_\_ immediately by sending my (3) \_\_\_\_\_\_\_\_ and a (4) \_\_\_\_\_\_\_\_. Since I met all the (5) \_\_\_\_\_\_\_\_, my application was (6) \_\_\_\_\_\_\_\_. I was then invited to (7) \_\_\_\_\_\_\_\_ before a panel of three. The interview went well and I (8) \_\_\_\_\_\_\_\_ two weeks later. I signed a two-year (9) \_\_\_\_\_\_\_\_ with a good (10) \_\_\_\_\_\_\_\_.
\tcblower
\textbf{Solution détaillée :}
\begin{enumerate}
    \item \textbf{advertised} (Passé simple, voix active : le sujet \textit{NGO} fait l'action).
    \item \textbf{applied} (Passé simple, 1ère pers. sg. : \eng{I applied for} — \textit{for} est sous-entendu ou à ajouter si la structure l'exige, ici transitif direct implicite ou \eng{applied for the post}). \textit{Note :} On attend \eng{applied for it/the post}, mais dans un texte à trous sans \textit{for}, on met souvent le verbe seul. Ici, le contexte appelle \eng{applied}.
    \item \textbf{CV} (Nom, objet de \textit{sending}).
    \item \textbf{cover letter} (Nom, second objet).
    \item \textbf{requirements} (Pluriel après \textit{all the}).
    \item \textbf{shortlisted} (Voix passive : \textit{was shortlisted} — le sujet \textit{application} subit l'action).
    \item \textbf{be interviewed} / \textbf{attend an interview} (Infinitif après \textit{invited to} — voix passive \eng{to be interviewed} ou active \eng{to attend an interview}). \textit{Les deux sont corrects.}
    \item \textbf{was hired} (Passif passé simple : \textit{I was hired}).
    \item \textbf{contract} (Nom, objet de \textit{signed}).
    \item \textbf{salary} (Nom, objet de \textit{with a good...}).
\end{enumerate}
\par\medskip
\textbf{Traduction du texte reconstitué :} « L'an dernier, une ONG locale a publié une offre pour un Coordinateur de terrain. J'ai postulé immédiatement en envoyant mon CV et une lettre de motivation. Comme je répondais à toutes les exigences, ma candidature a été présélectionnée. J'ai ensuite été invité à passer un entretien devant un jury de trois personnes. L'entretien s'est bien passé et j'ai été embauché deux semaines plus tard. J'ai signé un contrat de deux ans avec un bon salaire. »
\end{exoresolu}

\courssub{Exercice d'appropriation (À faire seul)}

\begin{exercice}[Production écrite guidée : Mon profil]
\textbf{Consigne :} Rédigez un paragraphe de 80--100 mots en anglais pour vous présenter comme candidat fictif à un poste de \eng{Marketing Assistant} dans une entreprise agro-industrielle à Douala. 
\par\medskip
\textbf{Contraintes obligatoires :}
\begin{enumerate}
    \item Utilisez au moins \textbf{10 mots} du vocabulaire de cette leçon (surlignez-les dans votre copie).
    \item Intégrez \textbf{3 collocations} de la \texttt{propriete} (p. ex. \eng{gain experience}, \eng{meet the requirements}...).
    \item Utilisez \textbf{deux dérivés} de la famille \eng{employ} ou \eng{apply} (ex. \eng{employee}, \eng{application}, \eng{unemployment}...).
    \item Évitez les faux-amis signalés dans l'\texttt{attention}.
\end{enumerate}
\par\medskip
\textbf{Amorce suggérée :} \eng{My name is ... I hold a ... in ... I am currently ...}
\end{exercice}

\begin{corrige}[Exercice 10 -- Proposition de correction]
\textit{Note : Plusieurs réponses sont possibles. Voici un modèle respectant toutes les contraintes.}
\par\medskip
\eng{My name is \textbf{Franck} and I hold a \textbf{Bachelor's degree} in Marketing from the University of Douala. After a period of \textbf{unemployment}, I \textbf{gained experience} during a six-month \textbf{internship} at a cocoa exporting company. I recently saw that your firm \textbf{advertised a vacancy} for a \textbf{Marketing Assistant}. I am writing to \textbf{apply for the post} because I \textbf{meet the requirements} perfectly: I am \textbf{computer-literate}, \textbf{fluent in English and French}, and very \textbf{dynamic}. My \textbf{application} includes my \textbf{CV} and a \textbf{cover letter}. I am ready to \textbf{attend an interview} at your convenience. I hope to \textbf{sign a contract} soon and contribute to your team as a \textbf{reliable employee}.}
\par\medskip
\textbf{Mots surlignés (13) :} Bachelor's degree, unemployment, gained experience, advertised a vacancy, Marketing Assistant, apply for the post, meet the requirements, computer-literate, dynamic, application, CV, cover letter, attend an interview, sign a contract, reliable employee.
\par\textbf{Collocations (3) :} gained experience, meet the requirements, attend an interview.
\par\textbf{Familles de mots (2) :} unemployment (employ), application (apply), employee (employ).
\end{corrige}

\coursec{Méthode d'écoute et exemple traité}

Dans la section précédente, tu as découvert le lexique du monde du recrutement : \eng{to apply for a job}, \eng{a candidate}, \eng{an interviewer}, \eng{to be hired}... Ce vocabulaire est désormais ton allié. Mais connaître des mots ne suffit pas pour réussir une épreuve de compréhension orale : il faut aussi une \cle{méthode}. Beaucoup d'élèves échouent non pas parce que l'enregistrement est trop difficile, mais parce qu'ils écoutent « au hasard », sans plan. Dans cette section, tu vas apprendre une démarche en \cle{trois étapes} — avant, pendant, après l'écoute — puis tu verras un exemple traité pas à pas, exactement comme à l'examen.

\courssub{La méthode en trois étapes : Before, While, After}

\begin{definition}[Qu'est-ce qu'une stratégie d'écoute ?]
Une \cle{stratégie d'écoute} (\eng{listening strategy}) est une action volontaire que l'auditeur accomplit pour comprendre plus efficacement. Comprendre un enregistrement, ce n'est pas subir le son : c'est \emph{prédire}, \emph{sélectionner} et \emph{vérifier}. Les bonnes stratégies se répartissent en trois moments : \eng{before listening} (avant l'écoute), \eng{while listening} (pendant l'écoute) et \eng{after listening} (après l'écoute).
\end{definition}

\begin{popfigure}
\begin{center}
\begin{tikzpicture}[
  node distance=0.9cm,
  etape/.style={rectangle, rounded corners=6pt, draw=popblue, fill=popblueL, thick, minimum width=3.6cm, minimum height=1.1cm, align=center, font=\bfseries},
  detail/.style={rectangle, rounded corners=4pt, draw=popline, fill=popcream, align=center, font=\small, text width=3.4cm},
  fleche/.style={-{Stealth[length=3mm]}, thick, popdark}
]
\node[etape, align=center] (before){BEFORE\\ Predict};
\node[etape, right=1.6cm of before, align=center] (while){WHILE\\ Note key words};
\node[etape, right=1.6cm of while, align=center] (after){AFTER\\ Write full answers};
\node[detail, below=0.7cm of before] (db) {Read the title and the questions. Predict the vocabulary. What am I looking for?};
\node[detail, below=0.7cm of while] (dw) {Note key words, figures and proper names. Use abbreviations. Never stop on an unknown word.};
\node[detail, below=0.7cm of after] (da) {Write complete sentences. Check grammar and the spelling of proper names.};
\draw[fleche] (before) -- (while);
\draw[fleche] (while) -- (after);
\draw[fleche, poporange] (after.north) to[bend left=25] node[above, font=\small\itshape, poporange]{second listening} (before.north);
\end{tikzpicture}
\legende{Organigramme de la méthode d'écoute en trois étapes}
\end{center}
\end{popfigure}

\courssub{Étape 1 — Avant l'écoute : prédire pour mieux entendre}

Observe cette situation. Deux élèves vont écouter le même entretien d'embauche. Le premier ferme les yeux et attend le son. Le second lit d'abord les questions : \eng{Where did Tabi study?}, \eng{When could he start work?}, \eng{How much experience does he have?} Lequel des deux entendra mieux ? Le second, évidemment : son cerveau sait déjà \emph{quoi chercher}.

\begin{propriete}[La règle d'or de l'anticipation]
Avant toute écoute, tu disposes toujours d'indices : le \cle{titre}, les \cle{questions}, parfois une image. Exploite-les en trois actions :
\begin{enumerate}
\item Lis le titre et déduis le thème (ici : \eng{a job interview}).
\item Lis toutes les questions et souligne les \cle{mots interrogatifs} : ils t'indiquent la \emph{nature} de l'information à capter.
\item Prédis le vocabulaire probable : pour un entretien d'embauche, attends-toi à entendre \eng{experience}, \eng{qualifications}, \eng{salary}, \eng{to start work}, \eng{references}...
\end{enumerate}
\end{propriete}

Le mot interrogatif est ta boussole. Voici le tableau de correspondance à connaître par cœur :

\begin{center}
\begin{tabularx}{\linewidth}{|l|X|X|}\hline
\rowcolor{popblueL} \textbf{Mot interrogatif} & \textbf{Ce que tu cherches} & \textbf{Exemple de réponse} \\ \hline
\eng{Who} & une personne, un nom propre & \eng{Mr Ngong, the manager.} \\ \hline
\eng{Where} & un lieu & \eng{In Buea.} \\ \hline
\eng{When} & une date, un moment & \eng{On the first of September.} \\ \hline
\eng{What time} & une heure précise & \eng{At 8 a.m.} \\ \hline
\eng{How much / How many} & une quantité, un chiffre & \eng{150,000 francs.} \\ \hline
\eng{How long} & une durée & \eng{For three years.} \\ \hline
\eng{Why} & une raison (souvent après \eng{because}) & \eng{Because he loves teaching.} \\ \hline
\end{tabularx}
\end{center}

\begin{exemplebox}[Prédire à partir d'une question]
Question de l'examen : \eng{When could Tabi start work?}

Avant même l'écoute, tu analyses : \eng{When} $\rightarrow$ je cherche une \textbf{date} ou un \textbf{moment}. Le verbe \eng{start} sera probablement prononcé dans l'enregistrement : c'est mon \cle{mot-clé} d'ancrage. Je prédis des réponses possibles : \eng{next week}, \eng{in September}, \eng{immediately}... Mon oreille est maintenant armée.
\end{exemplebox}

\courssub{Étape 2 — Pendant l'écoute : capter l'essentiel sans tout écrire}

Pendant l'écoute, le danger numéro un est de vouloir \emph{tout} noter. Impossible : pendant que tu écris une phrase, l'enregistrement continue et tu rates la suivante. La solution : ne note que les \cle{mots-clés}, les \cle{chiffres} et les \cle{noms propres}, avec des abréviations personnelles.

\begin{methode}[Noter les noms propres]
Les noms propres sont souvent \emph{épelés} lettre par lettre ou prononcés \emph{lentement} par le locuteur, car il sait qu'ils sont importants : \eng{“My name is Ngong — N-G-O-N-G.”} Dès que tu entends une épellation ou un ralentissement, note le nom \textbf{en majuscules} dans ta marge : NGONG, BUEA, HOPE FOUNDATION. Les majuscules t'évitent de confondre un nom propre avec un mot ordinaire.
\end{methode}

\begin{methode}[Noter les chiffres]
Écris les nombres en \textbf{chiffres arabes}, jamais en toutes lettres : 22, 150,000, 8 a.m. Pourquoi ? Parce que retranscrire \eng{one hundred and fifty thousand} en lettres prend du temps et multiplie les risques d'erreur. Attention aux confusions sonores classiques : \eng{thirteen} « teur-tine » / \eng{thirty} « teur-ti » ; \eng{fifteen} / \eng{fifty}. La terminaison accentuée « -tine » marque les nombres de 13 à 19.
\end{methode}

\begin{propriete}[Les abréviations de prise de notes]
Crée un code rapide et personnel. Exemples utiles pour un entretien d'embauche :
\begin{itemize}
\item \eng{exp} = experience ; \eng{qualif} = qualifications ; \eng{sal} = salary
\item \eng{mgr} = manager ; \eng{w/} = with ; \eng{b/c} = because
\item $\rightarrow$ = conduit à, entraîne ; + = et, en plus ; = = c'est-à-dire
\end{itemize}
L'important n'est pas le code lui-même, mais sa rapidité : une abréviation ne doit jamais prendre plus d'une seconde à écrire.
\end{propriete}

\begin{attention}[Ne traduis jamais mot à mot pendant l'écoute]
C'est l'erreur la plus fréquente des francophones : entendre \eng{“Well, I worked for three years in a secondary school”} et commencer à traduire mentalement « Bon, j'ai travaillé pendant trois ans dans une école secondaire... ». Pendant ce temps, la phrase suivante — qui contient la réponse à la question 3 — est déjà passée ! \textbf{Écoute les idées, pas les mots isolés.} Si un mot est inconnu, laisse-le tomber et reste accroché au sens général. Le contexte suffit presque toujours à répondre.
\end{attention}

\courssub{Les mots-repères : les panneaux indicateurs de l'enregistrement}

Un enregistrement n'est pas un flot continu : le locuteur balise son discours avec de petits mots qui annoncent une nouvelle information. Ce sont les \cle{signal words} (mots-repères). Les reconnaître, c'est savoir quand tendre l'oreille.

\begin{exemplebox}[Signal words en action]
\begin{itemize}
\item \eng{“Now, let's talk about your experience.”} — « Maintenant, parlons de votre expérience. » Le mot \eng{Now} annonce un \textbf{changement de sujet} : on passe à l'expérience. Si une question porte sur l'expérience, la réponse arrive dans les secondes qui suivent.
\item \eng{“So, you are twenty-two years old.”} — « Donc, vous avez vingt-deux ans. » \eng{So} introduit une \textbf{conclusion} ou une \textbf{vérification} : le chiffre 22 est confirmé.
\item \eng{“Well, I studied at the University of Buea.”} — « Eh bien, j'ai étudié à l'université de Buéa. » \eng{Well} marque le début d'une \textbf{réponse} du candidat.
\item \eng{“Finally, when could you start?”} — « Enfin, quand pourriez-vous commencer ? » \eng{Finally} annonce la \textbf{dernière question} de l'entretien.
\end{itemize}
\end{exemplebox}

\begin{center}
\begin{tabularx}{\linewidth}{|l|X|X|}\hline
\rowcolor{popblueL} \textbf{Signal word} & \textbf{Fonction} & \textbf{Équivalent français} \\ \hline
\eng{Now, ...} & changement de sujet & « Maintenant, ... » \\ \hline
\eng{Well, ...} & début d'une réponse, hésitation & « Eh bien, ... » \\ \hline
\eng{So, ...} & conclusion, résumé, vérification & « Donc, ... » \\ \hline
\eng{First / Then / Next} & ordre des étapes & « D'abord / Ensuite / Puis » \\ \hline
\eng{Finally, ...} & dernière information & « Enfin, ... » \\ \hline
\eng{Actually, ...} & correction ou précision & « En fait, ... » \\ \hline
\end{tabularx}
\end{center}

\begin{attention}[Faux ami : \eng{actually} ne veut pas dire « actuellement »]
\eng{Actually} signifie « en fait, en réalité » : \eng{“Actually, I have two years of experience, not three.”} — « En fait, j'ai deux ans d'expérience, pas trois. » Pour dire « actuellement », l'anglais utilise \eng{currently} ou \eng{at present} : \eng{I am currently working in Douala.} — « Je travaille actuellement à Douala. »
\end{attention}

\courssub{Étape 3 — Après l'écoute : rédiger et vérifier}

L'écoute est terminée : tes notes sont sur le brouillon. Il reste à transformer ces fragments en \cle{phrases complètes}. À l'examen, une réponse comme \eng{first of September} perd des points si la consigne exige une phrase entière.

\begin{propriete}[Rédiger la réponse]
\begin{enumerate}
\item Reprends le sujet et le verbe de la question : \eng{When could Tabi start work?} $\rightarrow$ \eng{He could start...}
\item Insère l'information notée : \eng{...on the first of September.}
\item Vérifie trois points : la \textbf{grammaire} (accord sujet-verbe, temps), les \textbf{prépositions} (\eng{on} + date, \eng{at} + heure, \eng{in} + mois) et l'\textbf{orthographe des noms propres} (majuscule : Buea, Ngong, September).
\end{enumerate}
\end{propriete}

\begin{attention}[Les prépositions du temps : piège majeur]
En français, on dit « le premier septembre », « à 8 heures », « en septembre » — une seule préposition française (« à », « en ») pour trois anglaises différentes :
\begin{itemize}
\item \eng{on} + jour ou date complète : \eng{on Monday}, \eng{on the first of September}
\item \eng{at} + heure : \eng{at 8 a.m.}, \eng{at noon}
\item \eng{in} + mois, année, saison : \eng{in September}, \eng{in 2024}
\end{itemize}
Ne dis jamais \eng{at the first of September} ni \eng{in Monday}.
\end{attention}

\courssub{Exemple traité pas à pas}

Appliquons toute la méthode sur un cas réel d'examen. Voici d'abord un extrait du script d'écoute, tel que le professeur le lit à la classe.

\begin{textebox}[Listening script — Tabi's job interview (extrait)]
Interviewer: So, Mr Tabi, you studied at the University of Buea, and you worked for three years in a secondary school in Bamenda. Now, let's talk about your availability. Finally, one last question: when could you start work?

Tabi: Well, I must give one month's notice to my current employer. So I could start on the first of September.

Interviewer: The first of September — that's perfect for us.
\end{textebox}

\begin{methode}[Résoudre une question d'écoute en 5 gestes]
Question : \eng{When could Tabi start work?}
\begin{enumerate}
\item \textbf{Avant} : \eng{When} $\rightarrow$ je cherche une date. Mot-clé d'ancrage : \eng{start}.
\item \textbf{Pendant} : j'entends le signal word \eng{Finally} $\rightarrow$ attention maximale, la question sur la date arrive. Puis j'entends \eng{“I could start on the first of September”}. Je note au brouillon : \emph{start $\rightarrow$ 1 Sept}.
\item \textbf{Après} : je rédige une phrase complète en reprenant le verbe de la question.
\item Je vérifie : modal \eng{could} conservé, préposition \eng{on} + date, majuscule à \eng{September}.
\item Réponse finale : \eng{He could start on the first of September.} — « Il pourrait commencer le premier septembre. »
\end{enumerate}
\end{methode}

\begin{exoresolu}[Application complète de la méthode]
Énoncé — Tu entends : \eng{“Well, I am twenty-two years old, and I worked for the Hope Foundation in Douala for two years. My salary was 150,000 francs a month.”} Réponds en phrases complètes : (a) \eng{How old is the candidate?} (b) \eng{Where did he work?} (c) \eng{How much did he earn?}
\tcblower
Solution détaillée —
\textbf{Avant l'écoute} : (a) \eng{How old} $\rightarrow$ un âge ; (b) \eng{Where} $\rightarrow$ un lieu ou une organisation ; (c) \eng{How much} $\rightarrow$ une somme d'argent.
\textbf{Pendant} : le signal word \eng{Well} annonce la réponse du candidat. Notes au brouillon : \emph{22 / HOPE FOUNDATION, Douala, 2 yrs / sal 150,000}.
\textbf{Après} : rédaction en phrases complètes avec vérification des majuscules et des chiffres en chiffres arabes.
\begin{itemize}
\item (a) \eng{He is twenty-two years old.} — « Il a vingt-deux ans. » (Attention : en anglais, l'âge se dit avec \eng{to be}, pas \eng{to have} : jamais \eng{He has twenty-two years}.)
\item (b) \eng{He worked for the Hope Foundation in Douala.} — « Il a travaillé pour la Fondation Hope à Douala. »
\item (c) \eng{He earned 150,000 francs a month.} — « Il gagnait 150 000 francs par mois. »
\end{itemize}
\end{exoresolu}

\courssub{Gérer le stress et exploiter la deuxième écoute}

Même avec la meilleure méthode, il arrive qu'on rate une information : un bruit dans la salle, une seconde d'inattention, un mot prononcé vite. C'est normal, et cela arrive aussi aux anglophones !

\begin{aretenir}[Les trois réflexes anti-panique]
\begin{enumerate}
\item \textbf{Ne t'arrête jamais sur ce que tu as raté.} Si tu rumines la question 2, tu rates aussi la question 3. Coupe court et reconcentre-toi immédiatement sur la suite : \emph{la phrase suivante est plus importante que celle que tu as perdue}.
\item \textbf{Exploite la deuxième écoute.} À l'examen, l'enregistrement est presque toujours passé deux fois. La première écoute sert à saisir le sens général et les réponses faciles ; la deuxième sert à \emph{combler les trous} et à vérifier. Note un point d'interrogation en face des questions non résolues pour savoir exactement quoi écouter la seconde fois.
\item \textbf{Ne laisse jamais une case vide.} Si tu n'es pas sûr, écris la réponse la plus probable d'après le contexte : une réponse plausible vaut toujours mieux qu'un blanc.
\end{enumerate}
\end{aretenir}

\begin{exercice}[À toi de jouer]
Ton professeur lit deux fois l'extrait suivant : \eng{“Now, Miss Fon, tell me about your education.” — “Well, I have a degree in management from the University of Yaoundé, and I can start work immediately.”} (a) Prédis, avant l'écoute, la nature des informations à capter. (b) Note les mots-clés avec des abréviations. (c) Rédige deux phrases complètes : sa formation, puis sa disponibilité.
\end{exercice}

\begin{corrige}[Exercice « À toi de jouer »]
(a) Le signal word \eng{Now} annonce le sujet \eng{education} $\rightarrow$ un diplôme et une université ; \eng{start work} $\rightarrow$ une date de disponibilité. (b) Notes possibles : \emph{degree mgmt, U. Yaoundé / start: immediately}. (c) \eng{She has a degree in management from the University of Yaoundé.} (« Elle a une licence en gestion de l'université de Yaoundé. ») \eng{She can start work immediately.} (« Elle peut commencer à travailler immédiatement. ») Vérification : majuscules à \eng{University of Yaoundé}, pas de préposition avant \eng{immediately}.
\end{corrige}

\begin{aretenir}[L'essentiel de la méthode]
\textbf{Before} : lis le titre et les questions, repère les mots interrogatifs, prédis le vocabulaire. \textbf{While} : note mots-clés, chiffres en chiffres arabes, noms propres en majuscules ; tends l'oreille après les signal words (\eng{Now}, \eng{Well}, \eng{So}, \eng{Finally}) ; ne traduis jamais mot à mot. \textbf{After} : rédige des phrases complètes, vérifie grammaire, prépositions et majuscules. Et si tu rates quelque chose : reste concentré, la deuxième écoute est ton amie.
\end{aretenir}

\coursec{Production guidée : réemploi du lexique}

Après avoir analysé des scripts d'entretiens d'embauche et maîtrisé le vocabulaire clé du recrutement (\emph{application, vacancy, strengths, weaknesses, experience, qualifications}), il est temps de passer de la compréhension à la production active. Cette section vous guide pour réinvestir ce lexique dans trois tâches progressives : la rédaction de votre profil personnel, l'écriture de phrases types et la simulation d'un entretien complet en binôme. Ces activités préparent directement l'épreuve orale du Baccalauréat et la rédaction d'une \eng{cover letter} (lettre de motivation).

\courssub{Tâche 1 : Ma fiche de candidat — \eng{My Candidate Profile}}

Avant de parler de soi à un recruteur, il faut structurer ses idées par écrit. Complétez la fiche ci-dessous en anglais. Utilisez le vocabulaire de la leçon (\eng{hard-working, punctual, team player, bachelor's degree, internship, part-time job}).

\begin{textebox}[Candidate Profile Form / Fiche de candidat]
\textbf{Personal Information}
\begin{itemize}
    \item \eng{Full name:} \underline{\hspace{8cm}}
    \item \eng{Age:} \underline{\hspace{2cm}} \eng{years old}
    \item \eng{City / Region:} \underline{\hspace{5cm}} (e.g., Yaoundé, Centre Region)
    \item \eng{Phone / Email:} \underline{\hspace{6cm}}
\end{itemize}
\textbf{Education \& Qualifications}
\begin{itemize}
    \item \eng{Current / Last School:} \underline{\hspace{6cm}}
    \item \eng{Highest Diploma / Expected:} \underline{\hspace{5cm}} (e.g., \eng{Baccalauréat Série A}, \eng{Bachelor's Degree in Law})
    \item \eng{Languages spoken:} \underline{\hspace{5cm}} (e.g., \eng{French (native), English (B2), Ewondo (native)})
    \item \eng{IT Skills:} \underline{\hspace{5cm}} (e.g., \eng{Microsoft Office, Internet research, Canva})
\end{itemize}
\textbf{Experience (Professional / Volunteer / School)}
\begin{itemize}
    \item \eng{Role 1:} \underline{\hspace{6cm}} \eng{Dates:} \underline{\hspace{3cm}}
    \item \eng{Main duties:} \underline{\hspace{8cm}}
    \item \eng{Role 2:} \underline{\hspace{6cm}} \eng{Dates:} \underline{\hspace{3cm}}
    \item \eng{Main duties:} \underline{\hspace{8cm}}
\end{itemize}
\textbf{Personal Strengths (\eng{Soft Skills})}
\begin{itemize}
    \item \eng{Strength 1:} \underline{\hspace{3cm}} \eng{Example:} \underline{\hspace{5cm}}
    \item \eng{Strength 2:} \underline{\hspace{3cm}} \eng{Example:} \underline{\hspace{5cm}}
    \item \eng{Strength 3:} \underline{\hspace{3cm}} \eng{Example:} \underline{\hspace{5cm}}
\end{itemize}
\textbf{Career Objective}
\begin{itemize}
    \item \eng{My dream job / Target position:} \underline{\hspace{6cm}}
    \item \eng{Why this job?} \underline{\hspace{8cm}}
\end{itemize}
\end{textebox}

\begin{exemplebox}[Exemple de fiche remplie (Élève de Tle A, Yaoundé)]
\textbf{Name:} \eng{Ngo Bitje Marie-Claire} \\
\textbf{Age:} \eng{19 years old} \\
\textbf{Diploma:} \eng{Baccalauréat Série A (expected June 2025)} \\
\textbf{Languages:} \eng{French (C2), English (B1), Bassa (Native)} \\
\textbf{Experience:} \eng{Volunteer at “Jeunesse Active” NGO (2023–2024) — Organised reading workshops for primary kids.} \\
\textbf{Strengths:} \eng{Organised (planned 10 events), Team player (led a group of 5), Adaptable.} \\
\textbf{Dream Job:} \eng{Project Assistant in an International NGO.} \eng{Because I want to help communities in the Far North region.}
\end{exemplebox}

\courssub{Tâche 2 : Cinq phrases personnelles — \eng{Five Personal Sentences}}

Rédigez cinq phrases complètes en anglais en réemployant \textbf{au moins 10 mots ou expressions} du lexique de la leçon (voir tableau récapitulatif page suivante). Surlignez ou mettez en gras le vocabulaire ciblé.

\begin{methode}[Comment rédiger ses phrases types]
\begin{enumerate}
    \item \textbf{Ciblez le vocabulaire :} Choisissez 2 à 3 mots par phrase dans la liste : \eng{apply for, vacancy, CV, cover letter, interview, experience, qualifications, strengths, weaknesses, hard-working, punctual, reliable, team player, salary, contract, full-time, part-time}.
    \item \textbf{Structurez :} Suivez l'ordre \eng{Subject + Verb + Complement}. Utilisez les temps adaptés : \eng{Present Simple} (habitudes, vérités), \eng{Present Perfect} (expérience vécue), \eng{Would like to} (souhait poli).
    \item \textbf{Personnalisez :} Ancrez la phrase dans VOTRE réalité (votre ville, votre école, vos stages).
    \item \textbf{Vérifiez :} Accord sujet-verbe (\eng{He has}, pas \eng{He have}), prépositions (\eng{apply FOR}, \eng{good AT}, \eng{interested IN}).
\end{enumerate}
\end{methode}

\begin{exemplebox}[Modèles de phrases (vocabulaire en \textbf{gras})]
\begin{enumerate}
    \item \eng{I would like to \textbf{apply for} the \textbf{vacancy} of \textbf{shop assistant} at the new supermarket in \textbf{Buea}.} (Je voudrais postuler au poste de vendeur au nouveau supermarché de Buéa.)
    \item \eng{My main \textbf{strengths} are being \textbf{hard-working} and \textbf{punctual}; I have never been late to school this year.} (Mes principaux atouts sont d'être travailleur et ponctuel ; je n'ai jamais été en retard à l'école cette année.)
    \item \eng{During my last \textbf{internship} at a law firm in \textbf{Douala}, I gained valuable \textbf{experience} in client reception.} (Lors de mon dernier stage dans un cabinet d'avocats à Douala, j'ai acquis une expérience précieuse en accueil client.)
    \item \eng{I have attached my \textbf{CV} and \textbf{cover letter} to this email as requested in the job \textbf{advertisement}.} (J'ai joint mon CV et ma lettre de motivation à cet email comme demandé dans l'annonce.)
    \item \eng{Although I have no professional \textbf{qualifications} yet, I am a fast \textbf{learner} and a reliable \textbf{team player}.} (Bien que je n'aie pas encore de qualifications professionnelles, j'apprends vite et je suis un équipier fiable.)
\end{enumerate}
\end{exemplebox}

\begin{exercice}[Vos 5 phrases]
Écrivez vos 5 phrases ci-dessous. Faites-les corriger par votre professeur ou un camarade.
\begin{enumerate}
    \item \underline{\hspace{14cm}}
    \item \underline{\hspace{14cm}}
    \item \underline{\hspace{14cm}}
    \item \underline{\hspace{14cm}}
    \item \underline{\hspace{14cm}}
\end{enumerate}
\end{exercice}

\courssub{Tâche 3 : Mini-dialogue guidé en binôme — \eng{Role-play: Job Interview}}

C'est l'activité centrale. En binôme, l'un joue le \textbf{Recruiter (A)}, l'autre le \textbf{Candidate (B)}. Suivez \textbf{strictement} la trame en 5 étapes ci-dessous. Préparez-vous 5 minutes, jouez 3 minutes, puis inversez les rôles.

\begin{popfigure}
\begin{tikzpicture}[
    node distance=1.2cm and 0.5cm,
    every node/.style={font=\small, align=center},
    stepS/.style={rectangle, rounded corners, draw=popblue, thick, fill=popblueL, minimum width=2.8cm, minimum height=1.4cm},
    arrow/.style={-Stealth, thick, popblue}
]
% Nœuds
\node[stepS, align=center] (greet){\textbf{1. Greeting}\\\eng{Good morning.}\\\eng{Please sit down.}};
\node[stepS, right=of greet, align=center] (self){\textbf{2. About Yourself}\\\eng{Tell me about}\\\eng{yourself.}};
\node[stepS, right=of self, align=center] (exp){\textbf{3. Experience}\\\eng{Do you have any}\\\eng{experience?}};
\node[stepS, right=of exp, align=center] (str){\textbf{4. Strengths}\\\eng{What are your}\\\eng{strengths?}};
\node[stepS, right=of str, align=center] (close){\textbf{5. Closing}\\\eng{Thank you. We will}\\\eng{call you.}};

% Flèches
\draw[arrow] (greet.east) -- (self.west);
\draw[arrow] (self.east) -- (exp.west);
\draw[arrow] (exp.east) -- (str.west);
\draw[arrow] (str.east) -- (close.west);

% Légendes rôle A / B
\node[above=0.2cm of greet, text=popdark, font=\bfseries\small] (la) {Recruiter (A)};
\node[below=0.2cm of greet, text=poporange, font=\bfseries\small] (lb) {Candidate (B)};
\node[above=0.2cm of self, text=popdark, font=\bfseries\small] {Recruiter (A)};
\node[below=0.2cm of self, text=poporange, font=\bfseries\small] {Candidate (B)};
\node[above=0.2cm of exp, text=popdark, font=\bfseries\small] {Recruiter (A)};
\node[below=0.2cm of exp, text=poporange, font=\bfseries\small] {Candidate (B)};
\node[above=0.2cm of str, text=popdark, font=\bfseries\small] {Recruiter (A)};
\node[below=0.2cm of str, text=poporange, font=\bfseries\small] {Candidate (B)};
\node[above=0.2cm of close, text=popdark, font=\bfseries\small] {Recruiter (A)};
\node[below=0.2cm of close, text=poporange, font=\bfseries\small] {Candidate (B)};
\end{tikzpicture}
\legende{Schéma du dialogue d'entretien en 5 étapes (Flux Recruiter $\rightarrow$ Candidate)}
\end{popfigure}

\begin{textebox}[Trame de dialogue imposée — \eng{Script Frame}]
\textbf{Step 1: Greeting}
\begin{itemize}
    \item \textbf{A (Recruiter):} \eng{Good morning. Please sit down.}
    \item \textbf{B (Candidate):} \eng{Good morning, Sir / Madam. Thank you.}
\end{itemize}
\textbf{Step 2: Tell me about yourself}
\begin{itemize}
    \item \textbf{A:} \eng{Tell me about yourself.}
    \item \textbf{B:} \eng{My name is [Name]. I am [Age] years old. I come from [City/Region]. I am a student in [Series/Field] at [School Name].}
\end{itemize}
\textbf{Step 3: Experience}
\begin{itemize}
    \item \textbf{A:} \eng{Do you have any experience? / Have you ever worked before?}
    \item \textbf{B:} \eng{Yes, I have worked as a [Job/Role] at [Place] for [Duration]. / No, I haven't had a job yet, but I did an internship at [Place] / I have volunteered at [Organisation].}
\end{itemize}
\textbf{Step 4: Strengths}
\begin{itemize}
    \item \textbf{A:} \eng{What are your strengths? / Why should we hire you?}
    \item \textbf{B:} \eng{I am [Adjective 1] and [Adjective 2]. For example, [Concrete example from school/life].}
\end{itemize}
\textbf{Step 5: Closing}
\begin{itemize}
    \item \textbf{A:} \eng{Thank you for coming. We will call you next week.}
    \item \textbf{B:} \eng{Thank you very much for your time. Goodbye, Sir / Madam.}
\end{itemize}
\end{textebox}

\begin{methode}[Parler de soi en 3 points — \eng{The "Tell me about yourself" Structure}]
Pour réussir l'étape 2 (\eng{Tell me about yourself}), structurez votre réponse en \textbf{trois blocs logiques}. Ne récitez pas votre CV, racontez une mini-histoire cohérente.

\begin{center}
\begin{tabularx}{\linewidth}{|X|X|X|}
\hline
\rowcolor{popblueL}
\textbf{Bloc 1 : Identity} & \textbf{Bloc 2 : Education} & \textbf{Bloc 3 : Qualities} \\
(\eng{Who are you?}) & (\eng{What did you learn?}) & (\eng{What is your value?}) \\
\hline
\eng{Name, Age, Origin} & \eng{Current school, Series, Diploma} & \eng{2 Adjectives + 1 Proof} \\
\eng{\textit{“I am Paul, 19, from Bamenda.”}} & \eng{\textit{“I study Literature at Lycée Bilingue.”}} & \eng{\textit{“I am organised and curious. I led the debate club.”}} \\
\hline
\end{tabularx}
\end{center}
\textbf{Durée idéale :} 45 à 60 secondes. Regardez le recruteur dans les yeux, souriez, parlez distinctement.
\end{methode}

\begin{attention}[Pièges classiques des francophones — \eng{Common Mistakes}]
\begin{itemize}
    \item \textbf{Âge :} \eng{WRONG: “I have 19 years.”} $\rightarrow$ \eng{RIGHT: “I am 19 years old.”} ou \eng{“I am 19.”} (Verbe \eng{BE}, pas \eng{HAVE}).
    \item \textbf{Postuler :} \eng{WRONG: “I apply the job.”} $\rightarrow$ \eng{RIGHT: “I apply \textbf{for} the job.”} (Préposition \textbf{for} obligatoire).
    \item \textbf{Expérience :} \eng{WRONG: “I have experience in this domain since 2 years.”} $\rightarrow$ \eng{RIGHT: “I have two years’ experience.” / “I have worked in this field for two years.”} (Present Perfect + \eng{for} ou structure nominale).
    \item \textbf{Politesse :} \eng{WRONG: “What you want?”} $\rightarrow$ \eng{RIGHT: “What \textbf{would you like}?”} ou \eng{“What \textbf{are you looking for}?”} (Formes polies).
    \item \textbf{Points forts :} \eng{WRONG: “My qualities are...”} $\rightarrow$ \eng{BETTER: “My \textbf{strengths} are...”} (Vocabulaire pro : \eng{strengths} vs \eng{weaknesses}).
\end{itemize}
\end{attention}

\courssub{Tableau récapitulatif : Lexique clé de l'entretien}

\begin{center}
\begin{tabularx}{\linewidth}{|X|X|X|}
\hline
\rowcolor{popblueL}
\textbf{English} & \textbf{Français} & \textbf{Contexte / Collocation} \\
\hline
\eng{To apply for} & Postuler à & \eng{apply for a job / a vacancy / a post} \\
\eng{Vacancy / Opening} & Poste vacant & \eng{There is a vacancy for a secretary.} \\
\eng{CV / Resume (US)} & Curriculum Vitae & \eng{attach my CV / update my CV} \\
\eng{Cover letter} & Lettre de motivation & \eng{write a cover letter / send a cover letter} \\
\eng{Interview} & Entretien & \eng{attend an interview / succeed in an interview (réussir)} \\
\eng{Experience} & Expérience & \eng{work experience / previous experience / gain experience} \\
\eng{Qualifications} & Qualifications / Diplômes & \eng{academic / professional qualifications} \\
\eng{Strengths} & Points forts / Atouts & \eng{My strengths are... / key strengths} \\
\eng{Weaknesses} & Points faibles & \eng{What are your weaknesses? / improve my weaknesses} \\
\eng{Hard-working} & Travailleur / Acharné & \eng{a hard-working student / employee} \\
\eng{Punctual / Reliable} & Ponctuel / Fiable & \eng{be punctual for / a reliable person} \\
\eng{Team player} & Esprit d'équipe & \eng{work as a team player / good team player} \\
\eng{Internship} & Stage & \eng{do an internship / internship report} \\
\eng{Salary / Wages} & Salaire & \eng{monthly salary / negotiate salary} \\
\eng{Contract} & Contrat & \eng{sign a contract / full-time / part-time contract} \\
\hline
\end{tabularx}
\end{center}

\courssub{Restitution devant la classe \& Grille d'observation}

Deux ou trois binômes volontaires (ou tirés au sort) présentent leur dialogue devant la classe. Les autres élèves observent en remplissant la grille ci-dessous. Cela développe l'écoute active (\eng{active listening}) et l'esprit critique.

\begin{experience}[Grille d'observation par les pairs — \eng{Peer Observation Grid}]
\textbf{Consigne :} Pendant la restitution, cochez \eng{Yes}, \eng{Partly} ou \eng{No} et notez un exemple entendu.
\begin{center}
\begin{tabularx}{\linewidth}{|X|c|c|c|X|}
\hline
\rowcolor{popblueL}
\textbf{Critère d'observation} & \textbf{Yes} & \textbf{Partly} & \textbf{No} & \textbf{Exemple entendu / Commentaire} \\
\hline
\textbf{Lexique} : Utilise au moins 5 mots du tableau (vacancy, strengths, experience...) & & & & \\
\hline
\textbf{Politesse} : Formules de politesse correctes (\eng{Sir/Madam, Please, Thank you, Would like}) & & & & \\
\hline
\textbf{Fluidité} : Parle sans trop d'hésitations, débit naturel, liaison correcte & & & & \\
\hline
\textbf{Prononciation} : Mots clés bien prononcés (\eng{interview, experience, qualifications, strengths}) & & & & \\
\hline
\textbf{Structure} : Respecte les 5 étapes (Greeting $\rightarrow$ Closing) & & & & \\
\hline
\textbf{Contact visuel} : Regarde le partenaire / la classe, ne lit pas sur sa feuille & & & & \\
\hline
\end{tabularx}
\end{center}
\textbf{Après chaque passage :} 1 minute de feedback oral constructif (\eng{“Good job on... Next time, try to...”}).
\end{experience}

\courssub{Complément Bac : De l'oral à l'écrit}

\begin{savaistu}[L'entretien oral et la \eng{Cover Letter} : Deux faces d'une même pièce]
L'épreuve orale du Baccalauréat (série A) inclut souvent une \textbf{simulation d'entretien} ou une \textbf{présentation personnelle}. La trame à 5 étapes que vous venez de travailler est la structure standard attendue.

Parallèlement, l'épreuve écrite peut demander une \textbf{lettre de motivation (\eng{cover letter})}. Les paragraphes de la lettre reprennent \textbf{exactement} la logique du dialogue oral :
\begin{center}
\begin{tabular}{|l|l|}
\hline
\rowcolor{popblueL}
\textbf{Étape Orale (Dialogue)} & \textbf{Paragraphe Écrit (Lettre)} \\
\hline
1. Greeting & En-tête + Objet : \eng{Application for the post of...} \\
2. About Yourself & \eng{Introduction :} Qui je suis, quel poste je vise. \\
3. Experience & \eng{Corps 1 :} Mon parcours, stages, bénévolat, compétences techniques. \\
4. Strengths & \eng{Corps 2 :} Mes qualités humaines (\eng{soft skills}) + preuves concrètes. \\
5. Closing & \eng{Conclusion :} Disponibilité pour entretien, formule de politesse formelle. \\
\hline
\end{tabular}
\end{center}
\textbf{Astuce Bac :} Rédigez votre \eng{cover letter} \textbf{après} avoir fait le role-play. Les idées sont déjà triées, il ne reste qu'à les mettre en forme écrite (style formel, pas de contractions \eng{I'm $\rightarrow$ I am}, paragraphes aérés).
\end{savaistu}

\begin{exoresolu}[Simulation complète : Réponse type pour un élève de Tle A]
\textbf{Contexte :} Poste de \eng{Library Assistant} à la Bibliothèque Municipale de Yaoundé. Candidats : \eng{Chantal (B)}.
\tcblower
\textbf{A:} \eng{Good morning. Please sit down.} \\
\textbf{B:} \eng{Good morning, Madam. Thank you.} \\
\textbf{A:} \eng{Tell me about yourself.} \\
\textbf{B:} \eng{My name is Chantal Atangana. I am 18 years old. I come from Mfoundi, Yaoundé. I am a Terminale A student at Lycée Leclerc. I study Literature and Languages.} \\
\textbf{A:} \eng{Do you have any experience?} \\
\textbf{B:} \eng{Yes, I have worked as a volunteer at my school library for two years. I helped students find books and I classified new arrivals. I also did a one-month internship at the National Library last summer.} \\
\textbf{A:} \eng{What are your strengths?} \\
\textbf{B:} \eng{I am very organised and punctual. For example, I created a new filing system for the school novels last term. I am also a good team player; I work well with the librarian teachers.} \\
\textbf{A:} \eng{Thank you for coming. We will call you next week.} \\
\textbf{B:} \eng{Thank you very much for your time. Goodbye, Madam.} \\
\\
\textbf{Analyse :} Respect de la trame. Vocabulaire ciblé utilisé (\eng{volunteer, classified, internship, organised, punctual, team player, filing system}). Grammaire correcte (\eng{I am 18}, \eng{I have worked}, \eng{I created}). Politesse parfaite.
\end{exoresolu}

\begin{exercice}[Entraînement individuel : Écriture de la \eng{Cover Letter}]
À partir de votre fiche de candidat (Tâche 1) et de votre dialogue (Tâche 3), rédigez une lettre de motivation formelle (\eng{cover letter}) de 120-150 mots pour le poste de vos rêves. Respectez la mise en page standard (expéditeur, destinataire, date, objet, corps en 3 paragraphes, formule de politesse finale, signature).
\begin{itemize}
    \item \textbf{Expéditeur :} Vos coordonnées (haut gauche).
    \item \textbf{Destinataire :} Nom du responsable / Entreprise (haut droite).
    \item \textbf{Objet :} \eng{Application for the position of [Job Title] - Ref: [Numéro si dispo]}.
    \item \textbf{Corps :} \eng{Dear Sir/Madam,} ... \eng{Yours faithfully,} [Votre Nom].
\end{itemize}
\end{exercice}

\begin{corrige}[Exercice : Cover Letter]
\textbf{Modèle type (adaptez les parties entre crochets) :}
\begin{textebox}[Cover Letter Model]
\eng{Chantal Atangana} \\
\eng{Quarter Mvog-Ada, Yaoundé} \\
\eng{+237 6XX XX XX XX} \\
\eng{chantal.atangana@email.com} \\
\vspace{0.5cm}
\eng{The Director} \\
\eng{Yaoundé Municipal Library} \\
\eng{P.O. Box 123, Yaoundé} \\
\vspace{0.5cm}
\eng{Yaoundé, 15th October 2025} \\
\vspace{0.5cm}
\eng{\textbf{Subject: Application for the position of Library Assistant}} \\
\vspace{0.5cm}
\eng{Dear Sir/Madam,} \\
\vspace{0.2cm}
\eng{I am writing to apply for the vacancy for Library Assistant advertised on your website. I am an 18-year-old Terminale A student at Lycée Leclerc, passionate about literature and information management.} \\
\vspace{0.2cm}
\eng{During the past two years, I have gained valuable experience as a volunteer at my school library. I assisted students with research, managed the book lending system, and classified new acquisitions. Last summer, I completed a one-month internship at the National Library where I learnt digital archiving techniques.} \\
\vspace{0.2cm}
\eng{My main strengths are organisation, punctuality and a strong sense of service. For instance, I designed a new colour-coded filing system for novels which reduced search time by 20\%. I am a reliable team player and I adapt quickly to new procedures.} \\
\vspace{0.2cm}
\eng{I am available for an interview at your earliest convenience. Thank you for considering my application.} \\
\vspace{0.5cm}
\eng{Yours faithfully,} \\
\vspace{1cm}
\eng{Chantal Atangana}
\end{textebox}
\end{corrige}

\begin{aretenir}[Check-list finale — Suis-je prêt pour l'oral Bac ?]
\begin{itemize}
    \item [ ] Je sais remplir une fiche de candidature en anglais sans erreur (\eng{I am 19}, pas \eng{I have 19}).
    \item [ ] Je connais par cœur 10 mots-clés du recrutement et leurs collocations (\eng{apply for, gain experience, key strengths}).
    \item [ ] Je peux me présenter en 50 secondes : Identité $\rightarrow$ Formation $\rightarrow$ 2 Qualités + Preuve.
    \item [ ] Je maîtrise la trame en 5 étapes du dialogue (Greeting $\rightarrow$ Closing).
    \item [ ] J'utilise la politesse formelle (\eng{Sir/Madam, Would like, Thank you for your time}).
    \item [ ] Je peux transformer mon oral en lettre de motivation structurée (3 paragraphes).
\end{itemize}
\end{aretenir}

\newpage

\coursec{Activité d'intégration}

\courssub{Objectifs de l'activité}

À l'issue de cette activité d'intégration, l'élève sera capable de :
\begin{itemize}
\item réinvestir les \cle{stratégies d'écoute} vues dans la leçon (repérer le thème, les mots-clés, les chiffres et les noms propres) dans une situation de communication réelle ;
\item utiliser correctement le \cle{lexique du recrutement} (\eng{applicant, employer, salary, experience, qualifications, punctual, to hire}) dans un échange oral ;
\item poser et comprendre les questions types d'un \cle{entretien d'embauche} (\eng{job interview}) ;
\item relever des informations précises (nom, âge, expérience, salaire demandé) en écoutant un camarade, comme lors d'une épreuve de \eng{listening} ;
\item formuler une \cle{recommandation} orale justifiée en anglais correct.
\end{itemize}

\courssub{Situation}

Le conseil des clubs d'anglais de ton lycée organise un \eng{Job Interview Day} : la salle de classe est transformée en cabinet de recrutement pour une matinée. Plusieurs petites entreprises de Yaoundé cherchent des jeunes diplômés pour des emplois à temps partiel pendant les grandes vacances. Une offre d'emploi est affichée au tableau. Tu vas, à tour de rôle, jouer le recruteur, le candidat et l'observateur qui prend des notes.

\begin{textebox}[Job advertisement — affichée au tableau]
WANTED!

Part-time secretary for “Connect House”, a busy cybercafé in Mvog-Ada, Yaoundé.

We are looking for a young, serious and friendly person to:
\begin{itemize}
\item welcome customers and answer the phone;
\item type documents and send e-mails;
\item keep the cash book every evening.
\end{itemize}

Requirements: at least 18 years old, good English and French, basic computer skills. Experience is an advantage but not compulsory.

Working hours: Monday to Friday, 4 p.m. -- 8 p.m.

Salary: 60,000 FCFA per month.

Interested? Apply now! Interviews will take place on Saturday 15th June at 9 a.m. Bring your CV and a recommendation letter.

Contact: Mr Etoundi, manager.
\end{textebox}

\begin{definition}[Glossaire de l'offre d'emploi]
\begin{itemize}
\item \eng{part-time} (adj.) : à temps partiel ; prononciation « parte-taïme ».
\item \eng{requirements} (n. pluriel) : les conditions requises, les exigences.
\item \eng{basic computer skills} : des notions de base en informatique.
\item \eng{an advantage} : un atout ; \eng{not compulsory} : pas obligatoire.
\item \eng{to apply} (v.) : poser sa candidature ; prononciation « a-plaï ».
\item \eng{a recommendation letter} : une lettre de recommandation.
\end{itemize}
\end{definition}

\courssub{Consignes}

Travail par groupes de trois élèves : un \cle{recruteur} (\eng{interviewer}), un \cle{candidat} (\eng{applicant}) et un \cle{observateur} (\eng{listener / note-taker}). Les rôles tournent à chaque tour, de sorte que chacun joue les trois rôles.

\begin{enumerate}
\item \textbf{Préparation (10 minutes).} Lis attentivement l'offre d'emploi au tableau. Souligne les informations essentielles : le poste, le lieu, les horaires, le salaire, les conditions requises. C'est la même démarche que l'anticipation avant une écoute : repérer le thème et les mots-clés.
\item \textbf{Le recruteur prépare ses questions.} Il écrit au moins \cle{cinq questions} d'entretien en s'aidant des modèles vus en leçon : \eng{What is your name? How old are you? Do you have any experience? Why do you want this job? What salary do you expect? When can you start?}
\item \textbf{Le candidat prépare sa fiche personnelle.} Il invente ou complète son identité : nom, âge, niveau d'études, expérience, deux qualités (\eng{punctual, hard-working, honest, friendly}) et le salaire qu'il demande.
\item \textbf{L'observateur prépare sa grille d'écoute.} Il recopie la fiche du candidat avec les rubriques : \eng{Name / Age / Experience / Qualities / Salary expected}. Pendant l'entretien, il écoute \textbf{sans interrompre} et remplit la fiche, exactement comme pendant un enregistrement de \eng{listening} : il note les chiffres, les noms propres et les mots-clés.
\item \textbf{Simulation de l'entretien (5 minutes par tour).} Le recruteur accueille le candidat poliment (\eng{Good morning, please sit down.}), pose ses questions, puis conclut (\eng{Thank you very much. We will call you soon.}). Le candidat répond en phrases complètes.
\item \textbf{Vérification et recommandation.} L'observateur lit sa fiche au candidat pour vérifier les informations (\eng{Is that correct?}). Puis il annonce oralement sa décision au groupe : \eng{I recommend Tabi because he is punctual and experienced.} ou \eng{I do not recommend her because she has no computer skills.} La justification est obligatoire.
\item \textbf{Rotation des rôles.} On change de rôle et on recommence avec un nouveau candidat, jusqu'à ce que chacun ait joué les trois rôles.
\item \textbf{Évaluation.} L'enseignant observe et note selon la grille : lexique du recrutement (5 pts), exactitude des informations relevées par l'observateur (5 pts), politesse et formules de l'entretien (5 pts), prononciation et fluidité (5 pts).
\end{enumerate}

\begin{attention}[Pièges à éviter pendant la simulation]
\begin{itemize}
\item Ne dis pas \eng{I have 19 years} : en anglais on utilise \eng{to be} : \eng{I am 19 years old.} (« J'ai 19 ans. »)
\item \eng{Salary} se prononce « sa-la-ri » ; ne confonds pas avec \eng{celery} (« céleri »).
\item Pour l'expérience, utilise le present perfect ou le past simple : \eng{I worked in a shop last year.} et non « I have worked last year » avec une date passée précise.
\item Reste poli : on dit \eng{Could you repeat, please?} et non « What? » si l'on n'a pas compris.
\end{itemize}
\end{attention}

\courssub{Ressources à mobiliser}

\begin{aretenir}[Boîte à outils de l'entretien d'embauche]
\textbf{Questions du recruteur :} \eng{What is your name? / How old are you? / Where do you live? / Do you have any experience? / What are your qualities? / Why do you want this job? / What salary do you expect? / When can you start?}

\textbf{Réponses du candidat :} \eng{My name is... / I am eighteen years old. / I live in Nkolbisson, Yaoundé. / Yes, I worked in my uncle's shop during the holidays. / I am punctual and hard-working. / I want this job because I need money for my studies. / I expect 60,000 FCFA a month. / I can start immediately.}

\textbf{Formules de politesse :} \eng{Good morning, sir/madam. / Please sit down. / Thank you for coming. / Could you repeat, please? / It was a pleasure to meet you.}

\textbf{Recommandation de l'observateur :} \eng{I recommend + nom + because + sujet + verbe...} ; \eng{I do not recommend + nom + because...}

\textbf{Stratégies d'écoute :} anticiper le thème, noter les chiffres et les noms propres, repérer les mots-clés, vérifier ses notes à la seconde écoute.
\end{aretenir}

\courssub{Proposition de production et corrigé commenté}

\begin{exoresolu}[Un entretien complet et une recommandation — modèle]
Voici la simulation réalisée par un groupe de Tle A. Relis-la et observe les choix de langue.

\textbf{Interviewer:} Good morning, please sit down. What is your name?

\textbf{Applicant:} Good morning, sir. My name is Tabi Ngong.

\textbf{Interviewer:} How old are you, Tabi?

\textbf{Applicant:} I am nineteen years old. I am a student in Terminale at Lycée de Nkolbisson.

\textbf{Interviewer:} Do you have any experience for this job?

\textbf{Applicant:} Yes, I do. I worked in my uncle's cybercafé in Buea last year, during the long holidays. I welcomed customers and typed documents.

\textbf{Interviewer:} Very good. What are your qualities?

\textbf{Applicant:} I am punctual, honest and friendly. I also speak English and French well.

\textbf{Interviewer:} Why do you want this job?

\textbf{Applicant:} I want this job because I need money for my university studies, and I like working with computers.

\textbf{Interviewer:} What salary do you expect?

\textbf{Applicant:} I expect 60,000 FCFA a month, as in your advertisement.

\textbf{Interviewer:} When can you start?

\textbf{Applicant:} I can start immediately, sir.

\textbf{Interviewer:} Thank you very much. We will call you soon.

\textbf{Applicant:} Thank you, sir. Goodbye.

\textbf{Observer's notes:} Name: Tabi Ngong. Age: 19. Experience: worked in a cybercafé in Buea last year. Qualities: punctual, honest, friendly. Salary expected: 60,000 FCFA.

\textbf{Observer's recommendation:} I recommend Tabi because he is punctual and experienced. He worked in a cybercafé before, and he accepts the salary of 60,000 FCFA.

\tcblower

\textbf{Commentaire des choix de langue :}
\begin{enumerate}
\item \textbf{Ouverture polie.} \eng{Good morning, please sit down.} installe un ton formel adapté à un entretien ; le candidat répond avec \eng{sir}, marque de respect attendue dans ce contexte.
\item \textbf{Âge avec \eng{to be}.} \eng{I am nineteen years old} respecte la règle anglaise (pas de « have » pour l'âge), erreur classique des francophones évitée ici.
\item \textbf{Past simple daté.} \eng{I worked in my uncle's cybercafé last year} : le past simple est correct car le moment (\eng{last year}) est précis et terminé. Le present perfect serait fautif avec cette date.
\item \textbf{Lexique de la leçon réemployé.} \eng{experience, qualities, punctual, salary, to expect, advertisement} : au moins six mots du champ lexical du recrutement sont mobilisés, ce qui remplit le critère « lexique » de la grille.
\item \textbf{Réponses en phrases complètes.} Le candidat ne répond jamais par un seul mot : \eng{I expect 60,000 FCFA a month} et non « 60,000 ». Cela entraîne la fluidité et la grammaire.
\item \textbf{Recommandation structurée.} L'observateur applique le schéma \eng{I recommend + nom + because + justification}, puis ajoute deux arguments précis tirés de ses notes (\eng{he worked in a cybercafé before} ; \eng{he accepts the salary}), preuve d'une écoute exacte des informations.
\end{enumerate}
\end{exoresolu}

\courssub{Bilan de l'activité}

Cette simulation a permis de transformer les compétences de la leçon en pratique réelle : tu as \cle{anticipé} un document (l'offre d'emploi), \cle{écouté activement} un partenaire pour relever des informations précises, \cle{réemployé le lexique} du recrutement dans un dialogue authentique et \cle{justifié une décision} à l'oral. La grille d'évaluation montre que la réussite d'un entretien repose sur quatre piliers : le vocabulaire, l'exactitude des informations, la politesse et la prononciation. Ces mêmes stratégies — repérer le thème, noter les chiffres et les noms propres, vérifier ses notes — te serviront à l'épreuve de \eng{listening} du baccalauréat comme dans un véritable entretien d'embauche. Pour aller plus loin, rédige à la maison ta propre fiche de candidat et une courte lettre de candidature (\eng{application letter}) en réponse à l'offre de \eng{Connect House} : ce sera la transition naturelle vers la leçon d'expression écrite.

\newpage

\coursec{Méthodes et exercices résolus}

\coursec{Lesson 1 — Listening: Texts about undergoing a job interview}

\courssub{Mise en route et anticipation}

\begin{experience}[Brainstorming : le vocabulaire de l'entretien]
En classe entière, le professeur écrit au tableau \eng{Job Interview}. Les élèves proposent en anglais tous les mots qui leur viennent à l'esprit (ex. \eng{application, CV, qualifications, salary, boss, nervous, handshake, dress code}). Le professeur regroupe les mots en catégories : \textbf{Documents} (\eng{CV, cover letter, application form}), \textbf{Acteurs} (\eng{interviewer, candidate, panel, HR manager}), \textbf{Déroulement} (\eng{greeting, questions, answers, salary negotiation, closing}), \textbf{Qualités} (\eng{punctuality, confidence, honesty, skills}). Cette carte mentale reste affichée pendant la leçon.
\end{experience}

\begin{methode}[Avant l'écoute : anticiper pour mieux comprendre]
\begin{enumerate}
\item \textbf{Lisez les questions} (ou le titre de l'exercice) avant le début de l'enregistrement. Soulignez les \cle{mots-clés} : noms propres, chiffres, dates, noms de poste, verbes d'action.
\item \textbf{Prédisez le contenu} : un entretien d'embauche suit presque toujours un schéma prévisible : \eng{greetings} $\rightarrow$ \eng{qualifications} $\rightarrow$ \eng{experience} $\rightarrow$ \eng{strengths/weaknesses} $\rightarrow$ \eng{salary expectations} $\rightarrow$ \eng{availability} $\rightarrow$ \eng{closing}.
\item \textbf{Visualisez la scène} : qui parle ? (un \eng{recruiter} / un \eng{candidate}), où ? (\eng{an office, a meeting room, online}), dans quel but ? (\eng{to fill a vacancy}).
\end{enumerate}
\end{methode}

\courssub{Le texte support : script d'écoute}

\begin{textebox}[Script d'écoute : Entretien chez Douala Trading Company]
\eng{Good morning, Miss Ekane. Please have a seat.}
\eng{Good morning, Sir. Thank you.}
\eng{So, you are applying for the post of secretary in our company. Can you tell me about your qualifications?}
\eng{Yes, Sir. I obtained my GCE Advanced Level certificate two years ago, and I hold a diploma in Computer Studies from a training centre in Douala.}
\eng{Have you got any professional experience?}
\eng{I have worked as a receptionist at the Hilton Hotel for eighteen months.}
\eng{Very good. Can you speak French and English?}
\eng{Yes, I speak both fluently.}
\eng{Fine. What are your salary expectations?}
\eng{I am open to discussion, but I know the market rate for a junior secretary.}
\eng{When can you start?}
\eng{I can start next Monday, Sir.}
\eng{Thank you for coming, Miss Ekane. We shall contact you soon.}
\eng{Thank you, Sir. Goodbye.}
\end{textebox}

\begin{savaistu}[Contexte camerounais : le marché de l'emploi]
Au Cameroun, les entretiens se déroulent souvent en \eng{English} dans les entreprises anglophones ou internationales (ex. \eng{CDC, SONARA, banks, NGOs}) et en français ailleurs. Le \eng{GCE Advanced Level} (niveau A4) est le diplôme de fin d'études secondaires anglophones, équivalent au Baccalauréat. Les postes de \eng{secretary}, \eng{receptionist}, \eng{driver}, \eng{sales representative} sont très demandés à Yaoundé, Douala, Buea, Bamenda. La \eng{punctuality} (ponctualité) et le \eng{dress code} (costume ou tailleur, tenue soignée) sont des critères éliminatoires. Le \eng{handshake} (poignée de main) ferme et le contact visuel (\eng{eye contact}) marquent le respect et la confiance.
\end{savaistu}

\courssub{Compréhension : questions variées}

\begin{methode}[Réussir une écoute sur un entretien d'embauche : la méthode en 5 étapes]
\begin{enumerate}
\item \textbf{Anticiper (avant l'écoute).} Lisez les questions, soulignez les mots-clés (\eng{post, diploma, how long, languages, when}). Prédisez le vocabulaire attendu.
\item \textbf{Saisir le thème (1\textsuperscript{re} écoute).} Ne notez rien. Identifiez : \eng{Who?} (interviewer / candidate), \eng{Where?} (office), \eng{What for?} (post of secretary).
\item \textbf{Noter les détails (2\textsuperscript{e} écoute).} Écrivez des \cle{mots isolés} : chiffres (\eng{18 months, 2 years}), noms propres (\eng{Miss Ekane, Douala, Hilton Hotel, Monday}), mots de liaison (\eng{but, and, so}).
\item \textbf{Répondre en anglais simple.} Reprenez la structure de la question : \eng{How long...?} $\rightarrow$ \eng{For eighteen months.} / \eng{She worked there for eighteen months.}
\item \textbf{Vérifier.} Accord sujet-verbe (\eng{She has}, pas \eng{She have}), temps cohérent (\eng{past} pour \eng{ago}, \eng{perfect} pour \eng{since/for} durée non terminée), majuscules aux noms propres (\eng{Monday, English, Douala}).
\end{enumerate}
\end{methode}

\begin{exoresolu}[Compréhension détaillée : Entretien chez Douala Trading Company]
\textbf{Questions :}
\begin{enumerate}
\item What post is Miss Ekane applying for?
\item What qualifications does she have?
\item Where did she get her diploma in Computer Studies?
\item How long did she work at the Hilton Hotel?
\item Which languages can she speak?
\item What was her previous position?
\item When can she start work?
\end{enumerate}
\tcblower
\textbf{Raisonnement :} On repère les mots-clés des questions : \eng{post} (1), \eng{qualifications} (2), \eng{where/diploma} (3), \eng{how long} (4 — attend une durée), \eng{languages} (5), \eng{previous position} (6), \eng{when} (7 — attend une date).

\textbf{Réponses rédigées :}
\begin{enumerate}
\item \eng{She is applying for the post of secretary.} (Transformation \eng{you are} $\rightarrow$ \eng{she is}).
\item \eng{She has the GCE Advanced Level and a diploma in Computer Studies.}
\item \eng{She got her diploma from a training centre in Douala.}
\item \eng{She worked there for eighteen months.} (\eng{How long} + \eng{for} + durée).
\item \eng{She can speak French and English fluently.}
\item \eng{She was a receptionist.}
\item \eng{She can start next Monday.} (Pas de préposition avant \eng{next Monday}).
\end{enumerate}
\textbf{Conclusion :} Les réponses reprennent le vocabulaire du script, adaptent les pronoms et les temps, et respectent la syntaxe anglaise (sujet + verbe + complément).
\end{exoresolu}

\courssub{Traiter les questions Vrai/Faux et QCM}

\begin{methode}[Stratégies pour Vrai/Faux et QCM]
\begin{enumerate}
\item \textbf{Vrai/Faux :} Localisez la phrase exacte dans le script. Citez-la pour justifier. \eng{False — the text says “eighteen months”.} Attention aux \cle{distracteurs} : un chiffre mentionné ailleurs dans le texte (\eng{two years} pour le diplôme) ne valide pas une affirmation sur l'expérience.
\item \textbf{QCM :} Lisez les 4 options avant d'écouter. Éliminez les impossibles. Méfiez-vous des mots entendus mais qui ne répondent pas à la question (ex. \eng{secretary} est le poste visé, pas l'ancien emploi).
\item \textbf{Chiffres et noms propres :} Écrivez les nombres en lettres (\eng{eighteen}) dans une phrase complète. Majuscules obligatoires : \eng{Douala, Monday, English, French, GCE, Hilton Hotel}.
\end{enumerate}
\end{methode}

\begin{exoresolu}[Vrai/Faux et QCM sur le script]
\textbf{A. Say True or False. Justify by quoting the text.}
\begin{enumerate}
\item Miss Ekane has a diploma in Computer Studies.
\item She worked at the Hilton Hotel for two years.
\item She can only speak English.
\item She obtained her Advanced Level certificate eighteen months ago.
\item The interview takes place at the Hilton Hotel.
\end{enumerate}

\textbf{B. Choose the correct answer (a, b, c or d).}
\begin{enumerate}
\item Miss Ekane obtained her Advanced Level certificate \dots{}
  \begin{enumerate}
  \item last year \item two years ago \item eighteen months ago \item next Monday
  \end{enumerate}
\item She worked as \dots{}
  \begin{enumerate}
  \item a secretary \item a receptionist \item a manager \item a teacher
  \end{enumerate}
\item She speaks \dots{}
  \begin{enumerate}
  \item only French \item only English \item French and English \item English and German
  \end{enumerate}
\item She can start \dots{}
  \begin{enumerate}
  \item immediately \item next Monday \item in two years \item next month
  \end{enumerate}
\end{enumerate}
\tcblower
\textbf{A. Vrai/Faux :}
\begin{enumerate}
\item \eng{True — the text says “I hold a diploma in Computer Studies”.}
\item \eng{False — the text says “for eighteen months”, not two years.} (Distracteur : \eng{two years} concerne le GCE).
\item \eng{False — she says “I speak both fluently”, that is French and English.}
\item \eng{False — the text says “two years ago”.}
\item \eng{False — the interview takes place at Douala Trading Company (implied by “our company”).}
\end{enumerate}

\textbf{B. QCM :}
\begin{enumerate}
\item (b) \eng{two years ago}. Script : \eng{“I obtained my GCE Advanced Level certificate two years ago.”}
\item (b) \eng{a receptionist}. Script : \eng{“I have worked as a receptionist…”} \eng{Secretary} est le poste visé (piège).
\item (c) \eng{French and English}. Script : \eng{“I speak both fluently.”}
\item (b) \eng{next Monday}. Script : \eng{“I can start next Monday.”}
\end{enumerate}
\end{exoresolu}

\courssub{Lexique thématique : le monde du recrutement}

\begin{definition}[Mots et collocations clés du recrutement]
\begin{center}
\begin{tabularx}{\linewidth}{|X|X|X|}
\hline
\rowcolor{popblueL} \textbf{English} & \textbf{Français} & \textbf{Collocation / Exemple} \\
\hline
\eng{to apply for a post / a job} & postuler à un poste & \eng{apply for} (préposition \textbf{for} obligatoire) \\
\eng{an application form} & un formulaire de candidature & \eng{fill in an application form} \\
\eng{a CV (Curriculum Vitae) / a résumé} & un CV & \eng{send / submit a CV} \\
\eng{a cover letter} & une lettre de motivation & \eng{write a cover letter} \\
\eng{qualifications} & diplômes, qualifications & \eng{academic / professional qualifications} \\
\eng{experience} (U) & expérience & \eng{work experience, previous experience, \textbf{no} experience} \\
\eng{an interview} & un entretien & \eng{attend an interview, pass an interview} \\
\eng{an interviewer / a recruiter} & un recruteur & \eng{the interviewer asked…} \\
\eng{a candidate / an applicant} & un candidat & \eng{shortlisted candidates} \\
\eng{to hire / to recruit / to employ} & embaucher, recruter & \eng{hire someone, be hired} \\
\eng{a vacancy / a post / a position} & un poste vacant & \eng{apply for a vacancy} \\
\eng{full-time / part-time} & temps plein / temps partiel & \eng{a full-time job} \\
\eng{salary / wages} & salaire & \eng{monthly salary, negotiate salary} \\
\eng{strengths / weaknesses} & points forts / points faibles & \eng{What are your strengths?} \\
\eng{availability} & disponibilité & \eng{immediate availability} \\
\hline
\end{tabularx}
\end{center}
\end{definition}

\begin{propriete}[Prépositions pièges à mémoriser]
\begin{center}
\begin{tabular}{|l|l|l|}
\hline
\rowcolor{popblueL} \textbf{Adjectif / Verbe} & \textbf{Préposition} & \textbf{Exemple} \\
\hline
\eng{apply} & \textbf{for} & \eng{apply for a job} \\
\eng{interested} & \textbf{in} & \eng{interested in the post} \\
\eng{good / bad} & \textbf{at} & \eng{good at typing} \\
\eng{qualified} & \textbf{for} & \eng{qualified for the job} \\
\eng{responsible} & \textbf{for} & \eng{responsible for sales} \\
\eng{listen} & \textbf{to} & \eng{listen to the interviewer} \\
\eng{wait} & \textbf{for} & \eng{wait for the result} \\
\hline
\end{tabular}
\end{center}
\end{propriete}

\begin{exoresolu}[Réemploi du lexique : complétion et production]
\textbf{A. Complete with the correct word or phrase:}
\eng{apply for — interview — qualifications — experience — hired — fill in — vacancy — strengths}
\begin{enumerate}
\item I would like to \dots{} the post of driver in your company.
\item There is a \dots{} for a secretary at the Ministry of Finance.
\item The \dots{} will take place on Friday at 10 a.m. in the main office.
\item She has excellent \dots{} : a Master's degree and two professional diplomas.
\item Have you got any \dots{} as a mechanic?
\item Please \dots{} this application form in capital letters.
\item Twenty workers were \dots{} last month after the expansion of the factory.
\item One of my main \dots{} is my ability to work under pressure.
\end{enumerate}

\textbf{B. Write two sentences of your own using \eng{apply for} and \eng{experience}.}
\tcblower
\textbf{A. Solutions :}
\begin{enumerate}
\item \eng{apply for} (base verbale après \eng{would like to}).
\item \eng{vacancy} (nom, synonyme de \eng{post}).
\item \eng{interview} (événement, convient avec \eng{take place}).
\item \eng{qualifications} (pluriel, désigne diplômes et titres).
\item \eng{experience} (indénombrable, pas \eng{an experience} pour l'expérience professionnelle).
\item \eng{fill in} (verbe à particule, \eng{fill in a form}).
\item \eng{hired} (participe passé, passif \eng{were hired}).
\item \eng{strengths} (pluriel, \eng{one of my strengths}).
\end{enumerate}

\textbf{B. Production personnelle (modèles) :}
\eng{I would like to apply for the post of administrative assistant advertised in the newspaper.}
\eng{My sister has three years' experience as a bilingual secretary in a law firm.}
\end{exoresolu}

\courssub{Point de grammaire : Present Perfect vs Past Simple}

\begin{propriete}[Choix du temps : Prétérit (Past Simple) ou Present Perfect ?]
Dans un entretien, l'interviewer alterne les deux temps selon qu'il demande une \textbf{date précise} (passé révolu) ou une \textbf{expérience / durée} (lien avec le présent).

\begin{center}
\begin{tabularx}{\linewidth}{|X|X|X|}
\hline
\rowcolor{popblueL} \textbf{Temps} & \textbf{Marqueurs typiques} & \textbf{Exemple (entretien)} \\
\hline
\textbf{Past Simple} & \eng{yesterday, last week/month/year, \dots{} ago, in 2022, when?} & \eng{I obtained my degree in 2022.} \\
\hline
\textbf{Present Perfect} & \eng{ever, never, already, just, yet, for, since, how long?} & \eng{I have worked here for two years.} \\
\hline
\end{tabularx}
\end{center}

\textbf{Règles clés :}
\begin{itemize}
\item \eng{For} + durée (\eng{for three years}) $\rightarrow$ Present Perfect.
\item \eng{Since} + point de départ (\eng{since 2021}) $\rightarrow$ Present Perfect.
\item \eng{Ago} / date précise $\rightarrow$ Past Simple.
\item Question \eng{Have you ever…?} $\rightarrow$ Réponse courte \eng{Yes, I have / No, I haven't.}
\item \eng{Just} (venir de) se place entre auxiliaire et participe passé : \eng{He has just arrived.}
\end{itemize}
\end{propriete}

\begin{exoresolu}[Prétérit ou Present Perfect ? Mettez le verbe au temps correct.]
\begin{enumerate}
\item I (obtain) my GCE Advanced Level two years ago.
\item She (work) as a receptionist since 2022.
\item (you / ever / attend) a job interview?
\item We (not / see) the manager yesterday.
\item He (just / finish) his training.
\item How long (you / know) Mr Ndi?
\item The company (hire) ten new employees last month.
\item I (never / write) a cover letter before.
\end{enumerate}
\tcblower
\textbf{Solutions et justifications :}
\begin{enumerate}
\item \eng{obtained} — \eng{two years ago} impose le Past Simple (verbe régulier : \eng{-ed}).
\item \eng{has worked} — \eng{since 2022} (point de départ) impose Present Perfect (\eng{she} $\rightarrow$ \eng{has}).
\item \eng{Have you ever attended} — \eng{ever} (expérience) appelle Present Perfect ; inversion sujet/auxiliaire.
\item \eng{did not see} (ou \eng{didn't see}) — \eng{yesterday} impose Past Simple ; négation avec auxiliaire \eng{did} + base verbale.
\item \eng{has just finished} — \eng{just} + Present Perfect ; \eng{just} entre \eng{has} et \eng{finished}.
\item \eng{have you known} — \eng{How long} + Present Perfect (durée non terminée). Verbe \eng{know} = état, pas d'action.
\item \eng{hired} — \eng{last month} (date passée) impose Past Simple.
\item \eng{have never written} — \eng{never} + Present Perfect (expérience de vie). Participe passé irrégulier : \eng{write - wrote - written}.
\end{enumerate}
\end{exoresolu}

\courssub{Production orale guidée : Simuler un mini-entretien}

\begin{methode}[Conduire un dialogue d'entretien en 4 temps]
\begin{enumerate}
\item \textbf{Salutations et installation} : \eng{Good morning, Sir/Madam. Please have a seat. / Thank you.}
\item \textbf{Questions types de l'interviewer} (à mémoriser) :
  \begin{itemize}
  \item \eng{Can you tell me about yourself?}
  \item \eng{What are your qualifications?}
  \item \eng{Have you got any experience? / What experience do you have?}
  \item \eng{What are your strengths and weaknesses?}
  \item \eng{Why do you want this job?}
  \item \eng{When can you start?}
  \item \eng{Do you have any questions?}
  \end{itemize}
\item \textbf{Réponses du candidat} : Phrase complète + détail concret. Utilisez \eng{I am good at…}, \eng{I have experience in…}, \eng{I can…}, \eng{I am available…}.
\item \textbf{Conclusion polie} : \eng{Thank you for your time. We will contact you soon. / Thank you, Sir. Goodbye.}
\end{enumerate}
\end{methode}

\begin{experience}[Jeu de rôle : Entretien simulé (binômes)]
\textbf{Consigne :} L'élève A est le recruteur (Directeur d'une entreprise de commerce à Douala), l'élève B est le candidat (titulaire d'un GCE A Level, diplôme en comptabilité, 1 an comme vendeur au Marché Mokolo).
\textbf{Durée :} 3 minutes par binôme. Changement de rôles.
\textbf{Grille d'évaluation (professeur) :} Salutations correctes, utilisation de \eng{apply for}, \eng{experience}, \eng{qualifications}, \eng{strengths}, temps verbaux justes (Past Simple / Present Perfect), politesse (\eng{Sir, Madam, please, thank you}), fluidité.
\end{experience}

\begin{exoresolu}[Complétez le dialogue d'entretien]
Complétez les répliques du candidat par des phrases complètes et correctes.

\eng{Interviewer: Good morning. Please sit down. Can you tell me about yourself?}
\eng{Candidate:} \dots{}

\eng{Interviewer: What are your qualifications?}
\eng{Candidate:} \dots{}

\eng{Interviewer: Have you got any experience?}
\eng{Candidate:} \dots{}

\eng{Interviewer: What are your strengths?}
\eng{Candidate:} \dots{}

\eng{Interviewer: When can you start?}
\eng{Candidate:} \dots{}
\tcblower
\textbf{Solution rédigée (modèle) :}
\begin{enumerate}
\item \eng{Good morning, Sir. My name is Alain Fouda. I am twenty years old and I live in Yaoundé.} (Présentation standard : nom, âge, lieu).
\item \eng{I have my GCE Advanced Level and a diploma in Accounting.} (\eng{diploma in} + matière).
\item \eng{Yes, I have. I worked as a shop assistant at Mokolo Market for one year.} (Réponse courte Present Perfect + détail Past Simple car durée close \eng{for one year}).
\item \eng{I am good at communication and I can speak English and French fluently. I am also punctual and hard-working.} (Vocabulaire qualités : \eng{hard-working, punctual, fluent}).
\item \eng{I can start immediately, Sir. / I can start next Monday.} (\eng{immediately} / \eng{next Monday} sans préposition).
\end{enumerate}
\end{exoresolu}

\courssub{Production écrite : Rédiger un CV et une lettre de motivation simples}

\begin{methode}[Écrire un CV (Résumé) et une Cover Letter (Lettre de motivation)]
\begin{enumerate}
\item \textbf{CV :} Structure standard : \eng{Personal Details} (Name, Address, Phone, Email), \eng{Education} (du plus récent au plus ancien), \eng{Work Experience} (Poste, Entreprise, Dates, Tâches), \eng{Skills} (Languages, IT, Driving licence), \eng{References}.
\item \textbf{Cover Letter :} 3 paragraphes.
  \begin{enumerate}
  \item \textbf{Objet / Accroche :} \eng{I am writing to apply for the post of… advertised in…}
  \item \textbf{Argumentaire :} \eng{I have… qualifications. I have… experience. I am… (qualités).}
  \item \textbf{Clôture :} \eng{I am available for an interview at your convenience. I look forward to hearing from you. Yours faithfully, [Name].}
  \end{enumerate}
\item \textbf{Conventions :} \eng{Yours faithfully} (si destinataire inconnu / \eng{Dear Sir/Madam}), \eng{Yours sincerely} (si nom connu / \eng{Dear Mr Ndi}). Date format : \eng{15th January 2025} ou \eng{January 15, 2025}.
\end{enumerate}
\end{methode}

\begin{exemplebox}[Modèle de Cover Letter]
\eng{Alain Fouda} \\
\eng{123 Rue de la Paix, Yaoundé} \\
\eng{Tel: +237 6XX XX XX XX} \\
\eng{Email: alain.fouda@email.com} \\
\eng{} \\
\eng{15th January 2025} \\
\eng{} \\
\eng{The Human Resources Manager} \\
\eng{Douala Trading Company} \\
\eng{P.O. Box 456, Douala} \\
\eng{} \\
\eng{Dear Sir/Madam,} \\
\eng{} \\
\eng{I am writing to apply for the post of Secretary advertised in the Cameroon Tribune of January 10th.} \\
\eng{I hold a GCE Advanced Level certificate and a Diploma in Computer Studies. I have worked as a receptionist at the Hilton Hotel for eighteen months, where I developed strong communication and organisational skills. I speak English and French fluently.} \\
\eng{I am hard-working, punctual and eager to learn. I am available for an interview at your convenience.} \\
\eng{} \\
\eng{Yours faithfully,} \\
\eng{} \\
\eng{Alain Fouda}
\end{exemplebox}

\begin{attention}[Les 5 erreurs qui coûtent des points au Bac]
\begin{enumerate}
\item \textbf{Faux ami :} \eng{To assist} $\neq$ « assister à ». \eng{To assist} = aider. « Assister à un entretien » = \eng{to attend an interview}. « Assister à une réunion » = \eng{to attend a meeting}.
\item \textbf{Prépositions obligatoires :} \eng{apply \textbf{for} a job} (jamais \eng{apply a job}), \eng{interested \textbf{in}} (jamais \eng{interested for}), \eng{good \textbf{at}} (jamais \eng{good in} pour une compétence), \eng{listen \textbf{to}} (jamais \eng{listen} seul + COD).
\item \textbf{Temps verbaux :} \eng{Ago, yesterday, last year, in 2022} $\rightarrow$ \textbf{Past Simple} (jamais Present Perfect). \eng{Wrong: I have obtained my degree two years ago.} $\rightarrow$ \eng{Right: I obtained my degree two years ago.}
\item \textbf{Ordre des mots (Question indirecte) :} Pas d'inversion après \eng{Can you tell me…} / \eng{I'd like to know…}. \eng{Wrong: Can you tell me what are your strengths?} $\rightarrow$ \eng{Right: Can you tell me what your strengths are?}
\item \textbf{Orthographe et Majuscules :} \eng{Experience} (1 \eng{e}, pas \eng{expérience}), \eng{Interview} (2 \eng{i}), \eng{Curriculum Vitae} (majuscules), \eng{English, French, Monday, January, Douala, Cameroon} (majuscules systématiques aux langues, jours, mois, noms propres).
\end{enumerate}
\end{attention}

\courssub{Exercices d'entraînement (Devoirs / Travail en classe)}

\begin{exercice}[Compréhension orale / écrite : Nouvel entretien]
\textbf{Écoutez (ou lisez) le script suivant et répondez aux questions.}

\eng{Interviewer: Good afternoon, Mr Atangana. Thanks for coming.}
\eng{Candidate: Good afternoon.}
\eng{Interviewer: You are applying for the position of sales representative. Tell me, what is your highest qualification?}
\eng{Candidate: I have a Bachelor's degree in Marketing from the University of Buea, obtained in 2021.}
\eng{Interviewer: Have you worked in sales before?}
\eng{Candidate: Yes, I have. I worked for a cocoa export company in Douala for two years, from 2021 to 2023.}
\eng{Interviewer: Why did you leave that job?}
\eng{Candidate: The contract ended. It was a fixed-term contract.}
\eng{Interviewer: What languages do you speak?}
\eng{Candidate: I speak English, French and Pidgin fluently.}
\eng{Interviewer: Do you have a driving licence?}
\eng{Candidate: Yes, I do. I have had it for three years.}
\eng{Interviewer: Good. When are you available?}
\eng{Candidate: I can start at the beginning of next month.}

\textbf{Questions :}
\begin{enumerate}
\item \textbf{True or False? Justify.}
  \begin{enumerate}
  \item Mr Atangana has a Master's degree.
  \item He worked in a cocoa export company for three years.
  \item He left his previous job because he was fired.
  \item He speaks three languages.
  \item He has had a driving licence since 2021.
  \end{enumerate}
\item \textbf{Answer in full sentences.}
  \begin{enumerate}
  \item What position is he applying for?
  \item Where did he study Marketing?
  \item How long did he work in Douala?
  \item Why is he available now?
  \item When can he start?
  \end{enumerate}
\end{enumerate}
\end{exercice}

\begin{exercice}[Grammaire : Temps verbaux et prépositions]
\begin{enumerate}
\item Put the verbs in brackets into the \textbf{Past Simple} or the \textbf{Present Perfect}.
  \begin{enumerate}
  \item I (send) my CV last week.
  \item She (work) in this company since 2020.
  \item (you / ever / write) a cover letter?
  \item They (not / hire) anyone last month.
  \item He (just / receive) a job offer.
  \end{enumerate}
\item Complete with the correct preposition: \eng{for, in, at, on, to, of}.
  \begin{enumerate}
  \item She is interested \dots{} the post of secretary.
  \item He applied \dots{} the job advertised in the paper.
  \item They are good \dots{} languages.
  \item The interview is \dots{} Monday \dots{} 10 a.m.
  \item I have worked here \dots{} five years.
  \end{enumerate}
\end{enumerate}
\end{exercice}

\begin{exercice}[Vocabulaire : Le bon mot]
Complétez les phrases avec un mot de la liste : \eng{qualifications, experience, vacancy, interviewer, hired, application form, strengths, salary}.
\begin{enumerate}
\item There is a \dots{} for a driver at the town hall.
\item Please fill in the \dots{} carefully.
\item What are your \dots{}? — I am punctual and honest.
\item The \dots{} asked difficult questions.
\item Ten people were \dots{} yesterday.
\item The \dots{} offered is 150,000 FCFA per month.
\item She has good \dots{} : a BTS and a Licence.
\item Do you have any \dots{} in teaching?
\end{enumerate}
\end{exercice}

\begin{exercice}[Expression écrite : Lettre de motivation]
Vous êtes \eng{Marie Njoya}, titulaire d'un \eng{GCE Advanced Level} et d'un \eng{Diploma in Secretarial Studies}. Vous avez travaillé 6 mois comme stagiaire (\eng{intern}) dans un cabinet d'avocats à Yaoundé. Vous parlez anglais et français. Vous répondez à une annonce pour un poste de \eng{Junior Secretary} à la \eng{Bamenda City Council}. Rédigez la \eng{cover letter} (120--150 mots). Respectez la structure en 3 paragraphes et les formules de politesse.
\end{exercice}

\begin{corrige}[Exercices d'entraînement]
\textbf{Exercice 1 (Compréhension) :}
\textbf{A. True/False :}
\begin{enumerate}
\item \eng{False — the text says “Bachelor's degree in Marketing”.}
\item \eng{False — the text says “for two years”.}
\item \eng{False — the text says “The contract ended. It was a fixed-term contract.”}
\item \eng{True — the text says “I speak English, French and Pidgin fluently”.}
\item \eng{False — the text says “I have had it for three years” (donc depuis 2022 si on est en 2025), pas depuis 2021.}
\end{enumerate}
\textbf{B. Réponses courtes :}
\begin{enumerate}
\item \eng{He is applying for the position of sales representative.}
\item \eng{He studied Marketing at the University of Buea.}
\item \eng{He worked there for two years.}
\item \eng{His fixed-term contract ended.}
\item \eng{He can start at the beginning of next month.}
\end{enumerate}

\textbf{Exercice 2 (Grammaire) :}
\textbf{1. Temps verbaux :}
\begin{enumerate}
\item \eng{sent} (Past Simple : \eng{last week}).
\item \eng{has worked} (Present Perfect : \eng{since 2020}).
\item \eng{Have you ever written} (Present Perfect : \eng{ever}).
\item \eng{did not hire} / \eng{didn't hire} (Past Simple : \eng{last month}).
\item \eng{has just received} (Present Perfect : \eng{just}).
\end{enumerate}
\textbf{2. Prépositions :}
\begin{enumerate}
\item \eng{in} (\eng{interested in}).
\item \eng{for} (\eng{apply for}).
\item \eng{at} (\eng{good at}).
\item \eng{on} (jour) \dots{} \eng{at} (heure précise).
\item \eng{for} (\eng{for five years} = durée).
\end{enumerate}

\textbf{Exercice 3 (Vocabulaire) :}
\begin{enumerate}
\item \eng{vacancy}
\item \eng{application form}
\item \eng{strengths}
\item \eng{interviewer}
\item \eng{hired}
\item \eng{salary}
\item \eng{qualifications}
\item \eng{experience}
\end{enumerate}

\textbf{Exercice 4 (Expression écrite — Modèle) :}
\eng{Marie Njoya} \\
\eng{Quarter Mvog-Ada, Yaoundé} \\
\eng{Tel: +237 6XX XX XX XX} \\
\eng{Email: marie.njoya@email.com} \\
\eng{} \\
\eng{20th January 2025} \\
\eng{} \\
\eng{The Secretary General} \\
\eng{Bamenda City Council} \\
\eng{P.O. Box 123, Bamenda} \\
\eng{} \\
\eng{Dear Sir/Madam,} \\
\eng{} \\
\eng{I am writing to apply for the post of Junior Secretary advertised in the Cameroon Tribune of January 18th.} \\
\eng{I hold a GCE Advanced Level certificate and a Diploma in Secretarial Studies. I have worked as an intern in a law firm in Yaoundé for six months, where I gained experience in filing, typing minutes and welcoming clients. I speak English and French fluently and I have good computer skills.} \\
\eng{I am dynamic, organised and ready to serve the Council. I am available for an interview at your earliest convenience.} \\
\eng{} \\
\eng{Yours faithfully,} \\
\eng{} \\
\eng{Marie Njoya}
\end{corrige}

\newpage

\coursec{Exercices}

\courssub{Je vérifie mes connaissances}

\begin{exercice}[QCM : Compréhension globale]
\textbf{Écoutez ou lisez le texte suivant, puis choisissez la bonne réponse (a, b ou c).}
\begin{textebox}[Script : A First Interview]
Good morning, Mr. Atangana. Thank you for coming today. We received your application for the position of \cle{junior accountant} last week. Your \cle{CV} looks impressive. You graduated from the University of Yaoundé II with a Master’s in Finance, didn't you?
\end{textebox}
\begin{enumerate}
\item What position is Mr. Atangana applying for?
      \begin{enumerate}
      \item Senior Manager
      \item Junior Accountant
      \item Human Resources Assistant
      \end{enumerate}
\item Where did he study?
      \begin{enumerate}
      \item University of Buea
      \item University of Yaoundé I
      \item University of Yaoundé II
      \end{enumerate}
\item What qualification does he hold?
      \begin{enumerate}
      \item Bachelor in Marketing
      \item Master’s in Finance
      \item PhD in Economics
      \end{enumerate}
\end{enumerate}
\end{exercice}

\begin{exercice}[Vrai ou Faux : Justifiez par une phrase du texte]
\textbf{Lisez le texte ci-dessous. Pour chaque affirmation, cochez V (Vrai) ou F (Faux) et recopiez la phrase du texte qui justifie votre réponse.}
\begin{textebox}[Interview Excerpt]
-- \textbf{Interviewer:} Do you have any previous work experience?
-- \textbf{Candidate:} Yes, I worked as an intern at \cle{SGBC} Bank in Douala for six months last year. I handled customer deposits and learned to use their accounting software.
-- \textbf{Interviewer:} That’s good. Are you available to start immediately?
-- \textbf{Candidate:} Well, I have to give one month's notice to my current employer, a small audit firm in Bastos.
\end{textebox}
\begin{enumerate}
\item The candidate has never worked in a bank.
\item He did an internship lasting six months.
\item He can start working tomorrow.
\item He currently works in an audit firm located in Bastos.
\end{enumerate}
\end{exercice}

\begin{exercice}[Vocabulaire : Le monde du recrutement]
\textbf{Reliez chaque mot ou expression de la colonne A à sa définition dans la colonne B.}
\begin{center}
\begin{tabularx}{\linewidth}{|X|X|}
\hline
\rowcolor{popblueL} \textbf{A : Vocabulaire} & \textbf{B : Définition (en anglais)} \\
\hline
1. \cle{Vacancy} & A. A document written by a previous employer describing your qualities. \\
\hline
2. \cle{Applicant} & B. The money you earn every month for doing a job. \\
\hline
3. \cle{Salary} & C. A person who applies for a job. \\
\hline
4. \cle{To hire} & D. A job that is available for someone to do. \\
\hline
5. \cle{Reference letter} & E. To give someone a job. \\
\hline
\end{tabularx}
\end{center}
\textbf{Complétez ensuite les phrases suivantes avec les mots du tableau (forme correcte).}
\begin{enumerate}
\item There is a \_\_\_\_\_\_ for a receptionist at the hotel near the \cle{rond-point Deïdo}.
\item Over fifty \_\_\_\_\_\_ sent their CVs for that single position.
\item The \_\_\_\_\_\_
\end{enumerate}
\end{exercice}


\newpage

\coursec{Corrigés des exercices}

\begin{corrige}[Exercice 1 — QCM : Compréhension globale]

\textbf{Réponses correctes :}

\begin{enumerate}
\item \textbf{b — Junior Accountant.} Justification (extrait du script) : \eng{“We received your application for the position of junior accountant last week.”} (« Nous avons reçu votre candidature pour le poste de comptable junior la semaine dernière. »)
\item \textbf{c — University of Yaoundé II.} Justification : \eng{“You graduated from the University of Yaoundé II with a Master's in Finance, didn't you?”} (« Vous avez obtenu votre diplôme à l'Université de Yaoundé II avec un master en finance, n'est-ce pas ? »)
\item \textbf{b — Master's in Finance.} Justification : même phrase que ci-dessus ; le diplôme mentionné est un \eng{Master's in Finance} (« un master en finance »).
\end{enumerate}

\textbf{Points de vigilance :}
\begin{itemize}
\item Ne pas confondre \eng{junior accountant} (« comptable débutant ») et \eng{senior manager} (« cadre supérieur ») : les deux mots \eng{junior} et \eng{senior} sont des mots-clés de l'écoute ; repérez l'adjectif qui précède le nom du poste.
\item Attention aux noms propres : \eng{Yaoundé I} et \eng{Yaoundé II} se ressemblent à l'oreille ; à l'écoute, concentrez-vous sur le chiffre prononcé (« one » / « two »).
\item \eng{Master's} (avec apostrophe) est l'abréviation de \eng{Master's degree} ; ne pas écrire \emph{Masters}.
\end{itemize}

\textbf{Barème indicatif (type évaluation) :} 1 point par bonne réponse, soit 3 points au total.

\end{corrige}

\begin{corrige}[Exercice 2 — Vrai ou Faux : Justifiez par une phrase du texte]

\begin{enumerate}
\item \textbf{F (Faux).} Justification (extrait du script) : \eng{“Yes, I worked as an intern at SGBC Bank in Douala for six months last year.”} (« Oui, j'ai travaillé comme stagiaire à la banque SGBC à Douala pendant six mois l'année dernière. ») Le candidat \emph{a bien} travaillé dans une banque.
\item \textbf{V (Vrai).} Justification : \eng{“I worked as an intern at SGBC Bank in Douala for six months last year.”} Un \eng{intern} est un stagiaire ; \eng{to work as an intern} = faire un stage. La durée est bien de six mois (\eng{for six months}).
\item \textbf{F (Faux).} Justification : \eng{“Well, I have to give one month's notice to my current employer.”} (« Eh bien, je dois donner un préavis d'un mois à mon employeur actuel. ») Il ne peut donc pas commencer demain : il doit respecter un préavis d'un mois.
\item \textbf{V (Vrai).} Justification : \eng{“...my current employer, a small audit firm in Bastos.”} (« ...mon employeur actuel, un petit cabinet d'audit à Bastos. ») \eng{Current} signifie « actuel » ; le cabinet est bien situé à Bastos.
\end{enumerate}

\textbf{Points de vigilance :}
\begin{itemize}
\item \eng{Notice} est ici un nom signifiant « préavis » ; \eng{to give notice} = donner sa démission avec préavis. Ne pas traduire par « remarquer ».
\item Faux ami : \eng{actually / current} signifient « en réalité / actuel » ; « actuellement » se dit \eng{currently}, pas \emph{actually}.
\item \eng{Firm} = « cabinet, entreprise » ; ne pas confondre avec l'adjectif \eng{firm} (« ferme, solide »).
\item Dans la justification, recopiez la phrase \emph{exactement} comme dans le texte, sans la modifier.
\end{itemize}

\textbf{Barème indicatif :} 0,5 point pour la réponse V/F correcte + 0,5 point pour la justification exacte, par item ; soit 4 points au total.

\end{corrige}

\begin{corrige}[Exercice 3 — Vocabulaire : Le monde du recrutement]

\textbf{Correspondances (reliement) :}

\begin{center}
\begin{tabularx}{\linewidth}{|X|X|X|}
\hline
\rowcolor{popblueL} \textbf{Vocabulaire} & \textbf{Définition} & \textbf{Traduction} \\
\hline
1. \cle{Vacancy} & D. A job that is available for someone to do. & un poste vacant \\
\hline
2. \cle{Applicant} & C. A person who applies for a job. & un candidat \\
\hline
3. \cle{Salary} & B. The money you earn every month for doing a job. & un salaire \\
\hline
4. \cle{To hire} & E. To give someone a job. & embaucher \\
\hline
5. \cle{Reference letter} & A. A document written by a previous employer describing your qualities. & une lettre de recommandation \\
\hline
\end{tabularx}
\end{center}

\textbf{Phrases complétées :}

\begin{enumerate}
\item \eng{There is a \textbf{vacancy} for a receptionist at the hotel near the Deïdo roundabout.} (« Il y a un poste vacant de réceptionniste à l'hôtel près du rond-point Deïdo. ») — Nom singulier après l'article indéfini \eng{a}.
\item \eng{Over fifty \textbf{applicants} sent their CVs for that single position.} (« Plus de cinquante candidats ont envoyé leur CV pour ce seul poste. ») — Pluriel obligatoire après \eng{fifty} ; le possessif \eng{their} confirme le pluriel.
\item Phrase proposée pour compléter l'énoncé : \eng{The \textbf{salary} offered for this job is very attractive.} (« Le salaire proposé pour ce poste est très intéressant. ») — Nom singulier avec l'article défini \eng{the}. Autres réponses possibles : \eng{The company decided to \textbf{hire} three new workers.} (verbe après \eng{to}) ; \eng{My former boss wrote me a \textbf{reference letter}.} (nom singulier après \eng{a}).
\end{enumerate}

\textbf{Points de vigilance :}
\begin{itemize}
\item Faux ami : \eng{applicant} = « candidat (à un poste) » ; « un candidat » à un examen se dit \eng{candidate}. Les deux sont possibles pour un emploi, mais \eng{applicant} vient de \eng{to apply for a job} (« postuler »).
\item \eng{Salary} est un nom dénombrable singulier ; on dit \eng{a good salary}, jamais \emph{a good money} (\eng{money} est indénombrable).
\item \eng{To hire} = « embaucher » ; attention, en Grande-Bretagne \eng{to hire} signifie aussi « louer » (une voiture), alors qu'aux États-Unis on préfère \eng{to rent}.
\item Orthographe : \eng{vacancy} (un seul « c », terminaison \emph{-ancy}), prononcé « vé-kan-si ».
\end{itemize}

\textbf{Barème indicatif :} 0,5 point par correspondance correcte (2,5 points) + 0,5 point par phrase correctement complétée (1,5 point) ; soit 4 points au total.

\end{corrige}

\newpage

\coursec{Fiche bilan}

\begin{popfigure}
\begin{tikzpicture}[mindmap, grow cyclic, every node/.style={concept, align=center, font=\small}, root concept/.append style={concept color=popblue, font=\bfseries\small}, level 1/.append style={level distance=3.2cm, sibling angle=51.4}, text=white]
\node[root concept, align=center]{Job Interview\\Listening}
  child[concept color=poporange]{ node[align=center]{Key\\vocabulary} }
  child[concept color=popgreen]{ node[align=center]{Listening\\strategies} }
  child[concept color=poppurple]{ node[align=center]{Questions\\types} }
  child[concept color=poppink]{ node[align=center]{Grammar\\focus} }
  child[concept color=popgold]{ node[align=center]{Interview\\stages} }
  child[concept color=popblue!70]{ node[align=center]{Common\\mistakes} }
  child[concept color=popgreen!70]{ node[align=center]{Reuse \&\\speaking} };
\end{tikzpicture}
\legende{Carte mentale de la leçon : comprendre un entretien d'embauche à l'écoute.}
\end{popfigure}

\begin{bilan}
\textbf{1. Lexique clé du recrutement (English / Français)}
\begin{center}
\begin{tabularx}{\linewidth}{|X|X|}\hline
\rowcolor{popblueL} English & Français \\ \hline
a job interview & un entretien d'embauche \\ \hline
an applicant / a candidate & un candidat \\ \hline
the interviewer / the panel & le recruteur / le jury \\ \hline
to apply for a job & postuler à un emploi \\ \hline
an application letter & une lettre de motivation \\ \hline
a CV (curriculum vitae) / a résumé & un CV \\ \hline
qualifications & les diplômes, les qualifications \\ \hline
work experience & l'expérience professionnelle \\ \hline
strengths and weaknesses & points forts et points faibles \\ \hline
to hire / to recruit & embaucher / recruter \\ \hline
a salary / wages & un salaire \\ \hline
a vacancy / an opening & un poste vacant \\ \hline
to be shortlisted & être présélectionné \\ \hline
punctual, confident, polite & ponctuel, confiant, poli \\ \hline
\end{tabularx}
\end{center}

\textbf{2. Grammaire à connaître par cœur}
\begin{center}
\begin{tabularx}{\linewidth}{|X|X|}\hline
\rowcolor{popblueL} Règle & Exemple \\ \hline
Prétérit simple pour raconter l'entretien passé & \eng{The interview started at 9 a.m.} \\ \hline
Present perfect pour l'expérience (durée non terminée) & \eng{I have worked here for two years.} \\ \hline
Questions indirectes : pas d'inversion & \eng{Could you tell me why you applied?} \\ \hline
Présent simple pour les habitudes et qualités & \eng{I work well under pressure.} \\ \hline
Futur avec \eng{will} / \eng{be going to} pour les projets & \eng{I am going to improve my English.} \\ \hline
\end{tabularx}
\end{center}

\textbf{3. Méthode express d'écoute}
\begin{itemize}
\item \cle{Avant} : lire les questions, souligner les mots-clés, anticiper le vocabulaire.
\item \cle{Pendant} : repérer le thème, noter les chiffres (heures, salaires, dates), les noms propres et les verbes d'action.
\item \cle{Après} : vérifier la cohérence des réponses, réemployer le lexique dans des phrases complètes.
\end{itemize}

\textbf{4. Réponses types aux questions d'entretien}
\begin{itemize}
\item \eng{Tell me about yourself.} $\rightarrow$ \eng{I am a final-year student with a passion for...}
\item \eng{Why do you want this job?} $\rightarrow$ \eng{Because I want to gain experience and serve my community.}
\item \eng{What are your strengths?} $\rightarrow$ \eng{I am punctual, hardworking and I learn fast.}
\end{itemize}
\end{bilan}

\begin{aretenir}[Checklist avant le Bac]
\begin{itemize}
\item[$\square$] Je connais le vocabulaire essentiel de l'entretien d'embauche (applicant, to apply, CV, salary, vacancy...).
\item[$\square$] Je sais lire les questions \emph{avant} l'écoute et souligner les mots-clés.
\item[$\square$] Je repère rapidement les chiffres, les heures, les dates et les noms propres dans un enregistrement.
\item[$\square$] Je distingue le thème général (gist) des détails précis.
\item[$\square$] Je réponds aux questions vrai/faux en justifiant par une phrase du script.
\item[$\square$] J'utilise correctement le prétérit pour raconter un entretien passé.
\item[$\square$] J'emploie le present perfect avec \eng{for} et \eng{since} pour parler d'expérience.
\item[$\square$] Je formule une question indirecte sans inversion : \eng{Could you tell me where you studied?}
\item[$\square$] Je fais attention aux faux amis : \eng{to apply} (postuler), \eng{sensible} (raisonnable), \eng{a salary} (un salaire).
\item[$\square$] Je sais rédiger trois phrases complètes réemployant le lexique du recrutement.
\item[$\square$] Je peux jouer un mini-entretien : salutations, présentation, qualités, remerciements.
\end{itemize}
\end{aretenir}

\findecours

\end{document}
